% Espaces vectoriels sur ℝ — Mathématiques 1ère E — Cours signé OnBuch+
% Programme officiel MINESEC. Compiler avec : tectonic N16-espaces-vectoriels.tex
\newcommand{\DOCMATIERE}{Mathématiques}
\newcommand{\DOCNIVEAU}{1ère E}
\newcommand{\DOCMODULE}{Configurations et transformations élémentaires du plan}
\newcommand{\DOCLECON}{N16}
\newcommand{\DOCTITRE}{Espaces vectoriels sur ℝ}
% OnBuch+ — préambule « cours signés » ENRICHI (compilé avec tectonic / XeLaTeX)
% Système visuel « Pop » d'OnBuch+ : fond crème (jamais blanc), bordures encre
% épaisses, ombres « dures » décalées, titres Archivo Black en capitales.
% Version MATHÉMATIQUES : graphes (pgfplots/tikz), boîtes pédagogiques
% (définition, propriété, méthode, attention, le savais-tu, exercice résolu,
% exercice, TP, bilan), page de garde et signature OnBuch+ ancrée en bas.
%
% Macros attendues dans main.tex AVANT \input{preamble} :
%   \DOCMATIERE  \DOCNIVEAU  \DOCMODULE  \DOCLECON (numéro)  \DOCTITRE
\documentclass[11pt]{article}

\usepackage{fontspec}
\usepackage[french]{babel}
\usepackage[a4paper,margin=2cm,top=3.4cm,bottom=2.8cm,headheight=1.4cm,footskip=1.3cm]{geometry}
\usepackage{amsmath,amssymb,amsfonts}
\usepackage{mathtools} % \xmapsto, \coloneqq... souvent utilisés par les modèles
\usepackage{cancel} % simplifications barrées
% \abs{z} (module, valeur absolue) et \norm{u} (norme), utilisés par les modèles
\providecommand{\abs}{}\let\abs\relax\DeclarePairedDelimiter{\abs}{\lvert}{\rvert}
\providecommand{\norm}{}\let\norm\relax\DeclarePairedDelimiter{\norm}{\lVert}{\rVert}
\usepackage[table]{xcolor} % [table] : fournit aussi \rowcolors (alternance de lignes)
\usepackage{pagecolor}
\usepackage{tikz}
\usetikzlibrary{arrows.meta,positioning,calc,shapes.geometric,shapes.misc,shapes.arrows,decorations.pathmorphing,decorations.pathreplacing,decorations.markings,patterns,shadows,mindmap,trees,fit,backgrounds,matrix,intersections,angles,quotes}
\usepackage{pgfplots}
\usepackage{pgf-pie} % diagrammes circulaires \pie (statistiques)
\pgfplotsset{compat=1.18}
\usepgfplotslibrary{fillbetween,groupplots} % aires entre courbes (calcul intégral)
\usepackage{enumitem}
\usepackage{fancyhdr}
% [most] charge toutes les bibliothèques tcolorbox (listings, minted, raster,
% documentation...) et gonfle inutilement la mémoire TeX sur un document long
% (~300 boîtes) : on ne charge que ce qui est réellement utilisé.
\usepackage[skins,breakable]{tcolorbox}
\usepackage{graphicx}
\usepackage{colortbl}
\usepackage{booktabs}
\usepackage{tabularx}
\usepackage{diagbox} % case d'en-tête diagonale des tableaux à double entrée
% Colonnes extensibles centrée / gauche / droite (C, L, R), souvent utilisées par les modèles
\newcolumntype{C}{>{\centering\arraybackslash}X}
\newcolumntype{L}{>{\raggedright\arraybackslash}X}
\newcolumntype{R}{>{\raggedleft\arraybackslash}X}
\usepackage{array}
\usepackage{multicol}
\usepackage{siunitx}
% parse-numbers=false : accepte aussi \SI{2,5 \times 10^{-3}}{...}
\sisetup{locale=FR,output-decimal-marker={,},group-minimum-digits=4,parse-numbers=false}
\DeclareSIUnit\ohm{\ensuremath{\symup{\Omega}}} % U+2126 absent de la police
\usepackage{qrcode}
% \degree : commande fréquemment utilisée par les modèles pour un angle ou une
% température en dehors de \SI{}{} (non fournie par les paquets chargés).
\providecommand{\degree}{^{\circ}}
% \qed (fin de démonstration), utilisé par les modèles sans amsthm
\providecommand{\qed}{\hfill$\square$}
% \barre : double barre des valeurs interdites dans un tableau de variations (tkz-tab)
\providecommand{\barre}{\ensuremath{\|}}
% environnement proof (amsthm non chargé) utilisé par les modèles
\ifdefined\proof\else\newenvironment{proof}[1][Démonstration]{\par\noindent\textit{#1.}\ }{\qed\par}\fi
\usepackage{lastpage}
\usepackage{hyperref}
\hypersetup{hidelinks}

% ============================================================
% Palette « Pop » OnBuch+ — mêmes hex que la classe P (plus_ui.dart)
% ============================================================
\definecolor{poporange}{HTML}{F59321}
\definecolor{popdark}{HTML}{E07A0C}
\definecolor{popcream}{HTML}{FFF6E8}
\definecolor{popink}{HTML}{1C1714}
\definecolor{poppurple}{HTML}{7A5AE0}
\definecolor{popgreen}{HTML}{2E9E6A}
\definecolor{poppink}{HTML}{E0457B}
\definecolor{popblue}{HTML}{2F7DE1}
\definecolor{popgold}{HTML}{C98A12}
\definecolor{popmuted}{HTML}{8A7F76}
\definecolor{popline}{HTML}{EADFD2}
\definecolor{popgray}{HTML}{8A7F76}
\colorlet{popgrey}{popgray}
% Alias tolérants (le modèle nomme parfois les couleurs différemment)
\colorlet{popwhite}{white}
\colorlet{popblack}{popink}
\colorlet{popred}{poppink}
\colorlet{popyellow}{popgold}
% Teintes claires pour fonds de tableaux / zones
\colorlet{popblueL}{popblue!10!white}
\colorlet{poporangeL}{poporange!14!white}
\colorlet{popgreenL}{popgreen!12!white}
\colorlet{poppurpleL}{poppurple!12!white}
\colorlet{poppinkL}{poppink!12!white}
\colorlet{popgoldL}{popgold!14!white}

\pagecolor{popcream}

% ============================================================
% Typographie
% ============================================================
\defaultfontfeatures{Ligatures=TeX}
\setmainfont{PlusJakartaSans-Medium.ttf}[Path=../../../fonts/,BoldFont=PlusJakartaSans-Bold.ttf]
% Maths en sans-serif (Fira Math) avec chiffres et lettres droites de Plus
% Jakarta, pour que les formules se fondent dans le texte.
\usepackage[math-style=ISO,bold-style=ISO]{unicode-math}
\setmathfont{FiraMath-Regular.otf}[Path=../../../fonts/]
\setmathfont{PlusJakartaSans-Medium.ttf}[Path=../../../fonts/,range={up/num,up/{Latin,latin},\mathup}]
\usepackage{icomma} % virgule décimale sans espace en mode math
% Caractères mathématiques Unicode absents des polices de texte (Plus Jakarta,
% Archivo Black) : rendus par leur équivalent mathématique, en texte comme en
% formule — sinon ils s'affichent en carré vide (ex. « Arithmétique dans ℤ »).
\usepackage{newunicodechar}
\newunicodechar{ℝ}{\ensuremath{\mathbb{R}}}
\newunicodechar{ℤ}{\ensuremath{\mathbb{Z}}}
\newunicodechar{ℂ}{\ensuremath{\mathbb{C}}}
\newunicodechar{ℕ}{\ensuremath{\mathbb{N}}}
\newunicodechar{ℚ}{\ensuremath{\mathbb{Q}}}
\newunicodechar{𝔻}{\ensuremath{\mathbb{D}}}
\newunicodechar{α}{\ensuremath{\alpha}}
\newunicodechar{β}{\ensuremath{\beta}}
\newunicodechar{γ}{\ensuremath{\gamma}}
\newunicodechar{δ}{\ensuremath{\delta}}
\newunicodechar{ε}{\ensuremath{\varepsilon}}
\newunicodechar{θ}{\ensuremath{\theta}}
\newunicodechar{λ}{\ensuremath{\lambda}}
\newunicodechar{μ}{\ensuremath{\mu}}
\newunicodechar{σ}{\ensuremath{\sigma}}
\newunicodechar{τ}{\ensuremath{\tau}}
\newunicodechar{φ}{\ensuremath{\varphi}}
\newunicodechar{ψ}{\ensuremath{\psi}}
\newunicodechar{ω}{\ensuremath{\omega}}
\newunicodechar{Σ}{\ensuremath{\Sigma}}
% majuscules produites par \MakeUppercase dans les titres (α-aminés -> Α)
\newunicodechar{Α}{\ensuremath{\alpha}}
\newunicodechar{Β}{\ensuremath{\beta}}
\newunicodechar{Γ}{\ensuremath{\Gamma}}
\newunicodechar{Δ}{\ensuremath{\Delta}}
\newunicodechar{Ε}{\ensuremath{\varepsilon}}
\newunicodechar{Θ}{\ensuremath{\Theta}}
\newunicodechar{Λ}{\ensuremath{\Lambda}}
\newunicodechar{Μ}{\ensuremath{\mu}}
\newunicodechar{Τ}{\ensuremath{\tau}}
\newunicodechar{Φ}{\ensuremath{\Phi}}
\newunicodechar{Ψ}{\ensuremath{\Psi}}
\newunicodechar{Ω}{\ensuremath{\Omega}}
\newunicodechar{↦}{\ensuremath{\mapsto}}
\newunicodechar{∈}{\ensuremath{\in}}
\newunicodechar{∉}{\ensuremath{\notin}}
\newunicodechar{∧}{\ensuremath{\wedge}}
\newunicodechar{⇔}{\ensuremath{\Leftrightarrow}}
\newunicodechar{⟺}{\ensuremath{\Longleftrightarrow}}
\newunicodechar{⇒}{\ensuremath{\Rightarrow}}
\newunicodechar{⟹}{\ensuremath{\Longrightarrow}}
\newunicodechar{∘}{\ensuremath{\circ}}
\newunicodechar{∪}{\ensuremath{\cup}}
\newunicodechar{∩}{\ensuremath{\cap}}
\newunicodechar{≡}{\ensuremath{\equiv}}
\newunicodechar{⊂}{\ensuremath{\subset}}
\newunicodechar{⊥}{\ensuremath{\perp}}
\newunicodechar{∃}{\ensuremath{\exists}}
\newunicodechar{∀}{\ensuremath{\forall}}
\newunicodechar{∛}{\ensuremath{\sqrt[3]{\,}}}
\newunicodechar{∜}{\ensuremath{\sqrt[4]{\,}}}
\newunicodechar{⃗}{\ensuremath{{}^{\to}}}
\newunicodechar{ⁿ}{\ifmmode^{n}\else\textsuperscript{\ensuremath{n}}\fi}
\newunicodechar{ˣ}{\ifmmode^{x}\else\textsuperscript{\ensuremath{x}}\fi}
\newunicodechar{ᵇ}{\ifmmode^{b}\else\textsuperscript{\ensuremath{b}}\fi}
\newunicodechar{ʳ}{\ifmmode^{r}\else\textsuperscript{\ensuremath{r}}\fi}
\newunicodechar{ᵘ}{\ifmmode^{u}\else\textsuperscript{\ensuremath{u}}\fi}
\newunicodechar{ʸ}{\ifmmode^{y}\else\textsuperscript{\ensuremath{y}}\fi}
\newunicodechar{ᵃ}{\ifmmode^{a}\else\textsuperscript{\ensuremath{a}}\fi}
\newunicodechar{ᵉ}{\ifmmode^{e}\else\textsuperscript{\ensuremath{e}}\fi}
\newunicodechar{ᵏ}{\ifmmode^{k}\else\textsuperscript{\ensuremath{k}}\fi}
\newunicodechar{ᵖ}{\ifmmode^{p}\else\textsuperscript{\ensuremath{p}}\fi}
\newunicodechar{ᴺ}{\ifmmode^{N}\else\textsuperscript{\ensuremath{N}}\fi}
\newunicodechar{ᶜ}{\ifmmode^{c}\else\textsuperscript{\ensuremath{c}}\fi}
\newunicodechar{ʰ}{\ifmmode^{h}\else\textsuperscript{\ensuremath{h}}\fi}
\newunicodechar{ᵠ}{\ifmmode^{q}\else\textsuperscript{\ensuremath{q}}\fi}
\newunicodechar{ₙ}{\ifmmode_{n}\else\textsubscript{\ensuremath{n}}\fi}
\newunicodechar{ₐ}{\ifmmode_{a}\else\textsubscript{\ensuremath{a}}\fi}
\newunicodechar{ᵢ}{\ifmmode_{i}\else\textsubscript{\ensuremath{i}}\fi}
\newunicodechar{ᵣ}{\ifmmode_{r}\else\textsubscript{\ensuremath{r}}\fi}
\newunicodechar{ₖ}{\ifmmode_{k}\else\textsubscript{\ensuremath{k}}\fi}
\newunicodechar{ᵦ}{\ifmmode_{\beta}\else\textsubscript{\ensuremath{\beta}}\fi}
\newunicodechar{ₓ}{\ifmmode_{x}\else\textsubscript{\ensuremath{x}}\fi}
\newfontfamily\popdisplay{ArchivoBlack-Regular.ttf}[Path=../../../fonts/]
\newfontfamily\poplogo{SpaceGrotesk-Bold.ttf}[Path=../../../fonts/]
\newfontfamily\popsemi{PlusJakartaSans-SemiBold.ttf}[Path=../../../fonts/]
\newfontfamily\popxbold{PlusJakartaSans-ExtraBold.ttf}[Path=../../../fonts/]
\newfontfamily\popxboldls{PlusJakartaSans-ExtraBold.ttf}[Path=../../../fonts/,LetterSpace=3.0]

\setlist[enumerate]{leftmargin=1.7em,itemsep=5pt,topsep=4pt}
\setlist[itemize]{leftmargin=1.7em,itemsep=4pt,topsep=4pt,label={\color{poporange}\textbullet}}
\setlength{\parindent}{0pt}
\setlength{\parskip}{5pt}
\renewcommand{\arraystretch}{1.25}

% pgfplots : style Pop par défaut
\pgfplotsset{
  every axis/.append style={axis line style={popink,line width=1pt},tick style={popink},
    grid style={popline,line width=0.5pt},label style={font=\small},tick label style={font=\footnotesize}},
  popcurve/.style={poporange,line width=1.8pt,no markers,smooth},
}
% Même style disponible dans un \draw TikZ simple (hors pgfplots)
\tikzset{popcurve/.style={poporange,line width=1.8pt}}

% ============================================================
% Pill / squiggle
% ============================================================
\newcommand{\poppill}[3]{%
  \tikz[baseline=-0.55ex]\node[draw=popink,line width=1.2pt,rounded corners=2.6pt,%
    fill=#2,inner xsep=6.5pt,inner ysep=3pt]{%
    \popxboldls\fontsize{7}{7}\selectfont\color{#3}\MakeUppercase{#1}};%
}
\newcommand{\popsquiggle}[1]{%
  \tikz[baseline=-0.4ex]{
    \draw[#1,line width=2.6pt,line cap=round]
      plot [smooth] coordinates {(0,0.09) (0.34,0.01) (0.68,0.21) (1.02,0.01) (1.36,0.10)};
  }%
}
% Mot-clé surligné (marqueur orange sous le texte)
\newcommand{\cle}[1]{{\popxbold\color{popdark}#1}}

% ============================================================
% En-tête / pied
% ============================================================
\pagestyle{fancy}
\fancyhf{}
\renewcommand{\headrulewidth}{0pt}
\renewcommand{\footrulewidth}{0pt}
\lhead{\raisebox{-0.15ex}{\poplogo\fontsize{13}{13}\selectfont\color{popink}OnBuch\color{poporange}+}}
\rhead{\poppill{\DOCMATIERE\ \textbullet\ \DOCNIVEAU\ \textbullet\ Leçon \DOCLECON}{white}{popink}}
\lfoot{\popxbold\fontsize{7.6}{7.6}\selectfont\color{popmuted}\MakeUppercase{Cours signé OnBuch+}\ \textbullet\ {\popsemi\DOCTITRE}}
\rfoot{\tikz[baseline=-0.6ex]\node[rounded corners=8.5pt,fill=popink,minimum height=17pt,minimum width=17pt,inner xsep=5pt,inner sep=0]{\popxbold\fontsize{8}{8}\selectfont\color{popcream}\thepage\,/\,\pageref*{LastPage}};}

% ============================================================
% Titres
% ============================================================
\newcommand{\chaptitre}[2]{%
  \begin{tcolorbox}[enhanced,colback=poporange,colframe=popink,boxrule=2.6pt,arc=11pt,
      drop shadow=popink,left=15pt,right=15pt,top=12pt,bottom=11pt,before skip=3pt,after skip=18pt]
    \poppill{#2}{popink}{popcream}\par\vspace{9pt}
    {\raggedright\hyphenpenalty=10000\exhyphenpenalty=10000\popdisplay\fontsize{22}{24}\selectfont\color{popcream}\MakeUppercase{#1}\par}
  \end{tcolorbox}
}
% Section numérotée : pastille orange avec le numéro + titre Archivo
\newcounter{popsec}
\newcommand{\coursec}[1]{%
  \stepcounter{popsec}\setcounter{popsub}{0}%
  \par\vspace{16pt}\noindent
  \begin{minipage}{\linewidth}
  \tikz[baseline=-0.7ex]\node[draw=popink,line width=1.6pt,fill=poporange,rounded corners=4pt,minimum size=22pt,inner sep=0]{\popdisplay\fontsize{12}{12}\selectfont\color{popcream}\thepopsec};%
  \hspace{8pt}{\popdisplay\fontsize{15}{17}\selectfont\color{popink}\MakeUppercase{#1}}\par
  \vspace{4pt}\noindent{\color{poporange}\rule{36pt}{3.6pt}}
  \end{minipage}\par\nopagebreak\vspace{8pt}
}
\newcounter{popsub}
\newcommand{\courssub}[1]{%
  \stepcounter{popsub}%
  \par\vspace{9pt}\noindent{\popxbold\fontsize{11.5}{13}\selectfont\color{popink}{\color{poporange}\thepopsec.\thepopsub}\hspace{6pt}#1}\par\nopagebreak\vspace{4pt}
}

% ============================================================
% Boîtes pédagogiques — même recette : fond blanc, bordure épaisse colorée,
% ombre dure encre, badge plein. Titre optionnel [#1].
% ============================================================
\tcbset{popbase/.style={breakable,enhanced,colback=white,boxrule=2.3pt,arc=9pt,
  shadow={3pt}{-3pt}{0pt}{popink},left=14pt,right=14pt,top=11pt,bottom=11pt,before skip=11pt,after skip=15pt,
  fontupper=\popsemi\fontsize{10.6}{14.5}\selectfont\color{popink},
  fontlower=\popsemi\fontsize{10.6}{14.5}\selectfont\color{popink}}}
% Coche dessinée (le glyphe ✓ n'existe pas dans les polices du document)
\newcommand{\popcheck}[1]{\tikz[baseline=-0.5ex]\draw[#1,line width=1.6pt,line cap=round,line join=round](-0.13,0.0)--(-0.04,-0.09)--(0.13,0.1);}
\newcommand{\popboxhead}[3]{\poppill{#1}{#2}{white}\ifx\relax#3\relax\else\hspace{6pt}{\popxbold\fontsize{10.6}{12}\selectfont\color{popink}#3}\fi\par\vspace{7pt}}

\newtcolorbox{aretenir}[1][]{popbase,colframe=popgreen,before upper={\popboxhead{À retenir}{popgreen}{#1}}}
\newtcolorbox{exemplebox}[1][]{popbase,colframe=poporange,before upper={\popboxhead{Exemple}{poporange}{#1}}}
\newtcolorbox{definition}[1][]{popbase,colframe=popblue,before upper={\popboxhead{Définition}{popblue}{#1}}}
\newtcolorbox{propriete}[1][]{popbase,colframe=poppurple,before upper={\popboxhead{Propriété}{poppurple}{#1}}}
\newtcolorbox{methode}[1][]{popbase,colframe=popgold,before upper={\popboxhead{Méthode}{popgold}{#1}}}
\newtcolorbox{attention}[1][]{popbase,colframe=poppink,before upper={\popboxhead{Attention}{poppink}{#1}}}
\newtcolorbox{experience}[1][]{popbase,colframe=popblue,colback=popblueL,before upper={\popboxhead{Expérience}{popblue}{#1}}}
% « Le savais-tu ? » : carte violette pleine (contexte, culture, Cameroun)
\newtcolorbox{savaistu}[1][]{popbase,colback=poppurple,colframe=popink,
  fontupper=\popsemi\fontsize{10.4}{14.2}\selectfont\color{white},
  before upper={\poppill{Le savais-tu\,?}{white}{poppurple}\ifx\relax#1\relax\else\hspace{6pt}{\popxbold\color{white}#1}\fi\par\vspace{7pt}}}
% Exercice résolu : énoncé en haut, \tcblower puis la solution sur fond crème
\newcounter{popexo}\newcounter{popexr}
\newtcolorbox{exoresolu}[1][]{popbase,colframe=poporange,colbacklower=popcream,
  segmentation style={popink,line width=1.2pt,dashed},
  before upper={\stepcounter{popexr}\popboxhead{Exercice résolu \thepopexr}{poporange}{#1}},
  before lower={\poppill{Solution}{popink}{popcream}\par\vspace{6pt}}}
% Exercice (non résolu en place) — la correction est dans la partie Corrigés
\newtcolorbox{exercice}[1][]{popbase,colframe=popink,
  before upper={\stepcounter{popexo}\popboxhead{Exercice \thepopexo}{popink}{#1}}}
% Corrigé
\newtcolorbox{corrige}[1][]{popbase,colframe=popgreen,colback=popgreenL,
  before upper={\popboxhead{Corrigé}{popgreen}{#1}}}
% Objectifs / prérequis / bilan
\newtcolorbox{objectifs}{popbase,colframe=popink,colback=poporangeL,
  before upper={\popboxhead{Objectifs de la leçon}{popink}{}}}
\newtcolorbox{prerequis}{popbase,colframe=popink,colback=popblueL,
  before upper={\popboxhead{Prérequis}{popblue}{}}}
\newtcolorbox{bilan}{popbase,colframe=popink,colback=white,boxrule=2.8pt,
  before upper={\popboxhead{Fiche bilan}{poporange}{L'essentiel à connaître pour l'examen}}}
% Figure encadrée avec légende
\newtcolorbox{popfigure}[1][]{enhanced,breakable=false,colback=white,colframe=popline,boxrule=1.4pt,arc=8pt,
  left=10pt,right=10pt,top=10pt,bottom=8pt,before skip=10pt,after skip=12pt,halign=center,
  lower separated=false,halign lower=center,
  fontlower=\popsemi\fontsize{9}{11.5}\selectfont\color{popmuted},
  #1}
\newcommand{\legende}[1]{\par\vspace{6pt}{\popsemi\fontsize{9}{11.5}\selectfont\color{popmuted}#1}}

% ============================================================
% Signature OnBuch+
% ============================================================
\newcommand{\poplogoinline}[1]{{\poplogo\fontsize{#1}{#1}\selectfont\color{popink}OnBuch\color{poporange}+}}
\newcommand{\signatureonbuch}{%
  \begin{tcolorbox}[enhanced,colback=popink,colframe=popink,boxrule=2.2pt,arc=10pt,
      drop shadow=poporange,left=14pt,right=14pt,top=10pt,bottom=10pt,before skip=0pt,after skip=0pt]
    \tikz[baseline=-0.7ex]\node[circle,fill=popgreen,draw=popcream,line width=1.3pt,minimum size=22pt,inner sep=0]{\popcheck{white}};%
    \hspace{9pt}\begin{minipage}[c]{0.62\linewidth}
      {\popxbold\fontsize{12}{13}\selectfont\color{popcream}Signé\ \textbullet\ {\poplogo OnBuch\color{poporange}+}}\par\vspace{2pt}
      {\raggedright\popsemi\fontsize{8}{10.5}\selectfont\color{popline}\MakeUppercase{Équipe pédagogique OnBuch+}\\\MakeUppercase{Conforme au programme officiel MINESEC}\par}
    \end{minipage}\hfill
    {\poplogo\fontsize{22}{22}\selectfont\color{popcream}OnBuch\color{poporange}+}
  \end{tcolorbox}%
}

% ============================================================
% Page de garde
% #1 titre · #2 matière · #3 niveau · #4 module · #5 n° leçon · #6 durée ·
% #7 description / objectifs (texte ou itemize)
% ============================================================
\newcommand{\pagegarde}[7]{%
  \thispagestyle{empty}
  \newgeometry{a4paper,margin=1.7cm,top=1.6cm,bottom=1.4cm}%
  \begin{tikzpicture}[remember picture,overlay]
    \fill[poporange!18] ([xshift=-3.2cm,yshift=-3.6cm]current page.north east) circle (4.6cm);
    \fill[poppurple!14] ([xshift=2.4cm,yshift=7.6cm]current page.south west) circle (3.8cm);
  \end{tikzpicture}%
  {\poplogo\fontsize{30}{30}\selectfont\color{popink}OnBuch\color{poporange}+}\hfill
  \raisebox{0.6ex}{\poppill{Cours signé \textbullet\ Premium}{popink}{popcream}}\par
  \vspace{1pt}\popsquiggle{poppurple}\par
  \vspace{8pt}
  \begin{tcolorbox}[enhanced,colback=poporange,colframe=popink,boxrule=2.8pt,arc=13pt,
      drop shadow=popink,left=18pt,right=18pt,top=12pt,bottom=13pt,after skip=10pt]
    \poppill{#2 \textbullet\ #3}{popink}{popcream}\hspace{5pt}\poppill{#4}{white}{popink}\par\vspace{8pt}
    {\popdisplay\fontsize{14}{14}\selectfont\color{popink}LEÇON #5}\par\vspace{4pt}
    {\raggedright\hyphenpenalty=10000\exhyphenpenalty=10000\popdisplay\fontsize{25}{27}\selectfont\color{popcream}\MakeUppercase{#1}\par}
    \vspace{8pt}
    {\popxbold\fontsize{9.5}{9.5}\selectfont\color{popink}\MakeUppercase{Durée conseillée : #6}}
  \end{tcolorbox}
  \begin{tcolorbox}[enhanced,colback=white,colframe=popink,boxrule=2.2pt,arc=10pt,
      drop shadow=popink,left=16pt,right=16pt,top=10pt,bottom=10pt,after skip=8pt]
    \poppill{Dans cette leçon}{poppurple}{white}\par\vspace{5pt}
    \popsemi\fontsize{9.3}{12.6}\selectfont\color{popink}#7
  \end{tcolorbox}
  \noindent\begin{minipage}[t]{0.487\linewidth}
    \begin{tcolorbox}[enhanced,colback=popgreen,colframe=popink,boxrule=2.2pt,arc=10pt,
        drop shadow=popink,left=11pt,right=11pt,top=10pt,bottom=10pt,height=3.1cm,valign=center]
      \tcbox[colback=white,colframe=popink,boxrule=1.3pt,arc=5pt,boxsep=0pt,nobeforeafter,
        left=3pt,right=3pt,top=3pt,bottom=3pt,valign=center]{\qrcode[height=1.9cm]{https://play.google.com/store/apps/details?id=com.luvvix.onbuch&hl=fr}}%
      \hspace{8pt}\begin{minipage}[c]{0.5\linewidth}
        {\popdisplay\fontsize{11}{12.5}\selectfont\color{white}\MakeUppercase{Télécharge OnBuch}}\par\vspace{4pt}
        {\raggedright\popsemi\fontsize{8.4}{11}\selectfont\color{white}Gratuit \textbullet\ Google Play\\Tuteur IA, annales, concours, résultats\par}
      \end{minipage}
    \end{tcolorbox}
  \end{minipage}\hfill
  \begin{minipage}[t]{0.487\linewidth}
    \begin{tcolorbox}[enhanced,colback=poppurple,colframe=popink,boxrule=2.2pt,arc=10pt,
        drop shadow=popink,left=11pt,right=11pt,top=10pt,bottom=10pt,height=3.1cm,valign=center]
      \tikz[baseline=-0.7ex]\node[circle,fill=white,draw=popink,line width=1.3pt,minimum size=1.6cm,inner sep=0]{\popdisplay\fontsize{24}{24}\selectfont\color{poppurple}+};%
      \hspace{8pt}\begin{minipage}[c]{0.6\linewidth}
        {\popdisplay\fontsize{11}{12.5}\selectfont\color{white}\MakeUppercase{Passe à OnBuch+}}\par\vspace{4pt}
        {\raggedright\popsemi\fontsize{8.4}{11}\selectfont\color{white}Tous les cours signés, Tuteur IA illimité, fiches PDF — onglet OnBuch+ de l'app\par}
      \end{minipage}
    \end{tcolorbox}
  \end{minipage}
  \vspace{8pt}
  \popcontenu
  \vfill
  \signatureonbuch
  \restoregeometry
  \newpage
}

% Rangée « Ce cours contient » (page de garde)
\newcommand{\poptile}[3]{% #1 couleur #2 gros texte #3 légende
  \begin{tcolorbox}[enhanced,colback=#1,colframe=popink,boxrule=2pt,arc=9pt,shadow={2.5pt}{-2.5pt}{0pt}{popink},
      left=6pt,right=6pt,top=9pt,bottom=9pt,halign=center,valign=center,height=1.75cm,width=\linewidth,nobeforeafter]
    {\popdisplay\fontsize{12.5}{13}\selectfont\color{white}\MakeUppercase{#2}}\par\vspace{3pt}
    {\centering\popsemi\fontsize{7.8}{9.5}\selectfont\color{white}#3\par}
  \end{tcolorbox}}
\newcommand{\popcontenu}{%
  \par\noindent{\popxboldls\fontsize{7.5}{7.5}\selectfont\color{popmuted}CE COURS SIGNÉ CONTIENT}\par\vspace{7pt}
  \noindent\begin{minipage}[t]{0.235\linewidth}\poptile{popblue}{Cours}{complet, illustré, pas à pas}\end{minipage}\hfill
  \begin{minipage}[t]{0.235\linewidth}\poptile{poporange}{Méthodes}{et exercices résolus}\end{minipage}\hfill
  \begin{minipage}[t]{0.235\linewidth}\poptile{poppink}{Exercices}{type examen + corrigés}\end{minipage}\hfill
  \begin{minipage}[t]{0.235\linewidth}\poptile{popgreen}{Fiche bilan}{l'essentiel à retenir}\end{minipage}\par}

% Sommaire
\newtcolorbox{sommaire}{enhanced,colback=white,colframe=popink,boxrule=2.2pt,arc=10pt,
  drop shadow=popink,left=18pt,right=18pt,top=13pt,bottom=13pt,before skip=6pt,after skip=20pt,
  before upper={\poppill{Sommaire}{poppurple}{white}\par\vspace{9pt}\popsemi\fontsize{10.6}{14.5}\selectfont\color{popink}}}

% Signature de fin de cours — poussée tout en bas de la dernière page
\newcommand{\findecours}{%
  \par\vfill\nopagebreak
  \begin{center}\popsquiggle{poporange}\end{center}\vspace{4pt}
  \signatureonbuch
}
\AtBeginDocument{\def\llbracket{\mathopen{[\mkern-3.5mu[}}\def\rrbracket{\mathclose{]\mkern-3.5mu]}}} % intervalles d'entiers (glyphe ⟦ absent de la police)
% Glyphes absents de Fira Math : redéfinis à partir de glyphes disponibles
\AtBeginDocument{%
  \def\setminus{\mathbin{\backslash}}\let\smallsetminus\setminus
  \def\vdots{\mathord{\vbox{\baselineskip3.2pt\lineskiplimit0pt\kern4pt\hbox{.}\hbox{.}\hbox{.}}}}%
  \def\ddots{\mathinner{\mkern1mu\raise6.5pt\hbox{.}\mkern2mu\raise3.6pt\hbox{.}\mkern2mu\raise0.7pt\hbox{.}\mkern1mu}}%
  \def\triangle{\mathord{\scalebox{0.9}{$\symup{\Delta}$}}}%
  \def\nabla{\mathord{\raisebox{\depth}{\scalebox{1}[-1]{$\symup{\Delta}$}}}}%
  \def\checkmark{\popcheck{popdark}}%
  \def\star{\mathbin{\ast}}\def\bigstar{\mathord{\ast}}%
  \def\diamond{\mathbin{\scalebox{0.6}{$\Diamond$}}}%
  \def\bigcup{\mathop{\mathchoice{\raisebox{-0.3em}{\scalebox{1.7}{$\cup$}}}{\scalebox{1.25}{$\cup$}}{\cup}{\cup}}\displaylimits}%
  \def\bigcap{\mathop{\mathchoice{\raisebox{-0.3em}{\scalebox{1.7}{$\cap$}}}{\scalebox{1.25}{$\cap$}}{\cap}{\cap}}\displaylimits}%
  \def\lesseqqgtr{\mathrel{\lessgtr}}\def\bigoplus{\mathop{\oplus}\displaylimits}%
}
\providecommand{\ssi}{si et seulement si} % abréviation fréquente
\AtBeginDocument{\providecommand{\wideparen}{\overparen}} % arc de cercle

%% FIGFIX-BEGIN v5 — correctifs globaux des figures (docs/onbuch-generation/tools/figfix)
\usepackage{adjustbox}\usepackage{etoolbox}\tcbuselibrary{raster}
% toute figure TikZ/pgfplots plus large que la ligne est réduite à la largeur disponible
\BeforeBeginEnvironment{tikzpicture}{\begin{adjustbox}{max width=\linewidth}}
\AfterEndEnvironment{tikzpicture}{\end{adjustbox}}
% légendes pgfplots lisibles par-dessus les courbes
\pgfplotsset{every axis/.append style={legend style={fill=white,fill opacity=0.92,text opacity=1,draw=black!15,inner sep=3pt}}}
% étiquettes d'arêtes (syntaxe edge["..."]) avec fond blanc
\tikzset{every edge quotes/.append style={fill=white,inner sep=1.2pt}}
%% FIGFIX-END
\begin{document}

\pagegarde{Espaces vectoriels sur ℝ}{Mathématiques}{1ère E}{Configurations et transformations élémentaires du plan}{N16}{10 h}{Cette leçon introduit la notion d'espace vectoriel sur ℝ, structure algébrique fondamentale qui unifie les vecteurs du plan, de l'espace et les ensembles ℝ² et ℝ³. Elle développe les outils essentiels — sous-espaces vectoriels, familles libres et génératrices, bases, dimension, applications linéaires, noyau et image — indispensables pour la suite du programme de Première et de Terminale.
\vspace{4pt}
\begin{multicols}{2}\small
\begin{itemize}
\item À la fin de cette leçon, je sais montrer qu'une loi de composition est externe sur un ensemble
\item À la fin de cette leçon, je sais montrer qu'un ensemble muni de deux lois est un espace vectoriel sur ℝ
\item À la fin de cette leçon, je sais montrer qu'une partie d'un espace vectoriel en est un sous-espace vectoriel
\item À la fin de cette leçon, je sais montrer qu'une famille finie de vecteurs est libre, liée ou génératrice
\item À la fin de cette leçon, je sais montrer qu'une famille est une base et déterminer la dimension d'un espace vectoriel
\item À la fin de cette leçon, je sais montrer qu'une application entre espaces vectoriels est linéaire
\end{itemize}\end{multicols}}

\begin{sommaire}
\begin{multicols}{2}
\begin{enumerate}
\item Lois de composition externe et espaces vectoriels sur ℝ
\item Sous-espaces vectoriels
\item Familles génératrices, familles libres, familles liées
\item Bases et dimension d'un espace vectoriel
\item Applications linéaires
\item Noyau et image d'une application linéaire
\item Activité d'intégration
\item Méthodes et exercices résolus
\item Exercices
\item Corrigés des exercices
\item Fiche bilan
\end{enumerate}
\end{multicols}
\end{sommaire}

\begin{objectifs}
\begin{itemize}
\item À la fin de cette leçon, je sais montrer qu'une loi de composition est externe sur un ensemble
\item À la fin de cette leçon, je sais montrer qu'un ensemble muni de deux lois est un espace vectoriel sur ℝ
\item À la fin de cette leçon, je sais montrer qu'une partie d'un espace vectoriel en est un sous-espace vectoriel
\item À la fin de cette leçon, je sais montrer qu'une famille finie de vecteurs est libre, liée ou génératrice
\item À la fin de cette leçon, je sais montrer qu'une famille est une base et déterminer la dimension d'un espace vectoriel
\item À la fin de cette leçon, je sais montrer qu'une application entre espaces vectoriels est linéaire
\item À la fin de cette leçon, je sais déterminer une équation caractéristique et une base du noyau d'une application linéaire
\item À la fin de cette leçon, je sais déterminer une équation caractéristique et une base de l'image d'une application linéaire
\item À la fin de cette leçon, je sais montrer qu'une application linéaire est bijective en exploitant son écriture analytique, son noyau ou son image
\end{itemize}
\end{objectifs}

\begin{prerequis}
\begin{itemize}
\item Vecteurs du plan et de l'espace : opérations, colinéarité, repères (classe de Seconde)
\item Résolution de systèmes linéaires à 2 ou 3 inconnues (méthode de substitution ou de combinaison)
\item Notion d'application, d'injectivité, de surjectivité et de bijectivité
\item Calculs dans ℝ² et ℝ³ : couples et triplets de réels, opérations coordonnée par coordonnée
\item Configurations et transformations élémentaires du plan (symétries, translations, homothéties)
\end{itemize}
\end{prerequis}

\newpage

\coursec{Lois de composition externe et espaces vectoriels sur $\mathbb{R}$}

À Douala, une entreprise de transport organise ses trajets : chaque trajet est un \emph{vecteur déplacement}. On peut additionner deux trajets successifs (aller de Bonabéri à Akwa, puis d'Akwa à Deïdo, revient à faire un seul trajet de Bonabéri à Deïdo) et multiplier un trajet par un coefficient pour ajuster les distances (deux fois le même trajet, la moitié d'un trajet, le trajet en sens inverse). À Yaoundé, un informaticien manipule des listes de nombres — coordonnées GPS, pixels d'une image — avec exactement les mêmes règles : additionner deux listes terme à terme, multiplier une liste par un réel. Comment unifier toutes ces situations sous une même théorie ? C'est toute l'idée des \cle{espaces vectoriels} : dégager la \emph{structure commune} qui permet de calculer, de décomposer et de transformer. Cette structure sera aussi la clé pour comprendre les transformations du plan étudiées en Terminale.

\courssub{Rappel : les lois de composition internes}

Tu connais déjà, sans le savoir, des lois de composition internes. Lorsque tu additionnes deux vecteurs du plan, le résultat est encore un vecteur du plan : l'addition « reste » dans l'ensemble.

\begin{definition}[Loi de composition interne]
Soit $E$ un ensemble non vide. Une \cle{loi de composition interne} sur $E$ est une application de $E \times E$ dans $E$. Autrement dit, à tout couple $(a,b)$ d'éléments de $E$, on associe un élément de $E$, noté $a \star b$.
\end{definition}

\begin{exemplebox}[Exemples de lois internes]
\begin{itemize}
\item L'addition des vecteurs du plan : $(\vec{u}, \vec{v}) \mapsto \vec{u} + \vec{v}$.
\item L'addition dans $\mathbb{R}^2$ : $(x,y) + (x',y') = (x+x',\, y+y')$.
\item La multiplication dans $\mathbb{R}$ : $(a,b) \mapsto ab$.
\end{itemize}
Dans chaque cas, on combine \emph{deux éléments du même ensemble} et le résultat \emph{reste dans cet ensemble}.
\end{exemplebox}

\courssub{Lois de composition externes}

Revenons à notre transporteur de Douala. Multiplier un trajet par $2$, par $0{,}5$ ou par $-1$, ce n'est pas combiner deux trajets : c'est faire agir un \emph{nombre réel} (un élément \emph{extérieur} à l'ensemble des vecteurs) sur un vecteur, pour obtenir un nouveau vecteur. C'est une loi de composition \emph{externe}.

\begin{definition}[Loi de composition externe]
Soit $E$ un ensemble non vide. Une \cle{loi de composition externe} sur $E$ à opérateurs dans $\mathbb{R}$ est une application de $\mathbb{R} \times E$ dans $E$ :
\[ \begin{array}{rcl} \mathbb{R} \times E & \longrightarrow & E \\ (\lambda, x) & \longmapsto & \lambda \cdot x \end{array} \]
Le réel $\lambda$ est appelé \emph{scalaire} ; l'élément $\lambda \cdot x$ doit appartenir à $E$.
\end{definition}

\begin{exemplebox}[Exemples fondamentaux]
\begin{itemize}
\item \textbf{Multiplication d'un vecteur du plan par un réel.} Si $\vec{u}$ est un vecteur du plan et $\lambda \in \mathbb{R}$, alors $\lambda \vec{u}$ est encore un vecteur du plan : même direction si $\lambda \neq 0$, même sens si $\lambda > 0$, sens contraire si $\lambda < 0$, et longueur multipliée par $|\lambda|$.
\item \textbf{Multiplication d'un triplet de $\mathbb{R}^3$ par un réel.} Pour $(x,y,z) \in \mathbb{R}^3$ et $\lambda \in \mathbb{R}$, on pose $\lambda \cdot (x,y,z) = (\lambda x, \lambda y, \lambda z)$, qui est bien un élément de $\mathbb{R}^3$.
\end{itemize}
\end{exemplebox}

\begin{popfigure}
\begin{center}
\begin{tikzpicture}[scale=1.1, >=Stealth]
\coordinate (O) at (0,0);
\coordinate (U) at (2,1);
\draw[->, thick, popblue] (O) -- (U) node[midway, above left] {$\vec{u}$};
\draw[->, thick, popgreen] (O) -- ($2*(U)$) node[midway, above left, pos=0.6] {$2\vec{u}$};
\draw[->, thick, poporange] (O) -- ($-1*(U)$) node[midway, below right] {$-\vec{u}$};
\draw[->, thick, poppurple] (O) -- ($0.5*(U)$) node[midway, below, pos=0.85] {$0.5\vec{u}$};
\fill[popdark] (O) circle (1.5pt) node[below left] {$O$};
\end{tikzpicture}
\end{center}
\legende{La loi externe en action : les multiples $2\vec{u}$, $-\vec{u}$ et $0.5\vec{u}$ d'un vecteur $\vec{u}$ du plan restent des vecteurs du plan.}
\end{popfigure}

\begin{attention}[Une loi externe doit « rester dans l'ensemble »]
Pour être une loi de composition externe sur $E$, l'application doit envoyer $\mathbb{R} \times E$ \textbf{dans} $E$. Contre-exemples :
\begin{itemize}
\item Sur $E = \mathbb{N}$, l'opération $(\lambda, n) \mapsto \lambda n$ n'est \textbf{pas} une loi externe : $0.5 \times 3 = 1.5 \notin \mathbb{N}$. L'ensemble d'arrivée est incorrect.
\item Sur $E = \mathbb{R}_+^*$, l'opération $(\lambda, x) \mapsto \lambda x$ n'est pas une loi externe : $(-2) \times 3 = -6 \notin \mathbb{R}_+^*$.
\end{itemize}
Vérifier que le résultat appartient à $E$ est la première chose à faire, et c'est souvent là que se cache le piège au Probatoire.
\end{attention}

\begin{methode}[Montrer qu'une loi est externe]
\begin{enumerate}
\item Identifier l'ensemble $E$ et l'opération proposée $(\lambda, x) \mapsto \lambda \cdot x$.
\item Prendre un réel $\lambda$ \emph{quelconque} et un élément $x$ \emph{quelconque} de $E$.
\item Calculer $\lambda \cdot x$ et vérifier que le résultat satisfait la condition d'appartenance à $E$.
\item Conclure : « la loi est une loi de composition externe sur $E$ ».
\end{enumerate}
\end{methode}

\courssub{Espaces vectoriels sur $\mathbb{R}$ : définition}

Nous avons maintenant les deux ingrédients : une addition interne et une multiplication par les réels. Mais il ne suffit pas de les juxtaposer : il faut qu'elles « coopèrent » selon des règles précises, celles que tu utilises spontanément depuis la Seconde avec les vecteurs. Ces règles sont les huit axiomes.

\begin{definition}[Espace vectoriel sur $\mathbb{R}$]
Soit $E$ un ensemble non vide muni :
\begin{itemize}
\item d'une loi de composition interne, l'\emph{addition} : $(x,y) \mapsto x + y$ ;
\item d'une loi de composition externe : $(\lambda, x) \mapsto \lambda \cdot x$.
\end{itemize}
On dit que $(E, +, \cdot)$ est un \cle{espace vectoriel sur $\mathbb{R}$} si les huit axiomes suivants sont vérifiés, pour tous $x, y, z \in E$ et tous $\lambda, \mu \in \mathbb{R}$ :
\begin{enumerate}
\item \textbf{Associativité de $+$ :} $x + (y + z) = (x + y) + z$ ;
\item \textbf{Commutativité de $+$ :} $x + y = y + x$ ;
\item \textbf{Élément neutre :} il existe $0_E \in E$ tel que $x + 0_E = x$ ;
\item \textbf{Opposé :} pour tout $x \in E$, il existe $x' \in E$ tel que $x + x' = 0_E$ (on note $x' = -x$) ;
\item \textbf{Distributivité par rapport aux scalaires :} $(\lambda + \mu) \cdot x = \lambda \cdot x + \mu \cdot x$ ;
\item \textbf{Distributivité par rapport aux vecteurs :} $\lambda \cdot (x + y) = \lambda \cdot x + \lambda \cdot y$ ;
\item \textbf{Compatibilité :} $\lambda \cdot (\mu \cdot x) = (\lambda \mu) \cdot x$ ;
\item \textbf{Unité :} $1 \cdot x = x$.
\end{enumerate}
Les éléments de $E$ sont appelés \emph{vecteurs}, ceux de $\mathbb{R}$ les \emph{scalaires}.
\end{definition}

\begin{popfigure}
\begin{center}
\begin{tikzpicture}[
  grow=down,
  level 1/.style={sibling distance=4.2cm, level distance=1.6cm},
  level 2/.style={sibling distance=2.6cm, level distance=1.5cm},
  every node/.style={draw, rounded corners, align=center, font=\small, fill=popcream},
  edge from parent/.style={draw=popmuted, thick}]
\node[fill=popblueL, font=\small\bfseries] {Espace vectoriel $(E,+,\cdot)$}
  child { node[fill=popgreenL, align=center] {Ensemble $E$\\ non vide} }
  child { node[fill=poporangeL, align=center] {Loi interne $+$\\ $E\times E \to E$}
    child { node[align=center] {4 axiomes\\ (groupe\\ commutatif)} } }
  child { node[fill=poppinkL, align=center] {Loi externe $\cdot$\\ $\mathbb{R}\times E \to E$}
    child { node[align=center] {4 axiomes\\ (distributivités,\\ compatibilité, $1\cdot x = x$)} } };
\end{tikzpicture}
\end{center}
\legende{Carte de la définition : un espace vectoriel, c'est un ensemble muni de deux lois vérifiant huit axiomes.}
\end{popfigure}

\begin{exemplebox}[Les exemples fondamentaux à connaître]
Munis de leurs opérations usuelles (addition et multiplication par un réel, composante par composante), sont des espaces vectoriels sur $\mathbb{R}$ :
\begin{itemize}
\item $\mathbb{R}$ lui-même ;
\item $\mathbb{R}^2 = \{(x,y) \mid x \in \mathbb{R},\ y \in \mathbb{R}\}$ ;
\item $\mathbb{R}^3 = \{(x,y,z) \mid x,y,z \in \mathbb{R}\}$ ;
\item l'ensemble des vecteurs du plan ;
\item l'ensemble des vecteurs de l'espace.
\end{itemize}
Ce sont les \emph{seuls} espaces vectoriels que nous étudierons cette année (conformément au programme).
\end{exemplebox}

\courssub{Vérification complète sur un exemple chiffré : $\mathbb{R}^2$}

\begin{exemplebox}[$\mathbb{R}^2$ est un espace vectoriel sur $\mathbb{R}$]
Munissons $\mathbb{R}^2$ de $(x,y) + (x',y') = (x+x',\, y+y')$ et $\lambda \cdot (x,y) = (\lambda x, \lambda y)$. Vérifions les axiomes avec $u = (x,y)$, $v = (x',y')$, $w = (x'',y'')$ et $\lambda, \mu \in \mathbb{R}$.

\textbf{Lois bien définies.} La somme de deux couples de réels est un couple de réels, et le produit d'un couple par un réel est un couple de réels : les deux lois sont bien interne et externe.

\textbf{Axiome 1 (associativité).} $(u+v)+w = (x+x'+x'',\ y+y'+y'') = u+(v+w)$, car l'addition des réels est associative.

\textbf{Axiome 2 (commutativité).} $u+v = (x+x',\ y+y') = (x'+x,\ y'+y) = v+u$.

\textbf{Axiome 3 (neutre).} $0_{\mathbb{R}^2} = (0,0)$ convient : $(x,y) + (0,0) = (x,y)$.

\textbf{Axiome 4 (opposé).} L'opposé de $(x,y)$ est $(-x,-y) \in \mathbb{R}^2$ : $(x,y) + (-x,-y) = (0,0)$.

\textbf{Axiome 5.} $(\lambda+\mu)\cdot(x,y) = \big((\lambda+\mu)x,\ (\lambda+\mu)y\big) = (\lambda x, \lambda y) + (\mu x, \mu y) = \lambda u + \mu u$.

\textbf{Axiome 6.} $\lambda(u+v) = \big(\lambda(x+x'),\ \lambda(y+y')\big) = (\lambda x, \lambda y) + (\lambda x', \lambda y') = \lambda u + \lambda v$.

\textbf{Axiome 7.} $\lambda(\mu u) = \lambda(\mu x, \mu y) = (\lambda\mu x, \lambda\mu y) = (\lambda\mu)u$.

\textbf{Axiome 8.} $1 \cdot (x,y) = (1 \times x, 1 \times y) = (x,y)$.

Les huit axiomes sont vérifiés : $\mathbb{R}^2$ est un espace vectoriel sur $\mathbb{R}$.
\end{exemplebox}

\begin{methode}[Montrer qu'un ensemble est un espace vectoriel]
\begin{enumerate}
\item Vérifier que l'addition est une loi \emph{interne} et la multiplication par un réel une loi \emph{externe} (le résultat reste dans l'ensemble).
\item Vérifier les huit axiomes, ou — méthode plus rapide que nous verrons dans la section suivante — reconnaître l'ensemble comme \emph{sous-espace vectoriel} d'un espace vectoriel connu ($\mathbb{R}$, $\mathbb{R}^2$, $\mathbb{R}^3$, vecteurs du plan ou de l'espace).
\item Rédiger la conclusion proprement : « $(E,+,\cdot)$ est un espace vectoriel sur $\mathbb{R}$ ».
\end{enumerate}
\end{methode}

\courssub{Règles de calcul dans un espace vectoriel}

Les axiomes ont des conséquences pratiques que tu utiliseras constamment.

\begin{propriete}[Règles de calcul]
Dans un espace vectoriel $E$ sur $\mathbb{R}$, pour tout $x \in E$ et tout $\lambda \in \mathbb{R}$ :
\begin{enumerate}
\item $0 \cdot x = 0_E$ ;
\item $\lambda \cdot 0_E = 0_E$ ;
\item $(-\lambda) \cdot x = -(\lambda \cdot x)$ ;
\item $\lambda \cdot x = 0_E \implies \lambda = 0 \text{ ou } x = 0_E$.
\end{enumerate}
\end{propriete}

\begin{experience}[Démonstration guidée de la règle 1]
Montrons que $0 \cdot x = 0_E$ pour tout $x \in E$.
\begin{enumerate}
\item Écris $0 \cdot x = (0+0) \cdot x$ (car $0 = 0+0$ dans $\mathbb{R}$).
\item Applique l'axiome 5 (distributivité) : $0 \cdot x = 0 \cdot x + 0 \cdot x$.
\item Ajoute l'opposé de $0 \cdot x$ aux deux membres : $0_E = 0 \cdot x$.
\end{enumerate}
La règle 4 se montre ainsi : si $\lambda \neq 0$, alors $\lambda$ admet un inverse $\frac{1}{\lambda}$ ; en multipliant $\lambda \cdot x = 0_E$ par $\frac{1}{\lambda}$ et en utilisant les axiomes 6, 7 et 8 (car $\frac{1}{\lambda} \cdot 0_E = 0_E$ par l'axiome 6), on obtient $x = \frac{1}{\lambda} \cdot 0_E = 0_E$.
\end{experience}

\begin{attention}[Ne pas confondre $0$ et $0_E$]
Le réel $0$ (scalaire) et le vecteur nul $0_E$ sont des objets de nature différente. Dans $\mathbb{R}^2$, $0_E = (0,0)$ : ce n'est pas le nombre $0$ ! De même, l'égalité $\lambda \cdot x = 0_E$ ne signifie pas automatiquement que $x = 0_E$ : c'est $\lambda = 0$ \textbf{ou} $x = 0_E$. Soigne cette distinction dans tes rédactions : les correcteurs du Probatoire y sont attentifs.
\end{attention}

\begin{exoresolu}[La loi proposée est-elle externe ?]
\textbf{Énoncé.} On munit $E = \mathbb{R}^2$ de l'opération $\lambda \star (x,y) = (\lambda x,\ y)$. Cette opération est-elle une loi de composition externe ? L'ensemble $(E, +, \star)$, où $+$ est l'addition usuelle, est-il un espace vectoriel sur $\mathbb{R}$ ?
\tcblower
\textbf{Solution.} Pour tout $\lambda \in \mathbb{R}$ et tout $(x,y) \in \mathbb{R}^2$, $\lambda \star (x,y) = (\lambda x, y)$ est bien un couple de réels : l'opération envoie $\mathbb{R} \times E$ dans $E$, c'est donc une loi de composition externe sur $E$.

Cependant, testons l'axiome 5 avec $\lambda = \mu = 1$ et $(x,y) = (0,1)$ :
\[ (1+1) \star (0,1) = 2 \star (0,1) = (0,1), \quad \text{alors que} \quad 1\star(0,1) + 1\star(0,1) = (0,1)+(0,1) = (0,2). \]
Or $(0,1) \neq (0,2)$ : l'axiome 5 est en défaut. Donc $(E, +, \star)$ \textbf{n'est pas} un espace vectoriel sur $\mathbb{R}$, bien que la loi soit externe. Être une loi externe est une condition nécessaire, mais pas suffisante : il faut les huit axiomes.
\end{exoresolu}

\begin{savaistu}[Giuseppe Peano, le père des espaces vectoriels]
La notion d'espace vectoriel a été formalisée en 1888 par le mathématicien italien \emph{Giuseppe Peano} (1858--1932), dans son ouvrage \emph{Calcolo geometrico}. Son idée révolutionnaire : au lieu d'étudier séparément les vecteurs géométriques, les nombres, les polynômes, on définit une \emph{structure abstraite} par des axiomes, et tout théorème démontré une fois vaut pour tous les exemples à la fois. C'est exactement la démarche que tu suis dans cette leçon : ce que tu prouves pour les vecteurs du plan sera vrai pour $\mathbb{R}^3$ sans effort supplémentaire. Peano est aussi célèbre pour ses axiomes de l'arithmétique et pour avoir créé une langue internationale, le \emph{latino sine flexione}.
\end{savaistu}

\begin{aretenir}[L'essentiel de la section]
\begin{itemize}
\item Une \cle{loi de composition externe} sur $E$ est une application de $\mathbb{R} \times E$ \textbf{dans} $E$ : $(\lambda, x) \mapsto \lambda \cdot x$.
\item Un \cle{espace vectoriel sur $\mathbb{R}$} est un ensemble muni d'une addition interne et d'une multiplication externe vérifiant les \textbf{huit axiomes}.
\item Exemples de référence : $\mathbb{R}$, $\mathbb{R}^2$, $\mathbb{R}^3$, les vecteurs du plan, les vecteurs de l'espace.
\item Règles de calcul : $0 \cdot x = 0_E$, $\lambda \cdot 0_E = 0_E$, $(-\lambda) \cdot x = -(\lambda \cdot x)$, et $\lambda \cdot x = 0_E \Rightarrow \lambda = 0 \text{ ou } x = 0_E$.
\end{itemize}
\end{aretenir}

\coursec{Sous-espaces vectoriels}

Dans la section précédente, nous avons défini les espaces vectoriels sur $\mathbb{R}$ : un ensemble $E$ muni d'une loi de composition interne $+$ et d'une loi de composition externe $\cdot$ (multiplication par un réel), vérifiant huit axiomes. Nous avons rencontré les exemples fondamentaux : $\mathbb{R}$, $\mathbb{R}^2$, $\mathbb{R}^3$, l'ensemble $\mathcal{V}$ des vecteurs du plan.

Une question se pose naturellement : à l'intérieur d'un espace vectoriel $E$, certaines parties $F \subset E$ héritent-elles automatiquement de la structure d'espace vectoriel ? Par exemple, dans le plan $\mathbb{R}^2$, une droite passant par l'origine se comporte-t-elle, à elle seule, comme un « petit espace vectoriel » ? La réponse est oui, et c'est toute la notion de \cle{sous-espace vectoriel}. L'intérêt pratique est immense : au lieu de revérifier les huit axiomes (fastidieux !), il suffira de vérifier \textbf{trois conditions simples}.

\courssub{Stabilité d'une partie par une loi}

Avant de définir les sous-espaces vectoriels, précisons le vocabulaire de la \cle{stabilité}, que tu as déjà croisé avec les lois de composition internes.

\begin{definition}[Partie stable par une loi]
Soit $E$ un ensemble muni d'une loi de composition interne $\star$ et $F$ une partie de $E$. On dit que $F$ est \cle{stable} par la loi $\star$ lorsque :
\[
\forall x \in F,\ \forall y \in F,\quad x \star y \in F.
\]
Soit $E$ muni d'une loi de composition externe $\cdot$ sur $\mathbb{R}$. On dit que $F$ est stable par cette loi externe lorsque :
\[
\forall \lambda \in \mathbb{R},\ \forall x \in F,\quad \lambda \cdot x \in F.
\]
\end{definition}

\begin{exemplebox}[Vérifications de stabilité]
\begin{enumerate}
\item Dans $\mathbb{R}$, l'ensemble $\mathbb{N}$ est stable par $+$ (la somme de deux entiers naturels est un entier naturel), mais n'est \textbf{pas} stable par multiplication par un réel : $1 \in \mathbb{N}$ mais $\dfrac{1}{2} \times 1 = \num{0,5} \notin \mathbb{N}$.
\item Dans $\mathbb{R}^2$, l'ensemble $F = \{(x,y) \in \mathbb{R}^2 \mid y = 2x\}$ est stable par $+$ : si $(x_1, 2x_1)$ et $(x_2, 2x_2)$ sont dans $F$, leur somme $(x_1+x_2,\ 2(x_1+x_2))$ est encore dans $F$.
\item Dans $\mathbb{R}^2$, l'ensemble $G = \{(x,y) \in \mathbb{R}^2 \mid x \geq 0\}$ est stable par $+$ mais pas par multiplication par un réel : $(1,0) \in G$ mais $(-1)\cdot(1,0) = (-1,0) \notin G$.
\end{enumerate}
\end{exemplebox}

\courssub{Définition d'un sous-espace vectoriel}

Revenons à notre intuition géométrique. Dans le plan muni d'un repère $(O ; \vec{i}, \vec{j})$, considérons une droite $(d)$ passant par $O$. Si $\vec{u}$ et $\vec{v}$ sont deux vecteurs portés par $(d)$, alors $\vec{u} + \vec{v}$ est encore porté par $(d)$, et $\lambda \vec{u}$ aussi. La droite $(d)$ est donc « fermée » pour les deux opérations : c'est un sous-espace vectoriel du plan. En revanche, une droite ne passant pas par $O$ ne contient pas le vecteur nul : elle ne peut pas être un sous-espace vectoriel.

\begin{definition}[Sous-espace vectoriel]
Soit $(E, +, \cdot)$ un espace vectoriel sur $\mathbb{R}$ et $F$ une partie de $E$. On dit que $F$ est un \cle{sous-espace vectoriel} de $E$ lorsque :
\begin{enumerate}
\item $F$ est \textbf{non vide} ;
\item $F$ est \textbf{stable par addition} : pour tous $x, y \in F$, on a $x + y \in F$ ;
\item $F$ est \textbf{stable par multiplication par un réel} : pour tout $\lambda \in \mathbb{R}$ et tout $x \in F$, on a $\lambda \cdot x \in F$.
\end{enumerate}
\end{definition}

\begin{propriete}[Tout sous-espace vectoriel est un espace vectoriel]
Si $F$ est un sous-espace vectoriel de $E$, alors $(F, +, \cdot)$, muni des lois induites par celles de $E$, est lui-même un espace vectoriel sur $\mathbb{R}$. En particulier, $F$ contient le vecteur nul $0_E$ (prendre $\lambda = 0$ : $0 \cdot x = 0_E \in F$) et l'opposé de chacun de ses éléments (prendre $\lambda = -1$).
\end{propriete}

\begin{aretenir}[Intérêt pratique]
Pour montrer qu'un ensemble $F$ est un espace vectoriel, il est \textbf{beaucoup plus rapide} de montrer que $F$ est un sous-espace vectoriel d'un espace vectoriel connu ($\mathbb{R}$, $\mathbb{R}^2$, $\mathbb{R}^3$, l'ensemble des vecteurs du plan) que de vérifier les huit axiomes un par un. C'est la stratégie à adopter systématiquement au Probatoire.
\end{aretenir}

\begin{exemplebox}[Sous-espaces vectoriels évidents]
Tout espace vectoriel $E$ admet au moins deux sous-espaces vectoriels, dits \emph{triviaux} :
\begin{itemize}
\item $\{0_E\}$, le sous-espace nul ;
\item $E$ lui-même.
\end{itemize}
\end{exemplebox}

\courssub{Caractérisation pratique}

Les trois conditions de la définition peuvent être regroupées en une seule, ce qui allège considérablement les rédactions.

\begin{propriete}[Caractérisation d'un sous-espace vectoriel]
Soit $E$ un espace vectoriel sur $\mathbb{R}$ et $F \subset E$. Alors :
\[
F \text{ est un sous-espace vectoriel de } E \iff
\begin{cases}
F \neq \varnothing \quad (\text{en pratique : } 0_E \in F),\\
\forall x, y \in F,\ \forall \lambda \in \mathbb{R},\quad \lambda x + y \in F.
\end{cases}
\]
\end{propriete}

\begin{experience}[Démonstration guidée de la caractérisation]
\textbf{Sens direct.} Supposons que $F$ est un sous-espace vectoriel de $E$.
\begin{enumerate}
\item Par définition, $F \neq \varnothing$.
\item Soient $x, y \in F$ et $\lambda \in \mathbb{R}$. Par stabilité par la loi externe, $\lambda x \in F$. Par stabilité par addition, $\lambda x + y \in F$. La condition est vérifiée.
\end{enumerate}
\textbf{Sens réciproque.} Supposons que $F \neq \varnothing$ et que $\lambda x + y \in F$ pour tous $x, y \in F$ et tout $\lambda \in \mathbb{R}$.
\begin{enumerate}
\item \emph{Stabilité par addition} : soient $x, y \in F$. En choisissant $\lambda = 1$, on obtient $1 \cdot x + y = x + y \in F$.
\item \emph{Stabilité par la loi externe} : soient $x \in F$ et $\lambda \in \mathbb{R}$. Comme $F \neq \varnothing$, il existe $x_0 \in F$. En prenant $\lambda = -1$ et $y = x_0$ dans l'hypothèse avec $x=x_0$, on a $(-1)x_0 + x_0 = 0_E \in F$. Alors, pour $x \in F$ et $\lambda \in \mathbb{R}$ quelconques : $\lambda x + 0_E = \lambda x \in F$.
\end{enumerate}
Les deux conditions de la définition sont donc vérifiées : $F$ est un sous-espace vectoriel de $E$. \hfill$\blacksquare$
\end{experience}

\begin{methode}[Montrer qu'une partie est un sous-espace vectoriel en 3 étapes]
Pour montrer qu'une partie $F$ d'un espace vectoriel connu $E$ (souvent $\mathbb{R}^2$ ou $\mathbb{R}^3$) est un sous-espace vectoriel :
\begin{enumerate}
\item \textbf{Non vide} : vérifier que $0_E \in F$ (les coordonnées nulles satisfont-elles la condition définissant $F$ ?). Cela prouve $F \neq \varnothing$.
\item \textbf{Éléments arbitraires} : prendre deux éléments quelconques $x, y \in F$ (écrire explicitement ce que cela signifie : leurs coordonnées vérifient la condition) et un réel $\lambda$ quelconque.
\item \textbf{Combinaison} : calculer $\lambda x + y$ et vérifier que ses coordonnées satisfont encore la condition définissant $F$. Conclure.
\end{enumerate}
\end{methode}

\begin{attention}[Le piège du vecteur nul]
Erreur fréquente au Probatoire : oublier de vérifier que $F$ est non vide, ou pire, ne pas remarquer que $0_E \notin F$. Or \textbf{tout sous-espace vectoriel contient obligatoirement le vecteur nul}. Conséquence pratique redoutable : si $0_E \notin F$, alors $F$ n'est \textbf{pas} un sous-espace vectoriel, et c'est terminé — inutile de tester la stabilité ! C'est le moyen le plus rapide de montrer qu'une partie n'est \emph{pas} un sous-espace vectoriel.
\end{attention}

\courssub{Exemples dans $\mathbb{R}^2$ et $\mathbb{R}^3$}

\begin{exoresolu}[Une droite vectorielle de $\mathbb{R}^2$]
Montrer que $F = \{(x,y) \in \mathbb{R}^2 \mid 2x - 3y = 0\}$ est un sous-espace vectoriel de $\mathbb{R}^2$.
\tcblower
\textbf{Solution détaillée.} On applique la méthode en trois étapes.

\textbf{Étape 1 : $F$ est non vide.} Le vecteur nul $(0,0)$ vérifie $2 \times 0 - 3 \times 0 = 0$, donc $(0,0) \in F$ et $F \neq \varnothing$.

\textbf{Étape 2 : éléments arbitraires.} Soient $u = (x_1, y_1) \in F$ et $v = (x_2, y_2) \in F$, et soit $\lambda \in \mathbb{R}$. Par définition de $F$ :
\[
2x_1 - 3y_1 = 0 \quad \text{et} \quad 2x_2 - 3y_2 = 0.
\]

\textbf{Étape 3 : combinaison.} On calcule :
\[
\lambda u + v = (\lambda x_1 + x_2,\ \lambda y_1 + y_2).
\]
Vérifions que ce couple appartient à $F$ :
\begin{align*}
2(\lambda x_1 + x_2) - 3(\lambda y_1 + y_2) &= \lambda(2x_1 - 3y_1) + (2x_2 - 3y_2)\\
&= \lambda \times 0 + 0 = 0.
\end{align*}
Donc $\lambda u + v \in F$.

\textbf{Conclusion :} $F$ est un sous-espace vectoriel de $\mathbb{R}^2$. Géométriquement, $F$ est la droite vectorielle d'équation $y = \dfrac{2}{3}x$, passant par l'origine.
\end{exoresolu}

\begin{propriete}[Sous-espaces vectoriels de $\mathbb{R}^2$ et $\mathbb{R}^3$ — admis]
\begin{itemize}
\item Dans $\mathbb{R}^2$, les sous-espaces vectoriels sont : $\{(0,0)\}$, les \cle{droites vectorielles} d'équation $ax + by = 0$ (avec $(a,b) \neq (0,0)$), et $\mathbb{R}^2$ tout entier.
\item Dans $\mathbb{R}^3$, les sous-espaces vectoriels sont : $\{(0,0,0)\}$, les droites vectorielles, les \cle{plans vectoriels} d'équation $ax + by + cz = 0$ (avec $(a,b,c) \neq (0,0,0)$), et $\mathbb{R}^3$ tout entier.
\end{itemize}
\end{propriete}

\begin{exemplebox}[Un plan vectoriel de $\mathbb{R}^3$]
Montrons que $P = \{(x,y,z) \in \mathbb{R}^3 \mid x + y - z = 0\}$ est un sous-espace vectoriel de $\mathbb{R}^3$.
\begin{enumerate}
\item $(0,0,0) \in P$ car $0 + 0 - 0 = 0$, donc $P \neq \varnothing$.
\item Soient $u = (x_1, y_1, z_1) \in P$, $v = (x_2, y_2, z_2) \in P$ et $\lambda \in \mathbb{R}$. On a $x_1 + y_1 - z_1 = 0$ et $x_2 + y_2 - z_2 = 0$.
\item $\lambda u + v = (\lambda x_1 + x_2,\ \lambda y_1 + y_2,\ \lambda z_1 + z_2)$ et :
\[
(\lambda x_1 + x_2) + (\lambda y_1 + y_2) - (\lambda z_1 + z_2) = \lambda(x_1 + y_1 - z_1) + (x_2 + y_2 - z_2) = 0.
\]
Donc $\lambda u + v \in P$.
\end{enumerate}
$P$ est un sous-espace vectoriel de $\mathbb{R}^3$ : c'est le plan vectoriel passant par l'origine, de vecteur normal $(1, 1, -1)$.
\end{exemplebox}

\begin{popfigure}
\begin{center}
\begin{tikzpicture}[scale=1.1, >=Stealth]
% Axes
\draw[->, popmuted] (-2.2,0) -- (2.6,0) node[below right] {$x$};
\draw[->, popmuted] (0,-2.2) -- (0,2.6) node[above left] {$y$};
\fill (0,0) circle (1.5pt) node[below left] {$O$};
% Droite vectorielle (sous-espace) : y = (2/3)x
\draw[popblue, very thick, domain=-2.2:2.2, samples=2] plot (\x, {2/3*\x});
\node[popblue, rotate=33.7] at (-1.5,-1.0) {$(d_1)$ : sous-espace};
\node[popblue, rotate=33.7] at (1.35,1.45) {$2x - 3y = 0$};
% Droite ne passant pas par O
\draw[poporange, very thick, domain=-1.5:2.4, samples=2] plot (\x, {0.8*\x + 1.4});
\node[poporange, rotate=38] at (-0.9,2.0) {$(d_2)$ : pas un sous-espace};
% Point sur d2 montrant que 0 n'y est pas
\fill[poporange] (0,1.4) circle (1.5pt) node[right] {$(0\,;\,1{,}4) \in (d_2)$};
\end{tikzpicture}
\end{center}
\legende{$(d_1)$ passe par $O$ : c'est un sous-espace vectoriel de $\mathbb{R}^2$ (équation $2x-3y=0$). $(d_2)$ ne passe pas par $O$ : elle ne contient pas le vecteur nul, ce n'est pas un sous-espace vectoriel.}
\end{popfigure}

\begin{popfigure}
\begin{center}
\begin{tikzpicture}[scale=1.6, >=Stealth]
% Axes en perspective cavalière
\draw[->, popmuted] (0,0) -- (2.6,0) node[below right] {$x$};
\draw[->, popmuted] (0,0) -- (0,2.4) node[above left] {$z$};
\draw[->, popmuted] (0,0) -- (-1.7,-1.2) node[below left] {$y$};
\fill (0,0) circle (1.2pt) node[above right] {$O$};
% Plan x + y - z = 0, c'est-à-dire z = x + y (représenté en perspective)
\coordinate (A) at (1.8, 1.0);
\coordinate (B) at (0.6, 1.9);
\coordinate (C) at (-1.4, 0.2);
\coordinate (D) at (-0.2, -0.7);
\fill[popblueL, opacity=0.55] (A) -- (B) -- (C) -- (D) -- cycle;
\draw[popblue, thick] (A) -- (B) -- (C) -- (D) -- cycle;
% Vecteurs du plan
\draw[->, popdark, thick] (0,0) -- (1.0, 0.55) node[midway, below right] {$\vec{u}$};
\draw[->, popdark, thick] (0,0) -- (-0.55, 0.75) node[midway, above left] {$\vec{v}$};
\draw[->, popgreen, thick] (0,0) -- (0.45, 1.3) node[above] {$\vec{u}+\vec{v}$};
\node[popblue] at (1.7, 1.75) {$P : x + y - z = 0$};
\end{tikzpicture}
\end{center}
\legende{Le plan vectoriel $P$ d'équation $x + y - z = 0$ dans $\mathbb{R}^3$ : il passe par $O$, et toute somme de vecteurs de $P$ reste dans $P$.}
\end{popfigure}

\courssub{Contre-exemples : parties qui ne sont pas des sous-espaces vectoriels}

\begin{exemplebox}[Une droite ne passant pas par l'origine]
Soit $D = \{(x,y) \in \mathbb{R}^2 \mid 2x - 3y = 6\}$. Le vecteur nul vérifie $2 \times 0 - 3 \times 0 = 0 \neq 6$, donc $(0,0) \notin D$. Par conséquent, $D$ \textbf{n'est pas} un sous-espace vectoriel de $\mathbb{R}^2$. On dit que $D$ est une droite \emph{affine}, mais pas vectorielle.
\end{exemplebox}

\begin{exemplebox}[Un demi-plan]
Soit $H = \{(x,y) \in \mathbb{R}^2 \mid x \geq 0\}$. On a $(0,0) \in H$ et $H$ est stable par addition. Mais $H$ n'est pas stable par multiplication par un réel : $u = (1, 0) \in H$ et pourtant $(-1) \cdot u = (-1, 0) \notin H$. Donc $H$ \textbf{n'est pas} un sous-espace vectoriel de $\mathbb{R}^2$. Moralité : il ne suffit pas que $0_E \in F$ et que $F$ soit stable par $+$ ; \textbf{les deux stabilités sont indispensables}.
\end{exemplebox}

\begin{attention}[Conditions nécessaires mais non suffisantes]
Contenir $0_E$ est une condition \emph{nécessaire} (si $0_E \notin F$, c'est fichu), mais pas \emph{suffisante} : l'ensemble $\{(x,y) \in \mathbb{R}^2 \mid xy = 0\}$ (la réunion des deux axes) contient $(0,0)$, mais $(1,0) + (0,1) = (1,1) \notin F$. Toujours vérifier les stabilités !
\end{attention}

\courssub{Intersection de deux sous-espaces vectoriels}

\begin{propriete}[Intersection de deux sous-espaces vectoriels]
Soient $F$ et $G$ deux sous-espaces vectoriels d'un même espace vectoriel $E$. Alors $F \cap G$ est un sous-espace vectoriel de $E$.
\end{propriete}

\begin{experience}[Démonstration guidée]
Soient $F$ et $G$ deux sous-espaces vectoriels de $E$. Montrons que $F \cap G$ vérifie la caractérisation.
\begin{enumerate}
\item \textbf{Non vide} : $F$ et $G$ sont des sous-espaces vectoriels, donc $0_E \in F$ et $0_E \in G$. Ainsi $0_E \in F \cap G$, donc $F \cap G \neq \varnothing$.
\item \textbf{Éléments arbitraires} : soient $x, y \in F \cap G$ et $\lambda \in \mathbb{R}$. Cela signifie que $x \in F$ \emph{et} $x \in G$, de même pour $y$.
\item \textbf{Combinaison} : comme $F$ est un sous-espace vectoriel, $\lambda x + y \in F$. Comme $G$ est un sous-espace vectoriel, $\lambda x + y \in G$. Donc $\lambda x + y \in F \cap G$.
\end{enumerate}
Conclusion : $F \cap G$ est un sous-espace vectoriel de $E$. \hfill$\blacksquare$
\end{experience}

\begin{exemplebox}[Intersection de deux plans vectoriels de $\mathbb{R}^3$]
Dans $\mathbb{R}^3$, soient $P_1 : x + y - z = 0$ et $P_2 : x - y = 0$. Leur intersection est l'ensemble des $(x,y,z)$ vérifiant le système $\begin{cases} x + y - z = 0 \\ x - y = 0 \end{cases}$, c'est-à-dire $y = x$ et $z = 2x$ : c'est la droite vectorielle engendrée par $(1, 1, 2)$. Conformément à la propriété, c'est bien un sous-espace vectoriel de $\mathbb{R}^3$.
\end{exemplebox}

\begin{attention}[La réunion ne marche pas !]
Contrairement à l'intersection, la \textbf{réunion} de deux sous-espaces vectoriels n'est en général \textbf{pas} un sous-espace vectoriel. Dans $\mathbb{R}^2$, prends l'axe des abscisses $F = \{(x,0)\}$ et l'axe des ordonnées $G = \{(0,y)\}$ : $(1,0) \in F \cup G$ et $(0,1) \in F \cup G$, mais $(1,0) + (0,1) = (1,1) \notin F \cup G$. Ne confonds jamais $\cap$ et $\cup$ dans cette propriété.
\end{attention}

\begin{exercice}[À toi de jouer]
\begin{enumerate}
\item Montrer que $F = \{(x,y) \in \mathbb{R}^2 \mid x + 4y = 0\}$ est un sous-espace vectoriel de $\mathbb{R}^2$.
\item L'ensemble $G = \{(x,y,z) \in \mathbb{R}^3 \mid x + y + z = 1\}$ est-il un sous-espace vectoriel de $\mathbb{R}^3$ ? Justifier.
\item Soit $H = \{(x,y) \in \mathbb{R}^2 \mid y = x^2\}$. Montrer que $H$ n'est pas un sous-espace vectoriel de $\mathbb{R}^2$.
\end{enumerate}
\end{exercice}

\begin{corrige}[Exercice]
\begin{enumerate}
\item $(0,0) \in F$ car $0 + 4 \times 0 = 0$. Soient $(x_1,y_1), (x_2,y_2) \in F$ et $\lambda \in \mathbb{R}$ : $(\lambda x_1 + x_2) + 4(\lambda y_1 + y_2) = \lambda(x_1 + 4y_1) + (x_2 + 4y_2) = 0$, donc $\lambda u + v \in F$. $F$ est un sous-espace vectoriel de $\mathbb{R}^2$.
\item $(0,0,0) \notin G$ car $0 + 0 + 0 = 0 \neq 1$. Donc $G$ n'est pas un sous-espace vectoriel de $\mathbb{R}^3$.
\item $(1,1) \in H$ car $1 = 1^2$, mais $2 \cdot (1,1) = (2,2) \notin H$ car $2 \neq 2^2 = 4$. $H$ n'est pas stable par multiplication par un réel : ce n'est pas un sous-espace vectoriel.
\end{enumerate}
\end{corrige}

\begin{aretenir}[L'essentiel de la section]
\begin{itemize}
\item $F$ est un sous-espace vectoriel de $E$ si et seulement si $F \neq \varnothing$ (vérifier $0_E \in F$) et $\lambda x + y \in F$ pour tous $x, y \in F$, $\lambda \in \mathbb{R}$.
\item Si $0_E \notin F$, alors $F$ n'est pas un sous-espace vectoriel : c'est le test le plus rapide.
\item Dans $\mathbb{R}^2$ : droites vectorielles $ax + by = 0$. Dans $\mathbb{R}^3$ : plans $ax + by + cz = 0$ et droites vectorielles.
\item L'intersection de deux sous-espaces vectoriels est un sous-espace vectoriel ; la réunion, en général, non.
\end{itemize}
\end{aretenir}

\coursec{Familles génératrices, familles libres, familles liées}

Dans la section précédente, nous avons appris à reconnaître les \cle{sous-espaces vectoriels} : des parties d'un espace vectoriel qui, elles-mêmes, se comportent comme des espaces vectoriels. Une question naturelle se pose alors : comment \emph{fabriquer} un sous-espace vectoriel à partir de quelques vecteurs donnés ? Et comment décrire un espace vectoriel entier avec le \emph{minimum} de vecteurs possible ?

Imagine un atelier de menuiserie à Douala. Avec deux planches de directions différentes, l'artisan peut reconstituer n'importe quel pan du plan de travail en découpant et assemblant des morceaux proportionnels à ces deux planches. Mais si les deux planches ont exactement la même direction, il ne pourra jamais sortir de cette direction-là. C'est exactement l'idée de cette section : à partir d'une famille de vecteurs, quels vecteurs peut-on \emph{construire} (famille génératrice), et les vecteurs de la famille sont-ils vraiment \emph{indispensables} les uns aux autres (famille libre) ou redondants (famille liée) ?

\courssub{Combinaisons linéaires}

\begin{definition}[Combinaison linéaire]
Soit $E$ un espace vectoriel sur $\mathbb{R}$ et $(x_1, x_2, \ldots, x_n)$ une famille finie de vecteurs de $E$. On appelle \cle{combinaison linéaire} des vecteurs $x_1, \ldots, x_n$ tout vecteur de $E$ de la forme
\[ \lambda_1 x_1 + \lambda_2 x_2 + \cdots + \lambda_n x_n, \]
où $\lambda_1, \lambda_2, \ldots, \lambda_n$ sont des réels, appelés les \cle{coefficients} de la combinaison linéaire.
\end{definition}

Une combinaison linéaire, c'est donc une « recette » : on prend un peu de $x_1$, un peu de $x_2$, etc., et on additionne. Par exemple, dans $\mathbb{R}^2$, le vecteur $(7, 4)$ est une combinaison linéaire de $x_1 = (1, 2)$ et $x_2 = (3, 1)$, car $(7, 4) = 1 \cdot x_1 + 2 \cdot x_2$. Vérifie : $(1,2) + 2(3,1) = (1+6, \, 2+2) = (7,4)$.

\begin{popfigure}
\begin{center}
\begin{tikzpicture}[scale=1.1, >=Stealth]
\draw[->, popmuted] (-0.5,0) -- (4.8,0) node[below left] {$x$};
\draw[->, popmuted] (0,-0.5) -- (0,3.2) node[below left] {$y$};
\draw[->, very thick, popblue] (0,0) -- (1,0.6) node[midway, below left] {$\vec{u}$};
\draw[->, very thick, poporange] (0,0) -- (0.6,1.2) node[midway, right] {$\vec{v}$};
\draw[->, thick, popblue!70, dashed] (0,0) -- (2,1.2) node[midway, below] {$2\vec{u}$};
\draw[->, thick, poporange!80, dashed] (2,1.2) -- (2.6,2.4) node[midway, right] {$\vec{v}$};
\draw[->, very thick, poppurple] (0,0) -- (2.6,2.4) node[midway, above left] {$\vec{w} = 2\vec{u} + \vec{v}$};
\fill[popdark] (2.6,2.4) circle (1.5pt);
\end{tikzpicture}
\end{center}
\legende{Le vecteur $\vec{w}$ est combinaison linéaire de $\vec{u}$ et $\vec{v}$ : on parcourt $2\vec{u}$ puis $\vec{v}$ (règle du parallélogramme).}
\end{popfigure}

\courssub{Familles génératrices et sous-espace engendré}

\begin{definition}[Famille génératrice, sous-espace engendré]
Soit $E$ un espace vectoriel sur $\mathbb{R}$ et $(x_1, \ldots, x_n)$ une famille de vecteurs de $E$.
\begin{itemize}
\item L'ensemble de \textbf{toutes} les combinaisons linéaires de $x_1, \ldots, x_n$ est un sous-espace vectoriel de $E$, appelé \cle{sous-espace engendré} par la famille, et noté
\[ \mathrm{Vect}(x_1, \ldots, x_n) = \left\{ \lambda_1 x_1 + \cdots + \lambda_n x_n \; ; \; \lambda_1, \ldots, \lambda_n \in \mathbb{R} \right\}. \]
\item On dit que la famille $(x_1, \ldots, x_n)$ est \cle{génératrice} de $E$ (ou qu'elle \emph{engendre} $E$) lorsque tout vecteur de $E$ est combinaison linéaire de $x_1, \ldots, x_n$, c'est-à-dire lorsque $\mathrm{Vect}(x_1, \ldots, x_n) = E$.
\end{itemize}
\end{definition}

\begin{exemplebox}[Sous-espace engendré par un ou deux vecteurs]
\begin{itemize}
\item Dans le plan, si $\vec{u} \neq \vec{0}$, alors $\mathrm{Vect}(\vec{u})$ est l'ensemble des vecteurs colinéaires à $\vec{u}$ : c'est la \textbf{droite vectorielle} dirigée par $\vec{u}$.
\item Dans l'espace, si $\vec{u}$ et $\vec{v}$ ne sont pas colinéaires, $\mathrm{Vect}(\vec{u}, \vec{v})$ est le \textbf{plan vectoriel} engendré par $\vec{u}$ et $\vec{v}$.
\item Dans $\mathbb{R}^2$, la famille $\big((1,0), (0,1)\big)$ est génératrice : tout vecteur $(x, y)$ s'écrit $x(1,0) + y(0,1)$.
\end{itemize}
\end{exemplebox}

\begin{methode}[Montrer qu'une famille est génératrice]
Pour montrer que $(x_1, \ldots, x_n)$ engendre $E$ :
\begin{enumerate}
\item On prend un vecteur \textbf{quelconque} $v$ de $E$ (avec des coordonnées littérales, par exemple $(x, y)$ dans $\mathbb{R}^2$).
\item On cherche des réels $\lambda_1, \ldots, \lambda_n$ tels que $v = \lambda_1 x_1 + \cdots + \lambda_n x_n$ : cela revient à résoudre un \textbf{système linéaire} d'inconnues $\lambda_1, \ldots, \lambda_n$.
\item Si le système admet au moins une solution \emph{pour tout} $v$, la famille est génératrice. Sinon, elle ne l'est pas.
\end{enumerate}
\end{methode}

\begin{exoresolu}[La famille $\big((1,2), (3,1)\big)$ engendre $\mathbb{R}^2$]
Montrer que la famille $\big(u, v\big)$ avec $u = (1, 2)$ et $v = (3, 1)$ est génératrice de $\mathbb{R}^2$.
\tcblower
Soit $(x, y)$ un vecteur quelconque de $\mathbb{R}^2$. On cherche $\lambda, \mu \in \mathbb{R}$ tels que $(x, y) = \lambda(1,2) + \mu(3,1)$, c'est-à-dire :
\[ \begin{cases} \lambda + 3\mu = x \\ 2\lambda + \mu = y \end{cases} \]
De la première équation, $\lambda = x - 3\mu$. En reportant dans la deuxième : $2(x - 3\mu) + \mu = y$, soit $-5\mu = y - 2x$, d'où $\mu = \dfrac{2x - y}{5}$ puis $\lambda = x - 3 \cdot \dfrac{2x - y}{5} = \dfrac{3y - x}{5}$.

Le système admet une solution pour \textbf{tout} $(x, y)$. Par exemple, pour $(x, y) = (7, 4)$ : $\mu = \frac{14 - 4}{5} = 2$ et $\lambda = \frac{12 - 7}{5} = 1$, ce qui redonne $(7,4) = u + 2v$. La famille est donc \textbf{génératrice} de $\mathbb{R}^2$.
\end{exoresolu}

\courssub{Familles libres et familles liées}

Voici maintenant la question de la \emph{redondance}. Si $\vec{w} = 2\vec{u} + \vec{v}$, alors le vecteur $\vec{w}$ n'apporte rien de nouveau à la famille $(\vec{u}, \vec{v}, \vec{w})$ : tout ce qu'on fabrique avec les trois vecteurs pouvait déjà l'être avec $\vec{u}$ et $\vec{v}$ seuls. La famille contient un « vecteur de trop ». La notion de famille libre capture cette idée d'indépendance.

\begin{definition}[Famille libre, famille liée]
Soit $E$ un espace vectoriel sur $\mathbb{R}$ et $(x_1, \ldots, x_n)$ une famille de vecteurs de $E$.
\begin{itemize}
\item La famille est dite \cle{libre} (ou ses vecteurs sont dits \emph{linéairement indépendants}) lorsque la seule combinaison linéaire nulle est celle dont tous les coefficients sont nuls :
\[ \lambda_1 x_1 + \lambda_2 x_2 + \cdots + \lambda_n x_n = 0_E \;\Longrightarrow\; \lambda_1 = \lambda_2 = \cdots = \lambda_n = 0. \]
\item La famille est dite \cle{liée} lorsqu'elle n'est pas libre, c'est-à-dire lorsqu'il existe des réels $\lambda_1, \ldots, \lambda_n$ \textbf{non tous nuls} tels que $\lambda_1 x_1 + \cdots + \lambda_n x_n = 0_E$.
\end{itemize}
\end{definition}

\begin{attention}[Le piège classique]
Une erreur très fréquente consiste à dire : « la famille est libre car aucun de ses vecteurs n'est le vecteur nul ». C'est \textbf{faux} ! Par exemple, la famille $\big((1,1), (2,2)\big)$ ne contient pas le vecteur nul, et pourtant elle est liée, car $2(1,1) - (2,2) = (0,0)$ avec des coefficients non nuls. Ce qui compte, c'est que la \textbf{combinaison linéaire nulle} force \textbf{tous les coefficients} à être nuls. Autre remarque : toute famille contenant le vecteur nul est automatiquement liée (prends le coefficient $1$ devant $0_E$ et $0$ ailleurs).
\end{attention}

\begin{methode}[Tester si une famille est libre ou liée]
\begin{enumerate}
\item On écrit l'équation $\lambda_1 x_1 + \cdots + \lambda_n x_n = 0_E$ avec des inconnues $\lambda_1, \ldots, \lambda_n$.
\item On traduit cette équation vectorielle en un \textbf{système homogène} (une équation par coordonnée).
\item On résout le système :
\begin{itemize}
\item si l'\textbf{unique} solution est $\lambda_1 = \cdots = \lambda_n = 0$, la famille est \textbf{libre} ;
\item s'il existe une solution \textbf{non nulle}, la famille est \textbf{liée}, et cette solution fournit une \emph{relation de liaison} entre les vecteurs.
\end{itemize}
\end{enumerate}
\end{methode}

\begin{exemplebox}[Deux familles dans $\mathbb{R}^2$]
\textbf{Famille libre.} Soient $u = (1, 2)$ et $v = (3, 1)$. L'équation $\lambda u + \mu v = (0,0)$ donne le système
\[ \begin{cases} \lambda + 3\mu = 0 \\ 2\lambda + \mu = 0 \end{cases} \]
De la première équation, $\lambda = -3\mu$ ; en reportant : $-6\mu + \mu = 0$, soit $-5\mu = 0$, donc $\mu = 0$ puis $\lambda = 0$. L'unique solution est nulle : la famille $\big((1,2), (3,1)\big)$ est \textbf{libre}.

\textbf{Famille liée.} Soient $a = (1, 1)$ et $b = (2, 2)$. L'équation $\lambda a + \mu b = (0,0)$ donne $\lambda + 2\mu = 0$ (deux fois la même équation). Le couple $(\lambda, \mu) = (2, -1)$ est une solution non nulle : la famille est \textbf{liée}, et $2a - b = 0$, c'est-à-dire $b = 2a$.
\end{exemplebox}

\begin{savaistu}[Liberté et colinéarité]
Pour une famille de \textbf{deux} vecteurs, être liée signifie exactement être \textbf{colinéaire} : si $\lambda u + \mu v = 0_E$ avec, disons, $\lambda \neq 0$, alors $u = -\dfrac{\mu}{\lambda} v$. Ainsi, la notion de famille libre \emph{généralise} à un nombre quelconque de vecteurs la notion de non-colinéarité que tu connais depuis la Seconde. De même, trois vecteurs de l'espace forment une famille liée si et seulement s'ils sont \textbf{coplanaires}. L'algèbre linéaire unifie toutes ces idées géométriques en un seul test : résoudre un système homogène !
\end{savaistu}

\begin{popfigure}
\begin{center}
\begin{tikzpicture}[scale=1.0, >=Stealth]
% Plan oblique simulé : base (2,0) et (0.5,1.5)
\coordinate (O) at (0,0);
\coordinate (U) at (2,0);
\coordinate (V) at (0.5,1.5);
\coordinate (W) at ($(U) + 0.5*(V)$); % W = U + 0.5 V = (2.25, 0.75)
\fill[popblueL, opacity=0.3] (O) -- (U) -- (W) -- (V) -- cycle;
\draw[popmuted] (O) -- (U) -- (W) -- (V) -- cycle;
\fill[popdark] (O) circle (1.5pt) node[below left, popdark] {$O$};
\draw[->, very thick, popblue] (O) -- (U) node[right] {$\vec{u}$};
\draw[->, very thick, poporange] (O) -- (V) node[above left] {$\vec{v}$};
\draw[->, very thick, poppurple] (O) -- (W) node[above right] {$\vec{w}$};
\draw[->, thick, popblue!70, dashed] (O) -- (U);
\draw[->, thick, poporange!80, dashed] (U) -- (W);
\node[popdark, align=center, font=\footnotesize] at (-1.5,0.5) {$\vec{w} = \vec{u} + \frac{1}{2}\vec{v}$\\ famille liée};
\end{tikzpicture}
\end{center}
\legende{Trois vecteurs coplanaires dans l'espace : $\vec{w}$ est combinaison linéaire de $\vec{u}$ et $\vec{v}$, donc la famille $(\vec{u}, \vec{v}, \vec{w})$ est liée.}
\end{popfigure}

\courssub{Exemples dans $\mathbb{R}^3$ et propriété fondamentale}

\begin{exoresolu}[Étude d'une famille de trois vecteurs de $\mathbb{R}^3$]
Dans $\mathbb{R}^3$, on considère $x_1 = (1, 0, 1)$, $x_2 = (1, 1, 0)$ et $x_3 = (0, 1, -1)$. La famille $(x_1, x_2, x_3)$ est-elle libre ou liée ?
\tcblower
Soient $\lambda_1, \lambda_2, \lambda_3 \in \mathbb{R}$ tels que $\lambda_1 x_1 + \lambda_2 x_2 + \lambda_3 x_3 = (0, 0, 0)$. On obtient le système homogène :
\[ \begin{cases} \lambda_1 + \lambda_2 = 0 \\ \lambda_2 + \lambda_3 = 0 \\ \lambda_1 - \lambda_3 = 0 \end{cases} \]
De la première équation : $\lambda_1 = -\lambda_2$. De la deuxième : $\lambda_3 = -\lambda_2$. La troisième donne alors $\lambda_1 - \lambda_3 = -\lambda_2 + \lambda_2 = 0$ : elle est automatiquement vérifiée. En choisissant $\lambda_2 = -1$, on obtient $(\lambda_1, \lambda_2, \lambda_3) = (1, -1, 1)$, solution \textbf{non nulle}.

La famille est donc \textbf{liée}, et on a la relation de liaison $x_1 - x_2 + x_3 = 0$, c'est-à-dire $x_3 = x_2 - x_1$. Vérifie : $(1,1,0) - (1,0,1) = (0, 1, -1) = x_3$. 
\end{exoresolu}

\begin{exemplebox}[La famille canonique de $\mathbb{R}^3$]
Soient $e_1 = (1, 0, 0)$, $e_2 = (0, 1, 0)$, $e_3 = (0, 0, 1)$. L'équation $\lambda_1 e_1 + \lambda_2 e_2 + \lambda_3 e_3 = (0,0,0)$ s'écrit directement $(\lambda_1, \lambda_2, \lambda_3) = (0, 0, 0)$ : la famille est \textbf{libre}. De plus, tout vecteur $(x, y, z) = x e_1 + y e_2 + z e_3$ : elle est aussi \textbf{génératrice}. C'est la \cle{famille canonique} de $\mathbb{R}^3$ — nous verrons dans la prochaine section qu'une famille à la fois libre et génératrice s'appelle une \emph{base}.
\end{exemplebox}

\begin{propriete}[Familles « trop grandes » dans $\mathbb{R}^n$ -- Démonstration exigible au Probatoire]
Toute famille de \textbf{trois vecteurs ou plus} de $\mathbb{R}^2$ est liée. Plus généralement, toute famille de $n+1$ vecteurs ou plus de $\mathbb{R}^n$ est liée.
\end{propriete}

\begin{experience}[Pourquoi trois vecteurs du plan sont toujours liés]
Justifions ce résultat dans $\mathbb{R}^2$ (la démonstration générale viendra avec la notion de dimension). Soient $u = (a, b)$, $v = (c, d)$, $w = (e, f)$ trois vecteurs de $\mathbb{R}^2$. L'équation $\lambda_1 u + \lambda_2 v + \lambda_3 w = (0,0)$ donne un système de \textbf{2 équations} à \textbf{3 inconnues} :
\[ \begin{cases} a\lambda_1 + c\lambda_2 + e\lambda_3 = 0 \\ b\lambda_1 + d\lambda_2 + f\lambda_3 = 0 \end{cases} \]
Un système homogène qui possède \emph{plus d'inconnues que d'équations} admet toujours une infinité de solutions, donc en particulier une solution non nulle. Par exemple, si $u$ et $v$ ne sont pas colinéaires, on sait (famille génératrice !) que $w$ s'écrit $w = \alpha u + \beta v$, d'où la relation $\alpha u + \beta v - w = 0$. Et si $u$ et $v$ sont déjà colinéaires, la sous-famille $(u, v)$ est déjà liée. Dans tous les cas, la famille $(u, v, w)$ est liée. Géométriquement : trois vecteurs du plan sont toujours « coplanaires », puisqu'ils vivent déjà dans un plan !
\end{experience}

\begin{aretenir}[L'essentiel de la section]
\begin{itemize}
\item Une \cle{combinaison linéaire} de $x_1, \ldots, x_n$ est un vecteur de la forme $\lambda_1 x_1 + \cdots + \lambda_n x_n$.
\item $\mathrm{Vect}(x_1, \ldots, x_n)$ est le sous-espace de toutes ces combinaisons ; la famille est \cle{génératrice} de $E$ si $\mathrm{Vect}(x_1, \ldots, x_n) = E$ — on le prouve en résolvant un système pour un vecteur \emph{quelconque}.
\item La famille est \cle{libre} si $\lambda_1 x_1 + \cdots + \lambda_n x_n = 0_E$ impose $\lambda_1 = \cdots = \lambda_n = 0$ ; sinon elle est \cle{liée}, et une solution non nulle du système homogène donne une relation de liaison.
\item Deux vecteurs liés $=$ deux vecteurs colinéaires ; trois vecteurs liés dans l'espace $=$ trois vecteurs coplanaires.
\item Dans $\mathbb{R}^2$, toute famille de $3$ vecteurs ou plus est liée ; dans $\mathbb{R}^n$, toute famille de $n+1$ vecteurs ou plus est liée.
\end{itemize}
\end{aretenir}

\begin{exercice}[À toi de jouer]
\begin{enumerate}
\item Montrer que la famille $\big((1, -1), (2, 3)\big)$ est libre et génératrice de $\mathbb{R}^2$.
\item Dans $\mathbb{R}^3$, montrer que la famille $\big((1, 2, 0), (0, 1, 1), (1, 3, 1)\big)$ est liée et donner une relation de liaison.
\item Le vecteur $(5, 1)$ appartient-il à $\mathrm{Vect}\big((1, 2), (2, 4)\big)$ ? Justifier.
\end{enumerate}
\end{exercice}

\begin{corrige}[Exercices proposés]
\textbf{1.} Liberté : $\lambda(1,-1) + \mu(2,3) = (0,0)$ donne $\lambda + 2\mu = 0$ et $-\lambda + 3\mu = 0$. En additionnant : $5\mu = 0$, donc $\mu = 0$ puis $\lambda = 0$ : famille libre. Génération : pour $(x, y)$ quelconque, le système $\lambda + 2\mu = x$, $-\lambda + 3\mu = y$ donne $\mu = \frac{x+y}{5}$ et $\lambda = \frac{3x - 2y}{5}$ : solution pour tout $(x,y)$, donc famille génératrice.

\textbf{2.} On remarque que $(1,2,0) + (0,1,1) = (1,3,1)$ : le troisième vecteur est la somme des deux premiers. La relation $x_1 + x_2 - x_3 = 0$ est une combinaison linéaire nulle à coefficients non tous nuls : la famille est liée.

\textbf{3.} Les vecteurs $(1,2)$ et $(2,4)$ sont colinéaires ($(2,4) = 2(1,2)$), donc $\mathrm{Vect}\big((1,2),(2,4)\big) = \mathrm{Vect}\big((1,2)\big)$ est la droite vectorielle des vecteurs de la forme $(\lambda, 2\lambda)$. Or $(5, 1)$ n'est pas de cette forme (car $1 \neq 2 \times 5$) : il n'appartient pas à ce sous-espace.
\end{corrige}

Tu sais désormais reconnaître les familles qui « suffisent » à décrire un espace (génératrices) et celles qui sont « sans gaspillage » (libres). Dans la prochaine section, nous réunirons ces deux qualités en une seule notion capitale : les \emph{bases}, qui permettront de définir la \emph{dimension} d'un espace vectoriel.

\coursec{Bases et dimension d'un espace vectoriel}

Dans la section précédente, nous avons rencontré deux types de familles de vecteurs : les \cle{familles génératrices}, qui permettent d'atteindre tous les vecteurs de l'espace par combinaisons linéaires, et les \cle{familles libres}, dans lesquelles aucun vecteur n'est « redondant ». Une famille génératrice peut contenir des vecteurs « en trop » ; une famille libre peut être « trop pauvre » pour engendrer tout l'espace. L'idée géniale de cette section est de réunir les deux qualités : une famille qui engendre tout l'espace \emph{sans gaspillage}. C'est exactement ce qu'on appelle une \cle{base}. Les bases sont les « systèmes de coordonnées » des espaces vectoriels : grâce à elles, chaque vecteur possède une adresse unique, comme chaque quartier de Yaoundé possède une adresse unique dans la ville.

\courssub{Notion de base}

\begin{definition}[Base d'un espace vectoriel]
Soit $E$ un espace vectoriel sur $\mathbb{R}$. Une famille $(\vec{e}_1, \vec{e}_2, \ldots, \vec{e}_n)$ de vecteurs de $E$ est une \cle{base} de $E$ si elle est à la fois :
\begin{itemize}
\item \textbf{libre} : la seule combinaison linéaire nulle est celle dont tous les coefficients sont nuls ;
\item \textbf{génératrice} : tout vecteur de $E$ est combinaison linéaire de $\vec{e}_1, \vec{e}_2, \ldots, \vec{e}_n$.
\end{itemize}
Autrement dit : une base est une famille génératrice \emph{minimale} (on ne peut rien lui retirer) et une famille libre \emph{maximale} (on ne peut rien lui ajouter).
\end{definition}

\begin{exemplebox}[Premiers exemples]
\begin{itemize}
\item La famille $(1)$ est une base de $\mathbb{R}$ : tout réel $x$ s'écrit $x = x \times 1$ (génératrice) et $\lambda \times 1 = 0$ impose $\lambda = 0$ (libre).
\item Dans $\mathbb{R}^2$, la famille $\big((1,0), (0,1)\big)$ est une base : tout couple $(x,y)$ s'écrit $(x,y) = x(1,0) + y(0,1)$, et $\alpha(1,0) + \beta(0,1) = (0,0)$ donne $(\alpha, \beta) = (0,0)$.
\item Dans le plan vectoriel, deux vecteurs $\vec{\imath}$ et $\vec{\jmath}$ non colinéaires forment une base : c'est le fondement de la notion de repère vue en Seconde.
\end{itemize}
\end{exemplebox}

\courssub{Décomposition unique d'un vecteur dans une base}

Voici la propriété qui rend les bases indispensables : dans une base, chaque vecteur possède une \emph{écriture unique}. C'est l'unicité qui distingue une base d'une simple famille génératrice.

\begin{propriete}[Décomposition unique dans une base — démonstration exigible]
Soit $E$ un espace vectoriel sur $\mathbb{R}$ et $\mathcal{B} = (\vec{e}_1, \vec{e}_2, \ldots, \vec{e}_n)$ une base de $E$. Alors tout vecteur $\vec{u}$ de $E$ s'écrit de \textbf{manière unique} :
\[ \vec{u} = x_1 \vec{e}_1 + x_2 \vec{e}_2 + \cdots + x_n \vec{e}_n \]
où $x_1, x_2, \ldots, x_n$ sont des réels.
\end{propriete}

\begin{experience}[Démonstration guidée : existence et unicité]
Nous allons prouver séparément l'\textbf{existence} puis l'\textbf{unicité}. Observe bien la stratégie : chaque hypothèse de la définition de base sert exactement une fois.

\textbf{Étape 1 : existence (on utilise le caractère générateur).}

Soit $\vec{u} \in E$. La famille $(\vec{e}_1, \ldots, \vec{e}_n)$ est génératrice : par définition, tout vecteur de $E$ est combinaison linéaire des $\vec{e}_i$. Il existe donc des réels $x_1, \ldots, x_n$ tels que :
\[ \vec{u} = x_1 \vec{e}_1 + x_2 \vec{e}_2 + \cdots + x_n \vec{e}_n. \]
L'écriture existe. Mais pourrait-il y en avoir une autre ?

\textbf{Étape 2 : unicité (on utilise la liberté).}

Supposons qu'il existe une \emph{autre} écriture :
\[ \vec{u} = y_1 \vec{e}_1 + y_2 \vec{e}_2 + \cdots + y_n \vec{e}_n. \]
En soustrayant membre à membre les deux égalités, on obtient :
\[ \vec{0} = (x_1 - y_1)\vec{e}_1 + (x_2 - y_2)\vec{e}_2 + \cdots + (x_n - y_n)\vec{e}_n. \]
Or la famille $(\vec{e}_1, \ldots, \vec{e}_n)$ est \textbf{libre} : la seule combinaison linéaire donnant $\vec{0}$ a tous ses coefficients nuls. Donc :
\[ x_1 - y_1 = 0, \quad x_2 - y_2 = 0, \quad \ldots, \quad x_n - y_n = 0, \]
c'est-à-dire $x_i = y_i$ pour tout $i$. Les deux écritures sont identiques : la décomposition est \textbf{unique}. $\blacksquare$

\textbf{Conclusion à retenir :} l'existence vient de « génératrice », l'unicité vient de « libre ». C'est pourquoi une base exige les deux propriétés.
\end{experience}

\begin{definition}[Coordonnées d'un vecteur dans une base]
Soit $\mathcal{B} = (\vec{e}_1, \ldots, \vec{e}_n)$ une base de $E$ et $\vec{u} \in E$. Les uniques réels $x_1, x_2, \ldots, x_n$ tels que
\[ \vec{u} = x_1 \vec{e}_1 + x_2 \vec{e}_2 + \cdots + x_n \vec{e}_n \]
sont appelés les \cle{coordonnées} (ou \cle{composantes}) de $\vec{u}$ dans la base $\mathcal{B}$. On écrit $\vec{u}\,(x_1, x_2, \ldots, x_n)$ dans la base $\mathcal{B}$, ou $\vec{u} = \begin{pmatrix} x_1 \\ x_2 \\ \vdots \\ x_n \end{pmatrix}_{\mathcal{B}}$.
\end{definition}

\begin{attention}[L'ordre des vecteurs compte !]
Une base est une famille \emph{ordonnée}. Dans $\mathbb{R}^2$, le vecteur $(3,5)$ a pour coordonnées $(3,5)$ dans la base $\big((1,0),(0,1)\big)$, mais $(5,3)$ dans la base $\big((0,1),(1,0)\big)$. Échanger l'ordre des vecteurs de la base échange les coordonnées. Précise toujours la base dans laquelle tu travailles.
\end{attention}

\begin{popfigure}
\begin{center}
\begin{tikzpicture}[scale=1.1, >=Stealth]
\draw[popline] (-0.5,0) -- (4.5,0) node[below left] {};
\draw[popline] (0,-0.5) -- (0,3.2);
\draw[->, very thick, popblue] (0,0) -- (1,0) node[midway, below] {$\vec{\imath}$};
\draw[->, very thick, popblue] (0,0) -- (0,1) node[midway, left] {$\vec{\jmath}$};
\draw[->, very thick, poporange] (0,0) -- (3,2) node[midway, above left] {$\vec{u}$};
\draw[->, thick, popgreen, dashed] (0,0) -- (3,0) node[midway, below] {$3\vec{\imath}$};
\draw[->, thick, popgreen, dashed] (3,0) -- (3,2) node[midway, right] {$2\vec{\jmath}$};
\draw[popmuted, dotted] (0,2) -- (3,2);
\fill[popdark] (0,0) circle (1.5pt) node[below left] {$O$};
\fill[poporange] (3,2) circle (1.5pt) node[above right] {$M(3\,;\,2)$};
\foreach \x in {1,2,3} \draw[popline] (\x,0.06) -- (\x,-0.06);
\foreach \y in {1,2} \draw[popline] (0.06,\y) -- (-0.06,\y);
\end{tikzpicture}
\end{center}
\legende{Dans le repère $(O\,;\,\vec{\imath},\vec{\jmath})$, le vecteur $\vec{u} = 3\vec{\imath} + 2\vec{\jmath}$ : ses coordonnées $(3\,;\,2)$ sont les coefficients uniques de la décomposition dans la base $(\vec{\imath}, \vec{\jmath})$.}
\end{popfigure}

\courssub{Bases canoniques}

Certains espaces vectoriels usuels possèdent une base « naturelle », dite \cle{canonique}, qu'il faut connaître par cœur.

\begin{aretenir}[Bases canoniques à connaître]
\begin{itemize}
\item \textbf{Dans $\mathbb{R}$} : la base canonique est $(1)$. Tout réel $x$ a pour coordonnée $x$ lui-même.
\item \textbf{Dans $\mathbb{R}^2$} : la base canonique est $(e_1, e_2)$ avec $e_1 = (1,0)$ et $e_2 = (0,1)$. On a $(x,y) = x\,e_1 + y\,e_2$.
\item \textbf{Dans $\mathbb{R}^3$} : la base canonique est $(e_1, e_2, e_3)$ avec $e_1 = (1,0,0)$, $e_2 = (0,1,0)$, $e_3 = (0,0,1)$. On a $(x,y,z) = x\,e_1 + y\,e_2 + z\,e_3$.
\item \textbf{Vecteurs du plan} : toute base $(\vec{\imath}, \vec{\jmath})$ formée de deux vecteurs non colinéaires ; la base canonique correspond au repère orthonormé habituel.
\item \textbf{Vecteurs de l'espace} : toute base $(\vec{\imath}, \vec{\jmath}, \vec{k})$ formée de trois vecteurs non coplanaires.
\end{itemize}
\end{aretenir}

\begin{exemplebox}[Lecture de coordonnées]
Dans $\mathbb{R}^3$ muni de sa base canonique, le vecteur $u = (2, -5, 7)$ se décompose :
\[ u = 2\,e_1 - 5\,e_2 + 7\,e_3 = 2(1,0,0) - 5(0,1,0) + 7(0,0,1). \]
Ses coordonnées dans la base canonique sont $(2, -5, 7)$ : dans une base canonique, les coordonnées « se lisent » directement. Attention : ce n'est plus vrai dans une base quelconque, comme le montre l'exercice résolu plus bas.
\end{exemplebox}

\courssub{Théorème de la dimension}

Une question naturelle se pose : un même espace vectoriel peut-il avoir des bases avec des nombres de vecteurs différents ? Par exemple, $\mathbb{R}^2$ admet la base $\big((1,0),(0,1)\big)$ à 2 vecteurs ; pourrait-il admettre une base à 3 vecteurs ? La réponse est non, et c'est l'un des théorèmes les plus importants de l'année.

\begin{propriete}[Théorème de la dimension — admis]
Si un espace vectoriel $E$ sur $\mathbb{R}$ admet une famille génératrice \textbf{finie}, alors $E$ admet des bases, et \textbf{toutes les bases de $E$ ont le même nombre de vecteurs}. Ce nombre entier est appelé la \cle{dimension} de $E$ et se note $\dim E$.

\emph{Ce théorème est admis au programme de Première : sa démonstration n'est pas exigible, mais son énoncé et son utilisation doivent être parfaitement maîtrisés.}
\end{propriete}

\begin{aretenir}[Dimensions usuelles]
\begin{center}
\begin{tabularx}{\linewidth}{|l|X|c|}
\hline
\rowcolor{popblueL} \textbf{Espace vectoriel} & \textbf{Base canonique} & \textbf{Dimension} \\
\hline
$\mathbb{R}$ & $(1)$ & $1$ \\
\hline
$\mathbb{R}^2$ & $\big((1,0),(0,1)\big)$ & $2$ \\
\hline
$\mathbb{R}^3$ & $\big((1,0,0),(0,1,0),(0,0,1)\big)$ & $3$ \\
\hline
Vecteurs du plan & $(\vec{\imath}, \vec{\jmath})$ & $2$ \\
\hline
Vecteurs de l'espace & $(\vec{\imath}, \vec{\jmath}, \vec{k})$ & $3$ \\
\hline
\end{tabularx}
\end{center}
Par convention, l'espace nul $\{\vec{0}\}$ est de dimension $0$ (sa base est la famille vide).
\end{aretenir}

\begin{propriete}[Caractérisations pratiques en dimension $n$ — admises]
Soit $E$ un espace vectoriel de dimension $n$ et $\mathcal{F}$ une famille de \textbf{exactement $n$ vecteurs} de $E$. Alors :
\begin{itemize}
\item si $\mathcal{F}$ est \textbf{libre}, alors $\mathcal{F}$ est une base de $E$ ;
\item si $\mathcal{F}$ est \textbf{génératrice}, alors $\mathcal{F}$ est une base de $E$.
\end{itemize}
En dimension connue, avec le bon nombre de vecteurs, \emph{une seule} des deux vérifications suffit. Ces résultats sont admis.
\end{propriete}

\begin{attention}[Une famille génératrice n'est pas toujours une base]
C'est l'erreur la plus fréquente au Probatoire. La famille $\big((1,0), (0,1), (1,1)\big)$ engendre $\mathbb{R}^2$ (elle contient la base canonique), mais elle n'est \textbf{pas} une base car elle est liée : $(1,1) = (1,0) + (0,1)$. Elle contient un vecteur « en trop ». De même, une famille libre n'est pas toujours une base : $(1,0)$ seul est libre dans $\mathbb{R}^2$ mais n'engendre pas $\mathbb{R}^2$ (le vecteur $(0,1)$ est hors d'atteinte). \textbf{Base = libre ET génératrice}, sauf si la dimension est connue et que le nombre de vecteurs est exactement $\dim E$.
\end{attention}

\begin{popfigure}
\begin{center}
\begin{tikzpicture}[
  node distance=0.9cm and 1.6cm,
  boite/.style={rectangle, rounded corners=4pt, draw=popblue, thick, fill=popblueL, text width=3.6cm, align=center, minimum height=1.1cm, font=\small},
  res/.style={rectangle, rounded corners=4pt, draw=poporange, very thick, fill=poporangeL, text width=4.2cm, align=center, minimum height=1.1cm, font=\small\bfseries},
  dim/.style={rectangle, rounded corners=4pt, draw=popgreen, very thick, fill=popgreenL, text width=5.4cm, align=center, minimum height=1cm, font=\small}
]
\node[boite, align=center] (libre) {Famille \textbf{libre}\\ (aucun vecteur redondant)};
\node[boite, below=of libre, align=center] (gen) {Famille \textbf{génératrice}\\ (elle atteint tout l'espace)};
\coordinate (mid) at ($(libre)!0.5!(gen)$);
\node[res, right=3cm of mid, align=center] (base) {BASE\\ libre $+$ génératrice};
\node[dim, right=4cm of base, align=center] (dimension) {$\dim E = $ nombre de vecteurs\\ de toute base de $E$};
\draw[->, very thick, popblue] (libre.east) -- (base.north west);
\draw[->, very thick, popblue] (gen.east) -- (base.south west);
\draw[->, very thick, poporange] (base.east) -- (dimension.west);
\end{tikzpicture}
\end{center}
\legende{Schéma de synthèse : une base est une famille à la fois libre et génératrice ; le nombre commun de vecteurs de toutes les bases est la dimension de l'espace.}
\end{popfigure}

\courssub{Méthodes et exemples résolus}

\begin{methode}[Montrer qu'une famille est une base]
Pour montrer qu'une famille $\mathcal{F} = (\vec{e}_1, \ldots, \vec{e}_p)$ est une base de $E$ :
\begin{enumerate}
\item \textbf{Si la dimension de $E$ est inconnue} : montrer que $\mathcal{F}$ est libre (l'équation $\lambda_1 \vec{e}_1 + \cdots + \lambda_p \vec{e}_p = \vec{0}$ n'a que la solution nulle) \emph{puis} que $\mathcal{F}$ est génératrice (tout vecteur de $E$ se décompose sur les $\vec{e}_i$).
\item \textbf{Si $\dim E = n$ est connue et $p = n$} : une seule vérification suffit — montrer que $\mathcal{F}$ est libre (le plus souvent) \emph{ou} qu'elle est génératrice.
\item \textbf{Contrôle rapide} : si $p > \dim E$, la famille est forcément liée ; si $p < \dim E$, elle ne peut pas être génératrice. Dans les deux cas, ce n'est pas une base.
\end{enumerate}
\end{methode}

\begin{exoresolu}[Une base de $\mathbb{R}^2$ et un changement de coordonnées]
On considère dans $\mathbb{R}^2$ les vecteurs $u = (1,1)$ et $v = (1,-1)$.
\begin{enumerate}
\item Montrer que $(u, v)$ est une base de $\mathbb{R}^2$.
\item Déterminer les coordonnées du vecteur $w = (3,5)$ dans la base $(u,v)$.
\end{enumerate}
\tcblower
\textbf{1.} On sait que $\dim \mathbb{R}^2 = 2$ et la famille $(u,v)$ compte exactement $2$ vecteurs : il suffit de montrer qu'elle est libre.

Soient $\alpha, \beta \in \mathbb{R}$ tels que $\alpha u + \beta v = (0,0)$. Alors :
\[ \alpha(1,1) + \beta(1,-1) = (\alpha + \beta,\ \alpha - \beta) = (0,0) \]
d'où le système :
\[ \begin{cases} \alpha + \beta = 0 \\ \alpha - \beta = 0 \end{cases} \]
En additionnant les deux équations : $2\alpha = 0$, donc $\alpha = 0$, puis $\beta = 0$. La famille $(u,v)$ est libre. Comme elle possède $2 = \dim \mathbb{R}^2$ vecteurs, c'est une \textbf{base} de $\mathbb{R}^2$.

\textbf{2.} On cherche les réels $a$ et $b$ tels que $w = a\,u + b\,v$, c'est-à-dire :
\[ (3,5) = a(1,1) + b(1,-1) = (a+b,\ a-b). \]
On résout :
\[ \begin{cases} a + b = 3 \\ a - b = 5 \end{cases} \]
Par addition : $2a = 8$, donc $a = 4$ ; par soustraction : $2b = -2$, donc $b = -1$.

\textbf{Vérification :} $4(1,1) - 1(1,-1) = (4,4) - (1,-1) = (3,5)$. \checkmark

Les coordonnées de $w$ dans la base $(u,v)$ sont donc $(4, -1)$, alors qu'elles sont $(3,5)$ dans la base canonique : \emph{les coordonnées dépendent de la base choisie}.
\end{exoresolu}

\begin{methode}[Déterminer la dimension d'un sous-espace vectoriel]
Pour déterminer la dimension d'un sous-espace vectoriel $F$ de $E$ :
\begin{enumerate}
\item Écrire un vecteur générique de $F$ à l'aide des conditions définissant $F$ (souvent une ou plusieurs équations).
\item Décomposer ce vecteur générique comme combinaison linéaire de vecteurs fixes : on obtient une famille \textbf{génératrice} de $F$.
\item Vérifier que cette famille est \textbf{libre} : c'est alors une base de $F$.
\item Conclure : $\dim F$ est le nombre de vecteurs de cette base.
\end{enumerate}
Cas usuels : une \textbf{droite vectorielle} (engendrée par un vecteur non nul) est de dimension $1$ ; un \textbf{plan vectoriel} (engendré par deux vecteurs non colinéaires) est de dimension $2$.
\end{methode}

\begin{exoresolu}[Dimension d'un sous-espace de $\mathbb{R}^3$]
On considère dans $\mathbb{R}^3$ l'ensemble $F = \{\,(x,y,z) \in \mathbb{R}^3 \mid x + y + z = 0\,\}$. On admet que $F$ est un sous-espace vectoriel de $\mathbb{R}^3$. Déterminer une base de $F$ et sa dimension.
\tcblower
Un vecteur $(x,y,z) \in F$ vérifie $z = -x - y$, où $x$ et $y$ sont libres. On écrit donc :
\[ (x,y,z) = (x,\ y,\ -x-y) = x(1,0,-1) + y(0,1,-1). \]
Posons $u = (1,0,-1)$ et $v = (0,1,-1)$. Tout vecteur de $F$ est combinaison linéaire de $u$ et $v$ : la famille $(u,v)$ est \textbf{génératrice} de $F$.

Montrons qu'elle est libre : si $\alpha u + \beta v = (0,0,0)$, alors $(\alpha, \beta, -\alpha-\beta) = (0,0,0)$, ce qui donne immédiatement $\alpha = 0$ et $\beta = 0$. La famille $(u,v)$ est donc \textbf{libre}.

Étant libre et génératrice, $(u,v)$ est une \textbf{base} de $F$, et par conséquent :
\[ \dim F = 2. \]
Géométriquement, $F$ est un plan vectoriel de $\mathbb{R}^3$ : il est normal qu'il soit de dimension $2$.
\end{exoresolu}

\begin{savaistu}[Pourquoi la dimension est-elle si importante ?]
La notion de dimension, formalisée au XIX\textsuperscript{e} siècle par Hermann Grassmann puis Giuseppe Peano, est l'une des plus fécondes des mathématiques. Elle permet de classifier les espaces vectoriels : deux espaces de même dimension se ressemblent « exactement » (on dira en section suivante qu'ils sont isomorphes). Dans la vie courante, la dimension est partout : un écran affiche des images en dimension $2$, les ingénieurs de CAMTEL modélisent les signaux dans des espaces de grande dimension, et les algorithmes de compression (photos, musique) reposent sur le choix de bonnes bases dans lesquelles l'information s'exprime avec très peu de coefficients non nuls.
\end{savaistu}

\begin{exercice}[À toi de jouer]
\begin{enumerate}
\item Montrer que la famille $\big((1,2), (3,1)\big)$ est une base de $\mathbb{R}^2$, puis déterminer les coordonnées du vecteur $(5,5)$ dans cette base.
\item Dans $\mathbb{R}^3$, on pose $G = \{\,(x,y,z) \in \mathbb{R}^3 \mid x = 2y \text{ et } z = -y\,\}$. Montrer que $G$ est un sous-espace vectoriel de dimension $1$ et en donner une base.
\item La famille $\big((1,0,0), (1,1,0), (1,1,1)\big)$ est-elle une base de $\mathbb{R}^3$ ? Justifier.
\end{enumerate}
\end{exercice}

\begin{corrige}[Exercices]
\textbf{1.} $\dim \mathbb{R}^2 = 2$ et la famille a $2$ vecteurs : montrons qu'elle est libre. Si $\alpha(1,2) + \beta(3,1) = (0,0)$, alors $\alpha + 3\beta = 0$ et $2\alpha + \beta = 0$. De la première équation, $\alpha = -3\beta$ ; en reportant : $-6\beta + \beta = -5\beta = 0$, donc $\beta = 0$ puis $\alpha = 0$. La famille est libre, donc c'est une base de $\mathbb{R}^2$.

Pour $(5,5)$ : on résout $a(1,2) + b(3,1) = (5,5)$, soit $a + 3b = 5$ et $2a + b = 5$. De la seconde, $b = 5 - 2a$ ; en reportant : $a + 15 - 6a = 5$, soit $-5a = -10$, $a = 2$ et $b = 1$. Vérification : $2(1,2) + (3,1) = (5,5)$. \checkmark{} Les coordonnées sont $(2,1)$.

\textbf{2.} Un vecteur de $G$ s'écrit $(2y, y, -y) = y(2,1,-1)$. Donc $G$ est l'ensemble des multiples de $u = (2,1,-1)$ : c'est la droite vectorielle engendrée par $u$, qui est un sous-espace vectoriel de $\mathbb{R}^3$. Comme $u \neq (0,0,0)$, la famille $(u)$ est libre et génératrice de $G$ : c'est une base de $G$, et $\dim G = 1$.

\textbf{3.} $\dim \mathbb{R}^3 = 3$ et la famille a $3$ vecteurs : il suffit de montrer qu'elle est libre. Si $\alpha(1,0,0) + \beta(1,1,0) + \gamma(1,1,1) = (0,0,0)$, alors :
\[ \begin{cases} \alpha + \beta + \gamma = 0 \\ \beta + \gamma = 0 \\ \gamma = 0 \end{cases} \]
En remontant : $\gamma = 0$, puis $\beta = 0$, puis $\alpha = 0$. La famille est libre à $3$ vecteurs dans un espace de dimension $3$ : c'est une \textbf{base} de $\mathbb{R}^3$.
\end{corrige}

\begin{aretenir}[L'essentiel de la section]
\begin{itemize}
\item Une \textbf{base} est une famille à la fois \textbf{libre} et \textbf{génératrice}.
\item Dans une base, tout vecteur se décompose de \textbf{manière unique} : les coefficients sont ses \textbf{coordonnées}.
\item Toutes les bases d'un même espace ont le \textbf{même nombre} de vecteurs : c'est la \textbf{dimension} de l'espace (théorème admis).
\item $\dim \mathbb{R} = 1$, $\dim \mathbb{R}^2 = 2$, $\dim \mathbb{R}^3 = 3$ ; le plan vectoriel est de dimension $2$, l'espace vectoriel de dimension $3$.
\item En dimension $n$ connue, toute famille libre (ou génératrice) de $n$ vecteurs est une base : une seule vérification suffit.
\item Pour trouver la dimension d'un sous-espace : exhiber une base en décomposant le vecteur générique.
\end{itemize}
\end{aretenir}

\coursec{Applications linéaires}

Dans la section précédente, nous avons appris à munir un espace vectoriel d'une \cle{base} et à mesurer sa \cle{dimension} : nous savons désormais décrire les vecteurs d'un espace par leurs coordonnées. Il est temps de faire \emph{voyager} ces vecteurs. Dans la vie courante, on transforme sans cesse des quantités : un commerçant de Mokolo applique une remise proportionnelle à tous ses prix, un opérateur téléphonique convertit des forfaits, un cartographe du Cameroun déforme le plan pour l'aplatir sur une feuille. En mathématiques, les transformations les plus simples et les plus utiles sont celles qui \emph{respectent la structure d'espace vectoriel} : elles envoient une somme sur une somme, et un multiple sur le même multiple. Ce sont les \cle{applications linéaires}, héroïnes de cette section.

\courssub{Définition et premières propriétés}

\textbf{Intuition.}\par
Imagine une machine $f$ qui prend des vecteurs de $E$ et fabrique des vecteurs de $F$. Cette machine est dite \emph{linéaire} si elle « commute » avec les deux opérations de l'espace vectoriel : additionner \emph{avant} ou \emph{après} passage dans la machine donne le même résultat, et multiplier par un réel \emph{avant} ou \emph{après} aussi. Autrement dit, la machine ne casse rien de la structure : elle la transporte.

\begin{definition}[Application linéaire]
Soient $E$ et $F$ deux espaces vectoriels sur $\mathbb{R}$. Une application $f : E \to F$ est dite \cle{linéaire} si :
\begin{enumerate}
\item pour tous vecteurs $x, y \in E$ : $f(x + y) = f(x) + f(y)$ ;
\item pour tout vecteur $x \in E$ et tout réel $\lambda \in \mathbb{R}$ : $f(\lambda x) = \lambda f(x)$.
\end{enumerate}
On note $\mathcal{L}(E, F)$ l'ensemble des applications linéaires de $E$ dans $F$.
\end{definition}

En pratique, au lieu de vérifier deux conditions séparément, on les regroupe en une seule.

\begin{propriete}[Caractérisation des applications linéaires -- Démonstration exigible au Probatoire]
Soient $E$ et $F$ deux espaces vectoriels sur $\mathbb{R}$ et $f : E \to F$ une application. Alors :
\[ f \text{ est linéaire} \iff \text{pour tous } x, y \in E \text{ et tout } \lambda \in \mathbb{R}, \quad f(\lambda x + y) = \lambda f(x) + f(y). \]
\end{propriete}

\begin{experience}[Démonstration de la caractérisation]
\textbf{Sens direct.} Supposons $f$ linéaire. Soient $x, y \in E$ et $\lambda \in \mathbb{R}$. Alors :
\[ f(\lambda x + y) = f(\lambda x) + f(y) \quad \text{(condition 1)} = \lambda f(x) + f(y) \quad \text{(condition 2)}. \]

\textbf{Sens réciproque.} Supposons que $f(\lambda x + y) = \lambda f(x) + f(y)$ pour tous $x, y \in E$ et tout $\lambda \in \mathbb{R}$.
\begin{enumerate}
\item On prend $x = y = 0_E$ et $\lambda = 1$ : $f(0_E) = f(0_E) + f(0_E)$, d'où $f(0_E) = 0_F$.
\item On prend $y = 0_E$ : $f(\lambda x) = \lambda f(x) + f(0_E) = \lambda f(x)$.
\item On prend $\lambda = 1$ : $f(x + y) = f(x) + f(y)$.
\end{enumerate}
Les deux conditions de la définition sont donc retrouvées, et l'équivalence est démontrée. $\blacksquare$
\end{experience}

\begin{propriete}[Conséquences immédiates -- Démonstration exigible au Probatoire pour $f(0_E)=0_F$]
Soit $f : E \to F$ une application linéaire. Alors :
\begin{enumerate}
\item $f(0_E) = 0_F$ ;
\item pour tout $x \in E$, $f(-x) = -f(x)$ ;
\item $f$ conserve les combinaisons linéaires : pour tous $x_1, \dots, x_n \in E$ et tous réels $\lambda_1, \dots, \lambda_n$,
\[ f(\lambda_1 x_1 + \lambda_2 x_2 + \dots + \lambda_n x_n) = \lambda_1 f(x_1) + \lambda_2 f(x_2) + \dots + \lambda_n f(x_n). \]
\end{enumerate}
\end{propriete}

\begin{experience}[Démonstration : $f(0_E) = 0_F$]
Cette preuve est courte, classique et \emph{exigible}. Puisque $f$ est linéaire, pour tout $x \in E$ :
\[ f(0_E) = f(0 \cdot x) = 0 \cdot f(x) = 0_F. \]
On peut aussi écrire : $f(0_E) = f(0_E + 0_E) = f(0_E) + f(0_E)$, puis en ajoutant l'opposé de $f(0_E)$ aux deux membres : $0_F = f(0_E)$. $\blacksquare$
\end{experience}

\begin{aretenir}[Critère rapide de non-linéarité]
La contraposée du point 1 fournit un test redoutablement efficace :
\[ \text{si } f(0_E) \neq 0_F, \text{ alors } f \text{ n'est pas linéaire.} \]
Attention à la réciproque : $f(0_E) = 0_F$ ne prouve \textbf{pas} que $f$ est linéaire (voir le contre-exemple $f(x,y) = (x^2, y)$ plus bas).
\end{aretenir}

\begin{popfigure}
\begin{tikzpicture}[scale=1.1]
% Espace E
\draw[thick, popblue] (0,0) ellipse (1.6 and 1.1);
\node[popblue, font=\bfseries] at (-1.1,0.85) {$E$};
\fill[popdark] (0,-0.2) circle (1.5pt) node[below left, popdark] {$0_E$};
\fill[poporange] (-0.5,0.3) circle (1.5pt) node[above left] {$x$};
\fill[popgreen] (0.6,0.4) circle (1.5pt) node[above right] {$y$};
% Espace F
\draw[thick, poppurple] (6,0) ellipse (1.6 and 1.1);
\node[poppurple, font=\bfseries] at (4.9,0.85) {$F$};
\fill[popdark] (6,-0.2) circle (1.5pt) node[below left, popdark] {$0_F$};
\fill[poporange] (5.5,0.3) circle (1.5pt) node[above left] {$f(x)$};
\fill[popgreen] (6.7,0.35) circle (1.5pt) node[above right] {$f(y)$};
% Flèches
\draw[-{Stealth}, very thick, popink] (1.8,0.6) to[bend left=20] node[above] {$f$ linéaire} (4.2,0.6);
\draw[-{Stealth}, thick, popdark, dashed] (0.15,-0.25) to[bend right=25] node[below] {$f(0_E)=0_F$} (5.8,-0.3);
\end{tikzpicture}
\legende{Une application linéaire $f : E \to F$ envoie toujours le vecteur nul de $E$ sur le vecteur nul de $F$.}
\end{popfigure}

\courssub{Exemples et contre-exemples}

\begin{exemplebox}[Exemples fondamentaux d'applications linéaires]
\begin{enumerate}
\item \textbf{Homothéties vectorielles.} Pour $k \in \mathbb{R}$ fixé, $h : E \to E$, $x \mapsto kx$ est linéaire : $h(\lambda x + y) = k(\lambda x + y) = \lambda (kx) + ky = \lambda h(x) + h(y)$.
\item \textbf{Projection sur un axe.} $p : \mathbb{R}^2 \to \mathbb{R}^2$, $(x, y) \mapsto (x, 0)$ est linéaire : pour $u = (x, y)$, $v = (x', y')$ et $\lambda \in \mathbb{R}$,
\[ p(\lambda u + v) = p(\lambda x + x',\ \lambda y + y') = (\lambda x + x',\ 0) = \lambda (x, 0) + (x', 0) = \lambda p(u) + p(v). \]
\item \textbf{Symétrie vectorielle.} $s : \mathbb{R}^2 \to \mathbb{R}^2$, $(x, y) \mapsto (x, -y)$ (symétrie par rapport à l'axe des abscisses) est linéaire (vérification identique).
\item \textbf{Application nulle} $x \mapsto 0_F$ et \textbf{identité} $x \mapsto x$ sont linéaires.
\end{enumerate}
\end{exemplebox}

\begin{exemplebox}[Contre-exemples : applications NON linéaires]
\begin{enumerate}
\item $f : \mathbb{R} \to \mathbb{R}$, $x \mapsto x + 1$. On a $f(0) = 1 \neq 0$ : par le critère rapide, $f$ \textbf{n'est pas linéaire}. (C'est une application \emph{affine}.)
\item $g : \mathbb{R}^2 \to \mathbb{R}^2$, $(x, y) \mapsto (x^2, y)$. Pourtant $g(0,0) = (0,0)$ ! Mais $g(2 \cdot (1, 0)) = g(2, 0) = (4, 0)$ alors que $2\,g(1, 0) = 2(1, 0) = (2, 0)$. Comme $(4,0) \neq (2,0)$, $g$ \textbf{n'est pas linéaire}. Cet exemple montre que le test $f(0_E) = 0_F$ est nécessaire mais \emph{pas suffisant}.
\end{enumerate}
\end{exemplebox}

\begin{attention}[Erreur fréquente : « ça marche sur mes exemples, donc c'est linéaire »]
Vérifier $f(\lambda x + y) = \lambda f(x) + f(y)$ sur deux ou trois exemples numériques ne prouve \textbf{rien} : la linéarité est une propriété \emph{universelle}, qui doit valoir pour \emph{tous} les vecteurs et \emph{tous} les réels. La preuve exige des éléments \textbf{quelconques} $x, y \in E$ et $\lambda \in \mathbb{R}$, manipulés littéralement. En revanche, pour montrer qu'une application n'est \emph{pas} linéaire, un seul contre-exemple numérique suffit.
\end{attention}

\begin{methode}[Montrer qu'une application est linéaire]
\begin{enumerate}
\item Identifier clairement l'espace de départ $E$ et l'espace d'arrivée $F$.
\item Prendre deux éléments \textbf{quelconques} $u, v \in E$ (avec des coordonnées littérales si $E = \mathbb{R}^2$ ou $\mathbb{R}^3$) et un réel quelconque $\lambda$.
\item Calculer $f(\lambda u + v)$ en développant à partir de la définition de $f$.
\item Calculer séparément $\lambda f(u) + f(v)$.
\item Conclure par égalité des deux expressions, valable pour tous $u, v, \lambda$.
\end{enumerate}
\end{methode}

\begin{exoresolu}[Montrer qu'une application est un endomorphisme de $\mathbb{R}^2$]
\textbf{Énoncé.} On considère l'application $f : \mathbb{R}^2 \to \mathbb{R}^2$ définie par $f(x, y) = (2x - y,\ x + 3y)$. Montrer que $f$ est un endomorphisme de $\mathbb{R}^2$.
\tcblower
\textbf{Solution.} $f$ va bien de $\mathbb{R}^2$ dans $\mathbb{R}^2$ ; montrons qu'elle est linéaire.

Soient $u = (x, y)$ et $v = (x', y')$ deux éléments \emph{quelconques} de $\mathbb{R}^2$, et $\lambda \in \mathbb{R}$ quelconque. Alors $\lambda u + v = (\lambda x + x',\ \lambda y + y')$, donc :
\begin{align*}
f(\lambda u + v) &= \big(2(\lambda x + x') - (\lambda y + y'),\ (\lambda x + x') + 3(\lambda y + y')\big) \\
&= \big(\lambda(2x - y) + (2x' - y'),\ \lambda(x + 3y) + (x' + 3y')\big) \\
&= \lambda\,(2x - y,\ x + 3y) + (2x' - y',\ x' + 3y') \\
&= \lambda f(u) + f(v).
\end{align*}
Cette égalité vaut pour tous $u, v \in \mathbb{R}^2$ et tout $\lambda \in \mathbb{R}$ : $f$ est donc \cle{linéaire}. Comme de plus l'espace de départ et l'espace d'arrivée sont le même espace $\mathbb{R}^2$, $f$ est un \cle{endomorphisme} de $\mathbb{R}^2$. $\blacksquare$
\end{exoresolu}

\courssub{Vocabulaire : endomorphisme, isomorphisme, automorphisme}

\begin{definition}[Endomorphisme, isomorphisme, automorphisme]
Soient $E$ et $F$ deux espaces vectoriels sur $\mathbb{R}$ et $f : E \to F$ une application linéaire.
\begin{itemize}
\item Si $E = F$, on dit que $f$ est un \cle{endomorphisme} de $E$. L'ensemble des endomorphismes de $E$ se note $\mathcal{L}(E)$.
\item Si $f$ est \emph{bijective}, on dit que $f$ est un \cle{isomorphisme} de $E$ sur $F$ ; on dit alors que $E$ et $F$ sont \emph{isomorphes}.
\item Un endomorphisme bijectif ($E = F$ et $f$ bijective) est appelé \cle{automorphisme} de $E$.
\end{itemize}
\end{definition}

\begin{exemplebox}[Illustrations du vocabulaire]
\begin{enumerate}
\item $f(x, y) = (2x - y,\ x + 3y)$ est un endomorphisme de $\mathbb{R}^2$ (vu ci-dessus) ; on verra qu'il est bijectif, donc c'est un automorphisme.
\item $g : \mathbb{R}^2 \to \mathbb{R}^2$, $(x, y) \mapsto (x, 0)$ est un endomorphisme \emph{non} bijectif : $(0, 5)$ et $(0, 7)$ ont la même image $(0,0)$.
\item $\varphi : \mathbb{R}^3 \to \mathbb{R}^3$, $(x, y, z) \mapsto (x + z,\ y,\ z)$ est un automorphisme de $\mathbb{R}^3$.
\end{enumerate}
\end{exemplebox}

\begin{methode}[Montrer qu'une application est un endomorphisme]
Pour affirmer que $f$ est un \cle{endomorphisme}, il faut vérifier \emph{deux} choses, et le dire explicitement dans la rédaction :
\begin{enumerate}
\item $f$ est une application de $E$ dans $E$ (\textbf{même} espace au départ et à l'arrivée) ;
\item $f$ est linéaire : $f(\lambda u + v) = \lambda f(u) + f(v)$ pour tous $u, v \in E$ et tout $\lambda \in \mathbb{R}$.
\end{enumerate}
Oublier la première condition est une erreur classique : une application linéaire de $\mathbb{R}^2$ dans $\mathbb{R}^3$ n'est \emph{pas} un endomorphisme.
\end{methode}

\courssub{Expression analytique d'une application linéaire dans une base}

\textbf{Intuition.}\par
Voici la bonne nouvelle qui couronne tout le travail sur les bases : une application linéaire est \emph{entièrement déterminée} par les images des vecteurs d'une base. En effet, si $(e_1, e_2)$ est une base de $E$, tout vecteur $x$ s'écrit $x = x_1 e_1 + x_2 e_2$, et par linéarité :
\[ f(x) = x_1 f(e_1) + x_2 f(e_2). \]
Connaître $f(e_1)$ et $f(e_2)$ suffit donc à connaître $f$ partout !

\begin{propriete}[Expression analytique dans $\mathbb{R}^2$]
Toute application linéaire $f : \mathbb{R}^2 \to \mathbb{R}^2$ s'écrit de manière unique sous la forme
\[ f(x, y) = (ax + by,\ cx + dy) \]
où $a, b, c, d$ sont des réels. Réciproquement, toute application de cette forme est linéaire.
\end{propriete}

\begin{experience}[Démonstration (formatrice)]
Soit $(e_1, e_2)$ la base canonique de $\mathbb{R}^2$, avec $e_1 = (1, 0)$ et $e_2 = (0, 1)$. Posons $f(e_1) = (a, c)$ et $f(e_2) = (b, d)$. Tout vecteur $(x, y)$ s'écrit $(x, y) = x e_1 + y e_2$. Par linéarité :
\[ f(x, y) = x f(e_1) + y f(e_2) = x(a, c) + y(b, d) = (ax + by,\ cx + dy). \]
Réciproquement, si $f(x, y) = (ax + by, cx + dy)$, la vérification de $f(\lambda u + v) = \lambda f(u) + f(v)$ se fait exactement comme dans l'exercice résolu précédent. $\blacksquare$
\end{experience}

\begin{aretenir}[Généralisation à $\mathbb{R}^3$]
Toute application linéaire $f : \mathbb{R}^3 \to \mathbb{R}^3$ s'écrit de manière unique
\[ f(x, y, z) = (ax + by + cz,\ a'x + b'y + c'z,\ a''x + b''y + c''z) \]
où les neuf coefficients sont des réels : chaque coordonnée de l'image est une \emph{combinaison linéaire} des coordonnées de départ, \textbf{sans terme constant}. De même, une application linéaire de $\mathbb{R}$ dans $\mathbb{R}$ est nécessairement de la forme $f(x) = ax$.
\end{aretenir}

\begin{attention}[Reconnaître une expression non linéaire]
Une expression analytique contenant un \textbf{terme constant} ($(2x - y + 3,\ x)$), un \textbf{carré} ou un \textbf{produit de coordonnées} ($(x^2, y)$, $(xy, x + y)$), une \textbf{valeur absolue} ne peut pas être linéaire. La forme imposée est : chaque coordonnée de l'image est une combinaison linéaire de $x, y$ (et $z$), sans constante.
\end{attention}

\begin{popfigure}
\begin{tikzpicture}[scale=0.85]
% Quadrillage de départ
\begin{scope}[shift={(0,0)}]
\draw[step=0.5, popline, very thin] (-1.5,-1.5) grid (1.5,1.5);
\draw[-{Stealth}, thick, popdark] (-1.6,0) -- (1.7,0);
\draw[-{Stealth}, thick, popdark] (0,-1.6) -- (0,1.7);
\draw[-{Stealth}, very thick, popblue] (0,0) -- (1,0) node[below] {$e_1$};
\draw[-{Stealth}, very thick, popblue] (0,0) -- (0,1) node[left] {$e_2$};
\fill[popdark] (0,0) circle (2pt) node[below left] {$O$};
\node[font=\bfseries, popdark] at (0,-2) {Avant : base canonique};
\end{scope}
% Flèche de transformation
\draw[-{Stealth}, ultra thick, poporange] (2.2,0) -- (3.8,0) node[midway, above] {$f$};
% Quadrillage transformé : f(e1)=(1.2,0.5), f(e2)=(-0.4,1.1)
\begin{scope}[shift={(6.5,0)}]
\pgfmathsetmacro{\ax}{1.2}\pgfmathsetmacro{\ay}{0.5}
\pgfmathsetmacro{\bx}{-0.4}\pgfmathsetmacro{\by}{1.1}
\foreach \i in {-2,-1,0,1,2}{
  \draw[popline, very thin] ({\i*\ax -2*\bx},{\i*\ay -2*\by}) -- ({\i*\ax +2*\bx},{\i*\ay +2*\by});
  \draw[popline, very thin] ({-2*\ax +\i*\bx},{-2*\ay +\i*\by}) -- ({2*\ax +\i*\bx},{2*\ay +\i*\by});
}
\draw[-{Stealth}, very thick, poporange] (0,0) -- (\ax,\ay) node[right] {$f(e_1)$};
\draw[-{Stealth}, very thick, poporange] (0,0) -- (\bx,\by) node[above] {$f(e_2)$};
\fill[popdark] (0,0) circle (2pt) node[below left] {$O=f(O)$};
\node[font=\bfseries, popdark] at (0.3,-2) {Après : parallélogrammes};
\end{scope}
\end{tikzpicture}
\legende{Effet d'une application linéaire du plan : le quadrillage de carrés devient un quadrillage de parallélogrammes, les droites restent des droites, les parallèles restent parallèles, et l'origine reste fixe.}
\end{popfigure}

\begin{savaistu}[Pourquoi « linéaire » ?]
Le mot vient du latin \emph{linea}, la ligne : une application linéaire du plan transforme toute droite en une droite (ou un point), et le quadrillage régulier d'une feuille en un quadrillage de parallélogrammes. C'est exactement ce que font les logiciels de cartographie utilisés pour établir les plans des villes de Douala ou Yaoundé : rotations, homothéties, projections sont des applications linéaires. En économie aussi, passer d'une production de $x$ tonnes de cacao et $y$ tonnes de café à un chiffre d'affaires $(p_1 x + p_2 y)$ est un modèle linéaire : les recettes s'ajoutent comme les productions.
\end{savaistu}

\begin{exercice}[À toi de jouer]
\begin{enumerate}
\item Montrer que $f : \mathbb{R}^2 \to \mathbb{R}^2$, $(x, y) \mapsto (x + 2y,\ 3x - y)$ est un endomorphisme de $\mathbb{R}^2$.
\item L'application $g : \mathbb{R}^2 \to \mathbb{R}^2$, $(x, y) \mapsto (x + 1,\ y)$ est-elle linéaire ? Justifier.
\item Déterminer l'expression analytique de l'endomorphisme $h$ de $\mathbb{R}^2$ tel que $h(1, 0) = (2, -1)$ et $h(0, 1) = (1, 3)$.
\end{enumerate}
\end{exercice}

\begin{corrige}[Exercices 1 à 3]
\begin{enumerate}
\item Soient $u = (x, y)$, $v = (x', y')$ quelconques dans $\mathbb{R}^2$ et $\lambda \in \mathbb{R}$. Alors $\lambda u + v = (\lambda x + x', \lambda y + y')$ et
\begin{align*}
f(\lambda u + v) &= \big((\lambda x + x') + 2(\lambda y + y'),\ 3(\lambda x + x') - (\lambda y + y')\big) \\
&= \lambda(x + 2y,\ 3x - y) + (x' + 2y',\ 3x' - y') = \lambda f(u) + f(v).
\end{align*}
$f$ est linéaire de $\mathbb{R}^2$ dans $\mathbb{R}^2$ : c'est un endomorphisme de $\mathbb{R}^2$.
\item $g(0, 0) = (1, 0) \neq (0, 0)$. Par le critère rapide, $g$ n'est \textbf{pas} linéaire (c'est une translation).
\item Par linéarité : $h(x, y) = x\,h(1,0) + y\,h(0,1) = x(2, -1) + y(1, 3) = (2x + y,\ -x + 3y)$.
\end{enumerate}
\end{corrige}

\begin{aretenir}[L'essentiel de la section]
\begin{itemize}
\item $f : E \to F$ est \cle{linéaire} si $f(\lambda x + y) = \lambda f(x) + f(y)$ pour tous $x, y \in E$, $\lambda \in \mathbb{R}$ — à prouver sur des éléments \emph{quelconques}, jamais sur des exemples.
\item Toute application linéaire vérifie $f(0_E) = 0_F$ ; si $f(0_E) \neq 0_F$, elle n'est pas linéaire.
\item \cle{Endomorphisme} : $E = F$ ; \cle{isomorphisme} : linéaire et bijective ; \cle{automorphisme} : les deux.
\item Une application linéaire est entièrement déterminée par les images des vecteurs d'une base ; dans $\mathbb{R}^2$, elle s'écrit $f(x, y) = (ax + by,\ cx + dy)$.
\end{itemize}
Dans la prochaine sous-section, nous associerons à toute application linéaire deux sous-espaces vectoriels remarquables — son \cle{noyau} et son \cle{image} — qui mesureront précisément ce qu'elle « écrase » et ce qu'elle « atteint ».
\end{aretenir}

\coursec{Noyau et image d'une application linéaire}

Dans la section précédente, nous avons appris à reconnaître une application linéaire $f : E \to F$ : c'est une application qui respecte l'addition et la multiplication par un réel. Nous avons aussi vu comment l'exprimer dans une base. Mais une question fondamentale reste ouverte : \emph{que fait réellement $f$ des vecteurs de $E$ ?} Certains vecteurs sont « écrasés » sur le vecteur nul de $F$ ; d'autres atteignent des vecteurs bien précis de $F$, tandis que certains vecteurs de $F$ ne sont jamais atteints. Deux ensembles vont nous permettre de photographier cette action : le \cle{noyau} (les vecteurs écrasés) et l'\cle{image} (les vecteurs atteints). Ces deux notions sont la clé pour décider si une application linéaire est injective, surjective ou bijective — une compétence exigible au Probatoire.

\courssub{Le noyau d'une application linéaire}

\textbf{Intuition.} Imagine un projecteur qui éclaire un mur. Tous les objets situés sur un même rayon lumineux donnent la même ombre sur le mur ; en particulier, tous les points du rayon qui passe par la source donnent une ombre nulle, confondue avec la source. L'ensemble des directions « perdues » par la projection, c'est exactement l'idée du noyau.

\begin{definition}[Noyau d'une application linéaire]
Soit $f : E \to F$ une application linéaire entre deux espaces vectoriels sur $\mathbb{R}$. On appelle \cle{noyau} de $f$, noté $\ker f$, l'ensemble des vecteurs de $E$ dont l'image est le vecteur nul de $F$ :
\[ \ker f = \{ x \in E \;/\; f(x) = 0_F \}. \]
\end{definition}

\begin{attention}[Le noyau n'est jamais vide]
Erreur très fréquente en classe : écrire « $\ker f = \varnothing$ » ou confondre $\ker f = \{0_E\}$ avec « le noyau n'existe pas ». Or toute application linéaire vérifie $f(0_E) = 0_F$, donc $0_E \in \ker f$ \textbf{toujours}. Le noyau contient \emph{au moins} le vecteur nul. Dire $\ker f = \{0_E\}$ signifie que le vecteur nul est le \emph{seul} vecteur écrasé — c'est le cas le plus favorable, comme nous le verrons.
\end{attention}

\begin{propriete}[Le noyau est un sous-espace vectoriel]
Pour toute application linéaire $f : E \to F$, $\ker f$ est un \cle{sous-espace vectoriel} de $E$.
\end{propriete}

\begin{experience}[Démonstration guidée : $\ker f$ est un sous-espace vectoriel de $E$]
Montrons les trois points de la caractérisation d'un sous-espace vectoriel.
\begin{enumerate}
\item \textbf{Non vide :} $f(0_E) = 0_F$ car $f$ est linéaire, donc $0_E \in \ker f$.
\item \textbf{Stabilité par addition :} soient $u, v \in \ker f$. Alors $f(u) = 0_F$ et $f(v) = 0_F$. Par linéarité :
\[ f(u+v) = f(u) + f(v) = 0_F + 0_F = 0_F, \]
donc $u + v \in \ker f$.
\item \textbf{Stabilité par multiplication par un réel :} soient $u \in \ker f$ et $\lambda \in \mathbb{R}$. Alors
\[ f(\lambda u) = \lambda f(u) = \lambda \cdot 0_F = 0_F, \]
donc $\lambda u \in \ker f$.
\end{enumerate}
Conclusion : $\ker f$ est un sous-espace vectoriel de $E$. $\square$
\end{experience}

\textbf{Conséquence pratique capitale.} Puisque $\ker f$ est un espace vectoriel, on peut en chercher une \emph{base}. Concrètement, déterminer $\ker f$ revient à résoudre le système homogène $f(x) = 0_F$ : c'est ce qu'on appelle l'\cle{équation caractéristique} du noyau.

\begin{methode}[Déterminer une équation caractéristique et une base du noyau]
\begin{enumerate}
\item Écrire $f(x) = 0_F$, c'est-à-dire traduire l'écriture analytique de $f$ en un \textbf{système homogène} sur les coordonnées de $x$.
\item Résoudre ce système : exprimer les inconnues « principales » en fonction des inconnues libres (paramètres).
\item Écrire la solution générale comme combinaison linéaire de vecteurs portés par les paramètres : $\ker f = \mathrm{Vect}(\ldots)$.
\item La famille obtenue est génératrice ; vérifier qu'elle est libre (elle l'est presque toujours automatiquement avec cette méthode) : c'est une \textbf{base de $\ker f$}.
\end{enumerate}
\end{methode}

\begin{exemplebox}[Noyau dans $\mathbb{R}^3$]
Soit $f : \mathbb{R}^3 \to \mathbb{R}^2$ définie par $f(x,y,z) = (x + y - z,\; 2x - y + z)$. Cherchons $\ker f$.
\[ f(x,y,z) = (0,0) \iff \begin{cases} x + y - z = 0 \\ 2x - y + z = 0 \end{cases} \]
En additionnant les deux lignes : $3x = 0$, donc $x = 0$. La première ligne donne alors $z = y$. Posons $y = t$, $t \in \mathbb{R}$. Les solutions sont les triplets $(0, t, t) = t(0,1,1)$. Donc
\[ \ker f = \mathrm{Vect}\big((0,1,1)\big), \]
droite vectorielle de base $\big((0,1,1)\big)$, de dimension $1$.
\end{exemplebox}

\begin{propriete}[Critère d'injectivité — démonstration exigible]
Soit $f : E \to F$ une application linéaire. Alors :
\[ f \text{ est injective} \iff \ker f = \{0_E\}. \]
\end{propriete}

\begin{experience}[Démonstration guidée du critère d'injectivité]
On démontre les deux implications.
\begin{enumerate}
\item[$\Rightarrow$] Supposons $f$ injective. Soit $x \in \ker f$ : alors $f(x) = 0_F$. Mais on sait aussi que $f(0_E) = 0_F$. Donc $f(x) = f(0_E)$, et par injectivité $x = 0_E$. Ainsi $\ker f = \{0_E\}$.
\item[$\Leftarrow$] Supposons $\ker f = \{0_E\}$. Soient $u, v \in E$ tels que $f(u) = f(v)$. Par linéarité :
\[ f(u - v) = f(u) - f(v) = 0_F, \]
donc $u - v \in \ker f = \{0_E\}$, c'est-à-dire $u - v = 0_E$, donc $u = v$. L'application $f$ est injective. $\square$
\end{enumerate}
Remarque bien l'étape clé de la réciproque : c'est la \emph{linéarité} qui permet de ramener l'égalité $f(u) = f(v)$ à une appartenance au noyau. C'est pour cela que ce critère, faux pour une application quelconque, est vrai pour les applications linéaires.
\end{experience}

\courssub{L'image d'une application linéaire}

\textbf{Intuition.} Si le noyau rassemble ce qui est « perdu », l'image rassemble ce qui est « atteint ». Reprends le projecteur : l'ensemble de toutes les ombres possibles sur le mur est une partie du mur — c'est l'image de la projection.

\begin{definition}[Image d'une application linéaire]
Soit $f : E \to F$ une application linéaire. On appelle \cle{image} de $f$, notée $\mathrm{Im}\,f$ (ou $f(E)$), l'ensemble des vecteurs de $F$ qui sont l'image d'au moins un vecteur de $E$ :
\[ \mathrm{Im}\,f = \{ f(x),\; x \in E \} = \{ y \in F \;/\; \exists\, x \in E,\; f(x) = y \}. \]
\end{definition}

\begin{propriete}[L'image est un sous-espace vectoriel]
Pour toute application linéaire $f : E \to F$, $\mathrm{Im}\,f$ est un \cle{sous-espace vectoriel} de $F$.
\end{propriete}

\begin{experience}[Démonstration guidée : $\mathrm{Im}\,f$ est un sous-espace vectoriel de $F$]
\begin{enumerate}
\item \textbf{Non vide :} $0_F = f(0_E) \in \mathrm{Im}\,f$.
\item \textbf{Stabilité par addition :} soient $y_1, y_2 \in \mathrm{Im}\,f$. Il existe $x_1, x_2 \in E$ tels que $y_1 = f(x_1)$ et $y_2 = f(x_2)$. Alors
\[ y_1 + y_2 = f(x_1) + f(x_2) = f(x_1 + x_2) \in \mathrm{Im}\,f, \]
car $x_1 + x_2 \in E$.
\item \textbf{Stabilité par multiplication par un réel :} soient $y \in \mathrm{Im}\,f$ et $\lambda \in \mathbb{R}$. Il existe $x \in E$ avec $y = f(x)$, et
\[ \lambda y = \lambda f(x) = f(\lambda x) \in \mathrm{Im}\,f. \]
\end{enumerate}
Conclusion : $\mathrm{Im}\,f$ est un sous-espace vectoriel de $F$. $\square$
\end{experience}

\begin{propriete}[L'image est engendrée par les images des vecteurs d'une base]
Soit $f : E \to F$ une application linéaire et $(e_1, e_2, \ldots, e_n)$ une base de $E$. Alors
\[ \mathrm{Im}\,f = \mathrm{Vect}\big(f(e_1), f(e_2), \ldots, f(e_n)\big). \]
\emph{Justification :} tout $x \in E$ s'écrit $x = x_1 e_1 + \cdots + x_n e_n$, donc par linéarité $f(x) = x_1 f(e_1) + \cdots + x_n f(e_n)$ : tout vecteur de $\mathrm{Im}\,f$ est combinaison linéaire des $f(e_i)$. Réciproquement, chaque $f(e_i)$ appartient à $\mathrm{Im}\,f$, donc toutes leurs combinaisons linéaires aussi (c'est un sous-espace vectoriel).
\end{propriete}

\begin{methode}[Déterminer une base de l'image]
\begin{enumerate}
\item Choisir une base $(e_1, \ldots, e_n)$ de $E$ (souvent la base canonique) et calculer $f(e_1), \ldots, f(e_n)$.
\item La famille $\big(f(e_1), \ldots, f(e_n)\big)$ engendre $\mathrm{Im}\,f$.
\item \textbf{Extraire une famille libre} : éliminer les vecteurs qui sont combinaisons linéaires des autres (par exemple les vecteurs colinéaires à un autre).
\item La famille restante, libre et génératrice, est une \textbf{base de $\mathrm{Im}\,f$} ; son nombre de vecteurs est $\dim \mathrm{Im}\,f$.
\end{enumerate}
Pour l'\cle{équation caractéristique de l'image}, on cherche les conditions sur $y \in F$ pour que le système $f(x) = y$ admette au moins une solution $x \in E$.
\end{methode}

\begin{popfigure}
\begin{tikzpicture}[scale=1.0]
% Patate E
\draw[thick, popblue, fill=popblueL] plot[smooth cycle, tension=0.8] coordinates {(0,0.2) (1.6,0) (2.4,1.2) (1.8,2.6) (0.4,2.8) (-0.4,1.4)};
\node[popblue, font=\bfseries] at (0.2,2.5) {$E$};
% Patate F
\draw[thick, poporange, fill=poporangeL] plot[smooth cycle, tension=0.8] coordinates {(6,0.2) (7.6,0) (8.4,1.2) (7.8,2.6) (6.4,2.8) (5.6,1.4)};
\node[poporange, font=\bfseries] at (6.2,2.5) {$F$};
% Ker f dans E
\draw[thick, poppurple, fill=poppurpleL] plot[smooth cycle, tension=0.8] coordinates {(0.5,0.5) (1.5,0.4) (1.9,1.1) (1.2,1.7) (0.4,1.3)};
\node[poppurple] at (1.15,1.05) {$\ker f$};
% Im f dans F
\draw[thick, popgreen, fill=popgreenL] plot[smooth cycle, tension=0.8] coordinates {(6.5,0.5) (7.5,0.4) (7.9,1.1) (7.2,1.7) (6.4,1.3)};
\node[popgreen!60!black] at (7.15,0.85) {$\mathrm{Im}\,f$};
% 0_F
\fill[popdark] (7.15,1.45) circle (0.05);
\node[popdark, above] at (7.15,1.5) {$0_F$};
% 0_E dans ker
\fill[popdark] (1.15,0.7) circle (0.05);
\node[popdark, below] at (1.15,0.68) {$0_E$};
% Flèche f
\draw[-{Stealth[length=3mm]}, very thick, popdark] (2.6,1.6) -- (5.4,1.6) node[midway, above]{$f$};
% Flèches ker -> 0_F
\draw[-{Stealth[length=2mm]}, poppurple, dashed] (1.9,1.2) to[bend left=15] (6.9,1.5);
\draw[-{Stealth[length=2mm]}, poppurple, dashed] (1.6,0.6) to[bend right=15] (6.95,1.35);
% Flèches E -> Im f
\draw[-{Stealth[length=2mm]}, popgreen!60!black] (2.2,2.2) to[bend left=20] (7.3,1.7);
\draw[-{Stealth[length=2mm]}, popgreen!60!black] (2.35,0.8) to[bend right=10] (6.7,0.55);
\end{tikzpicture}
\legende{Le noyau $\ker f$ (en violet) est entièrement envoyé sur $0_F$ ; l'image $\mathrm{Im}\,f$ (en vert) est la partie de $F$ effectivement atteinte par les flèches.}
\end{popfigure}

\courssub{Un exemple complet dans $\mathbb{R}^2$}

\begin{exoresolu}[Noyau, image et bijectivité de $f(x,y) = (x - y,\; 2x - 2y)$]
Soit $f : \mathbb{R}^2 \to \mathbb{R}^2$ définie par $f(x,y) = (x - y,\; 2x - 2y)$.
\begin{enumerate}
\item Montrer que $f$ est linéaire.
\item Déterminer $\ker f$, en donner une base et sa dimension.
\item Déterminer $\mathrm{Im}\,f$, en donner une base et sa dimension.
\item L'application $f$ est-elle injective ? surjective ? bijective ?
\end{enumerate}
\tcblower
\textbf{1. Linéarité.} Soient $u = (x,y)$, $v = (x',y')$ et $\lambda \in \mathbb{R}$.
\begin{align*}
f(u+v) &= f(x+x',\, y+y') = \big((x+x') - (y+y'),\; 2(x+x') - 2(y+y')\big) \\
&= (x - y,\; 2x - 2y) + (x' - y',\; 2x' - 2y') = f(u) + f(v) ;
\end{align*}
\[ f(\lambda u) = f(\lambda x, \lambda y) = (\lambda x - \lambda y,\; 2\lambda x - 2\lambda y) = \lambda (x - y,\; 2x - 2y) = \lambda f(u). \]
Donc $f$ est une application linéaire (c'est un \cle{endomorphisme} de $\mathbb{R}^2$).

\textbf{2. Noyau.} On résout le système homogène (équation caractéristique du noyau) :
\[ f(x,y) = (0,0) \iff \begin{cases} x - y = 0 \\ 2x - 2y = 0 \end{cases} \iff x = y. \]
(La deuxième équation est le double de la première : elle n'apporte rien.) Les solutions sont les couples $(t, t) = t(1,1)$, $t \in \mathbb{R}$. Donc
\[ \ker f = \{ (x,y) \in \mathbb{R}^2 \;/\; x = y \} = \mathrm{Vect}\big((1,1)\big). \]
Le noyau est la \textbf{droite vectorielle d'équation $x = y$}, de base $\big((1,1)\big)$, et $\dim \ker f = 1$.

\textbf{3. Image.} Utilisons la base canonique $(e_1, e_2)$ de $\mathbb{R}^2$ :
\[ f(e_1) = f(1,0) = (1, 2), \qquad f(e_2) = f(0,1) = (-1, -2). \]
Donc $\mathrm{Im}\,f = \mathrm{Vect}\big((1,2), (-1,-2)\big)$. Or $(-1,-2) = -(1,2)$ : les deux vecteurs sont colinéaires, on élimine le second. Ainsi
\[ \mathrm{Im}\,f = \mathrm{Vect}\big((1,2)\big), \]
droite vectorielle de base $\big((1,2)\big)$, de dimension $1$. Équation caractéristique : $(a,b) \in \mathrm{Im}\,f \iff$ il existe $(x,y)$ avec $a = x - y$ et $b = 2(x-y)$, c'est-à-dire $\boxed{b = 2a}$.

\textbf{4. Conclusion.} $\ker f \neq \{(0,0)\}$ (il contient par exemple $(1,1)$), donc $f$ \textbf{n'est pas injective}. $\mathrm{Im}\,f \neq \mathbb{R}^2$ (c'est une droite, pas le plan entier), donc $f$ \textbf{n'est pas surjective}. A fortiori, $f$ n'est pas bijective. Remarque au passage : $\dim \ker f + \dim \mathrm{Im}\,f = 1 + 1 = 2 = \dim \mathbb{R}^2$ — ce n'est pas un hasard, voir le « Le savais-tu ? » plus bas.
\end{exoresolu}

\begin{popfigure}
\begin{tikzpicture}[scale=1.1]
% Repère
\draw[-{Stealth[length=2.5mm]}, popmuted] (-2.6,0) -- (2.6,0) node[below left]{$x$};
\draw[-{Stealth[length=2.5mm]}, popmuted] (0,-2.6) -- (0,2.6) node[below left]{$y$};
% Image : droite de base (1,2)
\draw[very thick, popgreen] (-1.15,-2.3) -- (1.15,2.3);
\node[popgreen!60!black, right] at (1.0,2.15) {$\mathrm{Im}\,f$ : droite de base $(1,2)$};
% Noyau : droite x = y
\draw[very thick, poppurple, dashed] (-2.3,-2.3) -- (2.3,2.3);
\node[poppurple, above left] at (2.1,2.1) {$\ker f$ : $x = y$};
% Vecteur (1,1) et son image
\draw[-{Stealth[length=2.5mm]}, very thick, popblue] (0,0) -- (1,1) node[midway, below right]{\small $(1,1)$};
\draw[-{Stealth[length=2.5mm]}, very thick, poporange] (0,0) -- (0.5,1) node[pos=0.9, right]{\small $(1,2)$};
% Flèche de projection
\draw[-{Stealth[length=2mm]}, popdark, dotted] (1,1) -- (0.5,1);
\fill[popdark] (0,0) circle (0.06);
\node[popdark, below left] at (-0.05,-0.05) {$O$};
\end{tikzpicture}
\legende{Pour $f(x,y) = (x-y,\, 2x-2y)$ : tout vecteur de la droite $\ker f$ (en violet, pointillés) est écrasé sur $O$ ; toutes les images tombent sur la droite $\mathrm{Im}\,f$ (en vert). Le vecteur $(1,1)$ du noyau est envoyé sur $O$, tandis que $(1,0)$ est envoyé sur $(1,2)$.}
\end{popfigure}

\courssub{Noyau, image et bijectivité}

\begin{propriete}[Critères de bijectivité]
Soit $f : E \to F$ une application linéaire.
\begin{itemize}
\item $f$ est \textbf{bijective} si et seulement si $\ker f = \{0_E\}$ \textbf{et} $\mathrm{Im}\,f = F$.
\item Si de plus $\dim E = \dim F$ (en particulier pour un endomorphisme en dimension finie), alors les trois propriétés sont équivalentes (résultat admis) :
\[ f \text{ injective} \iff f \text{ surjective} \iff f \text{ bijective}. \]
Dans ce cas, un endomorphisme bijectif est appelé \cle{automorphisme}.
\end{itemize}
\end{propriete}

\begin{attention}[Bien choisir sa stratégie]
Pour montrer qu'un endomorphisme de $\mathbb{R}^2$ ou $\mathbb{R}^3$ est bijectif, \textbf{une seule} des deux vérifications suffit : soit $\ker f = \{0\}$, soit $\mathrm{Im}\,f = F$. Inutile de faire les deux ! En général, le calcul du noyau (système homogène) est le plus rapide. À l'inverse, pour montrer que $f$ n'est \emph{pas} bijective, exhiber un vecteur non nul du noyau est souvent le chemin le plus court.
\end{attention}

\begin{methode}[Conclure sur la bijectivité à partir de l'écriture analytique]
Trois approches au choix pour $f : E \to F$ avec $\dim E = \dim F$ :
\begin{enumerate}
\item \textbf{Par le noyau :} résoudre $f(x) = 0_F$. Si l'unique solution est $x = 0_E$, alors $f$ est injective, donc bijective.
\item \textbf{Par l'image :} montrer que les $f(e_i)$ forment une famille libre de $F$ (autant de vecteurs que $\dim F$) ; alors $\mathrm{Im}\,f = F$ et $f$ est bijective.
\item \textbf{Par résolution directe :} pour $y \in F$ quelconque, résoudre le système $f(x) = y$ d'inconnue $x$. Si pour tout $y$ il existe une \emph{unique} solution $x$, alors $f$ est bijective (et on obtient au passage l'expression de $f^{-1}$).
\end{enumerate}
\end{methode}

\begin{exoresolu}[Bijectivité par résolution directe du système $f(x) = y$]
Soit $g : \mathbb{R}^2 \to \mathbb{R}^2$ définie par $g(x,y) = (2x + y,\; x - y)$. Montrer que $g$ est un automorphisme de $\mathbb{R}^2$ et déterminer $g^{-1}$.
\tcblower
$g$ est linéaire (somme d'applications coordonnées linéaires — vérification immédiate laissée au lecteur). Soit $(a,b) \in \mathbb{R}^2$ un vecteur quelconque. Résolvons $g(x,y) = (a,b)$ :
\[ \begin{cases} 2x + y = a \\ x - y = b \end{cases} \iff \begin{cases} 2x + y = a \\ 3x = a + b \end{cases} \quad (L_2 \leftarrow L_1 + L_2) \]
d'où $x = \dfrac{a+b}{3}$, puis $y = x - b = \dfrac{a + b - 3b}{3} = \dfrac{a - 2b}{3}$.
Pour tout $(a,b) \in \mathbb{R}^2$, le système admet une \textbf{unique} solution :
\[ (x,y) = \left(\dfrac{a+b}{3},\; \dfrac{a - 2b}{3}\right). \]
Donc $g$ est bijective : c'est un \cle{automorphisme} de $\mathbb{R}^2$, et sa réciproque est
\[ g^{-1}(a,b) = \left(\dfrac{a+b}{3},\; \dfrac{a-2b}{3}\right). \]
\emph{Vérification par le noyau (plus rapide) :} $g(x,y) = (0,0)$ donne $3x = 0$ puis $y = x = 0$, donc $\ker g = \{(0,0)\}$ : $g$ est injective, et comme c'est un endomorphisme de $\mathbb{R}^2$ (dimension finie), elle est bijective.
\end{exoresolu}

\begin{exemplebox}[Un exemple dans $\mathbb{R}^3$]
Soit $h : \mathbb{R}^3 \to \mathbb{R}^3$ définie par $h(x,y,z) = (x + z,\; y + z,\; x + y)$. Noyau :
\[ h(x,y,z) = (0,0,0) \iff \begin{cases} x + z = 0 \\ y + z = 0 \\ x + y = 0 \end{cases} \iff \begin{cases} x = -z \\ y = -z \\ -2z = 0 \end{cases} \iff x = y = z = 0. \]
Donc $\ker h = \{(0,0,0)\}$ : $h$ est injective, et comme $h$ est un endomorphisme de $\mathbb{R}^3$, c'est un \textbf{automorphisme}. Son image est alors $\mathbb{R}^3$ tout entier : les trois vecteurs $h(e_1) = (1,0,1)$, $h(e_2) = (0,1,1)$, $h(e_3) = (1,1,0)$ forment une base de $\mathbb{R}^3$.
\end{exemplebox}

\begin{attention}[Pièges de rédaction au Probatoire]
\begin{itemize}
\item Ne jamais écrire $\ker f = \varnothing$ : le noyau contient toujours $0_E$.
\item $\ker f$ est un sous-espace de $E$ (l'ensemble de \emph{départ}), $\mathrm{Im}\,f$ un sous-espace de $F$ (l'ensemble d'\emph{arrivée}) : ne pas les inverser !
\item L'équivalence « injective $\iff$ surjective » n'est valable que si $\dim E = \dim F$ (en particulier pour un endomorphisme). Pour $f : \mathbb{R}^3 \to \mathbb{R}^2$, elle ne s'applique pas.
\item Dire « $\ker f = \{0_E\}$ » ne suffit pas à conclure à la bijectivité si les dimensions diffèrent.
\end{itemize}
\end{attention}

\begin{savaistu}[Le théorème du rang : la balance parfaite]
Dans tous les exemples de cette leçon, tu as pu constater que $\dim \ker f + \dim \mathrm{Im}\,f = \dim E$. Ce résultat général s'appelle le \cle{théorème du rang} :
\[ \dim E = \dim \ker f + \dim \mathrm{Im}\,f. \]
Il dit que la dimension de l'espace de départ se répartit exactement entre ce qui est « écrasé » (noyau) et ce qui est « atteint » (image). Ce théorème sera étudié en Terminale ; il n'est pas exigible au Probatoire de Première, mais il explique pourquoi, en dimension finie égale, injectivité et surjectivité vont de pair : si $\ker f = \{0_E\}$, alors $\dim \mathrm{Im}\,f = \dim E = \dim F$, donc $\mathrm{Im}\,f = F$. La dimension de l'image, $\dim \mathrm{Im}\,f$, s'appelle le \emph{rang} de $f$ — une notion centrale en algèbre linéaire, utilisée partout, du traitement d'images aux réseaux de neurones qui font fonctionner les applications de reconnaissance vocale sur ton téléphone.
\end{savaistu}

\begin{exercice}[Pour s'entraîner]
Soit $f : \mathbb{R}^3 \to \mathbb{R}^3$ définie par $f(x,y,z) = (x - y,\; y - z,\; z - x)$.
\begin{enumerate}
\item Montrer que $f$ est un endomorphisme de $\mathbb{R}^3$.
\item Déterminer une équation caractéristique et une base de $\ker f$.
\item Déterminer une base de $\mathrm{Im}\,f$ et son équation caractéristique.
\item $f$ est-elle bijective ? Justifier.
\end{enumerate}
\end{exercice}

\begin{corrige}[Exercice 1]
\textbf{1.} Pour $u = (x,y,z)$, $v = (x',y',z')$ et $\lambda \in \mathbb{R}$ : $f(u+v) = f(u) + f(v)$ et $f(\lambda u) = \lambda f(u)$ par distributivité des opérations coordonnée par coordonnée. Donc $f$ est linéaire, de $\mathbb{R}^3$ dans $\mathbb{R}^3$ : c'est un endomorphisme.

\textbf{2.} $f(x,y,z) = (0,0,0) \iff x = y$, $y = z$, $z = x$, c'est-à-dire $x = y = z$. Les solutions sont $(t,t,t) = t(1,1,1)$. Donc $\ker f = \mathrm{Vect}\big((1,1,1)\big)$, de base $\big((1,1,1)\big)$, dimension $1$.

\textbf{3.} $f(e_1) = (1,0,-1)$, $f(e_2) = (-1,1,0)$, $f(e_3) = (0,-1,1)$. Or $f(e_3) = -f(e_1) - f(e_2)$ : on élimine $f(e_3)$. Les vecteurs $(1,0,-1)$ et $(-1,1,0)$ ne sont pas colinéaires, donc forment une famille libre. Ainsi $\mathrm{Im}\,f = \mathrm{Vect}\big((1,0,-1), (-1,1,0)\big)$, plan vectoriel de dimension $2$. Équation caractéristique : $(a,b,c) \in \mathrm{Im}\,f \iff$ le système $x - y = a$, $y - z = b$, $z - x = c$ a une solution ; en additionnant les trois lignes : $a + b + c = 0$.

\textbf{4.} $\ker f \neq \{(0,0,0)\}$, donc $f$ n'est pas injective ; $\mathrm{Im}\,f \neq \mathbb{R}^3$, donc $f$ n'est pas surjective. Conclusion : $f$ n'est \textbf{pas bijective}. (On vérifie : $1 + 2 = 3 = \dim \mathbb{R}^3$.)
\end{corrige}

\newpage

\coursec{Activité d'intégration}

\coursec{Leçon 16 : Espaces vectoriels sur $\mathbb{R}$ — Configurations et transformations élémentaires du plan}

\courssub{Objectifs de l'activité}

À l'issue de cette activité d'intégration, tu seras capable de :
\begin{itemize}
\item justifier qu'un ensemble concret de déplacements, muni de l'addition et de la multiplication par un réel, possède une structure d'\cle{espace vectoriel sur} $\mathbb{R}$ ;
\item reconnaître et exploiter une \cle{base} de $\mathbb{R}^2$ pour décomposer un vecteur en \cle{combinaison linéaire} ;
\item démontrer qu'une application donnée par son écriture analytique est un \cle{endomorphisme} de $\mathbb{R}^2$ ;
\item déterminer le \cle{noyau} d'une application linéaire et en déduire sa \cle{bijectivité} ;
\item déterminer l'\cle{image} d'une application linéaire et résoudre un problème inverse (retrouver l'itinéraire initial à partir de l'itinéraire transformé).
\end{itemize}

\courssub{Situation}

La ville de Ngaoundéré, chef-lieu de la région de l'Adamaoua, est traversée chaque jour par des centaines de motos-taxis (les « bendskins ») qui relient les quartiers : le marché central, la gare ferroviaire, le quartier Dang, le campus universitaire de Ngaoundéré, le rond-point Bamyanga. Pour organiser ses itinéraires, la jeune compagnie \emph{Adamaoua Express} modélise chaque déplacement en taxi-moto par un couple de réels $(x\,;\,y)$, où :
\begin{itemize}
\item $x$ désigne la composante du déplacement selon l'axe \emph{Est--Ouest}, en kilomètres (positive vers l'Est, négative vers l'Ouest) ;
\item $y$ désigne la composante du déplacement selon l'axe \emph{Nord--Sud}, en kilomètres (positive vers le Nord, négative vers le Sud).
\end{itemize}
Ainsi, le trajet « marché central $\to$ gare » est codé par le vecteur $\vec{v}=(4\,;\,-2)$ : \SI{4}{km} vers l'Est puis \SI{2}{km} vers le Sud. On note $E$ l'ensemble de tous les déplacements possibles, muni des deux opérations usuelles :
\[
(x\,;\,y)+(x'\,;\,y')=(x+x'\,;\,y+y') \qquad \text{et} \qquad \lambda\cdot(x\,;\,y)=(\lambda x\,;\,\lambda y)\quad (\lambda\in\mathbb{R}).
\]
L'addition de deux déplacements correspond à l'enchaînement des deux trajets ; la multiplication par un réel $\lambda$ correspond à la répétition ($\lambda>0$) ou à l'inversion ($\lambda<0$) d'un trajet.

Pour réorganiser son réseau (regroupement des courses, optimisation des itinéraires), la compagnie applique à chaque trajet $(x\,;\,y)$ la transformation définie par :
\[
f(x\,;\,y)=(2x+y\,;\,x-y).
\]
Le couple obtenu code le nouvel itinéraire effectivement parcouru par la moto après réorganisation. La direction de la compagnie se pose plusieurs questions :
\begin{enumerate}
\item L'ensemble $E$ des déplacements possède-t-il une structure qui permette de faire des calculs (sommes, combinaisons) ?
\item Les deux trajets élémentaires $\vec{e}_1=(1\,;\,0)$ (\SI{1}{km} vers l'Est) et $\vec{e}_2=(0\,;\,1)$ (\SI{1}{km} vers le Nord) suffisent-ils à engendrer tous les déplacements ? Comment décomposer le trajet « marché central $\to$ gare » ?
\item La transformation $f$ respecte-t-elle les opérations du réseau (enchaînements et répétitions de trajets) ?
\item Deux itinéraires initiaux distincts peuvent-ils donner le même itinéraire transformé ? Autrement dit, la réorganisation est-elle réversible ?
\item Un chauffeur a parcouru l'itinéraire transformé $(5\,;\,1)$. Quel était l'itinéraire initialement commandé par le client ?
\end{enumerate}

\courssub{Consignes}

Travail en groupes de quatre élèves. Chaque groupe rédige ses réponses sur une feuille, puis un rapporteur présente les résultats au tableau.

\begin{enumerate}
\item \textbf{Structure de l'ensemble des déplacements.}
\begin{enumerate}
\item Vérifier que l'addition définie sur $E$ est une loi de composition interne commutative et associative, qu'elle admet un élément neutre (le préciser) et que tout déplacement admet un opposé (l'interpréter concrètement).
\item Vérifier que la multiplication par un réel est une loi de composition externe sur $E$ et que, pour tous réels $\lambda,\mu$ et tous déplacements $\vec{u},\vec{v}$ :
\[
\lambda\cdot(\vec{u}+\vec{v})=\lambda\cdot\vec{u}+\lambda\cdot\vec{v},\quad (\lambda+\mu)\cdot\vec{u}=\lambda\cdot\vec{u}+\mu\cdot\vec{u},\quad \lambda\cdot(\mu\cdot\vec{u})=(\lambda\mu)\cdot\vec{u},\quad 1\cdot\vec{u}=\vec{u}.
\]
Conclure sur la structure de $(E,+,\cdot)$.
\end{enumerate}
\item \textbf{Base du réseau.}
\begin{enumerate}
\item Montrer que la famille $(\vec{e}_1,\vec{e}_2)$ est génératrice de $E$, puis qu'elle est libre. Conclure.
\item Décomposer le trajet $\vec{v}=(4\,;\,-2)$ « marché central $\to$ gare » sur $(\vec{e}_1,\vec{e}_2)$ et interpréter le résultat.
\item Quelle est la dimension de $E$ ?
\end{enumerate}
\item \textbf{Étude de la transformation $f$.}
\begin{enumerate}
\item Montrer que pour tous $\vec{u},\vec{v}\in E$ et tout réel $\lambda$ : $f(\vec{u}+\vec{v})=f(\vec{u})+f(\vec{v})$ et $f(\lambda\cdot\vec{u})=\lambda f(\vec{u})$. Conclure que $f$ est un endomorphisme de $\mathbb{R}^2$.
\item Calculer $f(\vec{e}_1)$ et $f(\vec{e}_2)$.
\end{enumerate}
\item \textbf{Noyau et bijectivité.}
\begin{enumerate}
\item Déterminer le noyau $\ker f$, c'est-à-dire l'ensemble des itinéraires transformés en déplacement nul.
\item En déduire que $f$ est injective, puis, en justifiant, que $f$ est bijective (automorphisme de $\mathbb{R}^2$). Interpréter : « tout itinéraire transformé provient d'un unique itinéraire initial ».
\end{enumerate}
\item \textbf{Image et problème inverse.}
\begin{enumerate}
\item Montrer que $\mathrm{Im}\,f=\mathbb{R}^2$ (on pourra utiliser les questions 3.b et 4).
\item Déterminer l'itinéraire initial $(x\,;\,y)$ dont l'itinéraire transformé est $(5\,;\,1)$, en résolvant un système. Vérifier le résultat.
\end{enumerate}
\end{enumerate}

\courssub{Ressources à mobiliser}

\begin{aretenir}[Boîte à outils de l'activité]
\begin{itemize}
\item \textbf{Espace vectoriel sur $\mathbb{R}$} : ensemble muni d'une loi interne $+$ (groupe commutatif) et d'une loi externe $\cdot$ vérifiant les quatre propriétés de compatibilité (distributivités, associativité mixte, $1\cdot\vec{u}=\vec{u}$).
\item \textbf{Famille génératrice} : tout vecteur de $E$ s'écrit comme combinaison linéaire des vecteurs de la famille. \textbf{Famille libre} : $\lambda_1\vec{u}_1+\cdots+\lambda_p\vec{u}_p=\vec{0}\Rightarrow \lambda_1=\cdots=\lambda_p=0$. \textbf{Base} : famille à la fois libre et génératrice ; toutes les bases ont le même nombre de vecteurs : la \cle{dimension}.
\item \textbf{Application linéaire} : $f(\vec{u}+\vec{v})=f(\vec{u})+f(\vec{v})$ et $f(\lambda\vec{u})=\lambda f(\vec{u})$. Un \cle{endomorphisme} est une application linéaire de $E$ dans lui-même ; un endomorphisme bijectif est un \cle{automorphisme}.
\item \textbf{Noyau} : $\ker f=\{\vec{u}\in E\mid f(\vec{u})=\vec{0}\}$. $f$ est injective si et seulement si $\ker f=\{\vec{0}\}$.
\item \textbf{Image} : $\mathrm{Im}\,f=\{f(\vec{u})\mid \vec{u}\in E\}$ ; c'est le sous-espace vectoriel engendré par les images des vecteurs d'une base. En dimension finie : $\dim\ker f+\dim\mathrm{Im}\,f=\dim E$ (théorème du rang).
\end{itemize}
\end{aretenir}

\courssub{Démarche et corrigé commenté}

\begin{exoresolu}[Le réseau de transport de Ngaoundéré — corrigé complet]
\textbf{1.} Structure de l'ensemble des déplacements. \textbf{2.} Base et décomposition du trajet $(4\,;\,-2)$. \textbf{3.} $f$ est un endomorphisme. \textbf{4.} Noyau et bijectivité. \textbf{5.} Image et itinéraire initial de l'itinéraire transformé $(5\,;\,1)$.
\tcblower
\textbf{1.a) Loi interne.} Soient $\vec{u}=(x\,;\,y)$, $\vec{v}=(x'\,;\,y')$, $\vec{w}=(x''\,;\,y'')$ des éléments de $E$.
\begin{itemize}
\item $\vec{u}+\vec{v}=(x+x'\,;\,y+y')\in E$ car $x+x'$ et $y+y'$ sont des réels : $+$ est une loi de composition \cle{interne}.
\item Commutativité : $\vec{u}+\vec{v}=(x+x'\,;\,y+y')=(x'+x\,;\,y'+y)=\vec{v}+\vec{u}$, car l'addition des réels est commutative.
\item Associativité : $(\vec{u}+\vec{v})+\vec{w}=\big((x+x')+x''\,;\,(y+y')+y''\big)=\big(x+(x'+x'')\,;\,y+(y'+y'')\big)=\vec{u}+(\vec{v}+\vec{w})$.
\item Élément neutre : $\vec{0}=(0\,;\,0)$, le « déplacement nul » (la moto reste sur place) : $\vec{u}+\vec{0}=\vec{u}$.
\item Opposé : $-\vec{u}=(-x\,;\,-y)$, le trajet retour (même distance, sens inverse) : $\vec{u}+(-\vec{u})=\vec{0}$.
\end{itemize}
Donc $(E,+)$ est un \textbf{groupe commutatif}.

\textbf{1.b) Loi externe.} Pour $\lambda\in\mathbb{R}$, $\lambda\cdot\vec{u}=(\lambda x\,;\,\lambda y)\in E$ : la loi $\cdot:\mathbb{R}\times E\to E$ est bien une loi de composition \cle{externe}. Vérifions la première distributivité (les autres se traitent de même, coordonnée par coordonnée) :
\begin{align*}
\lambda\cdot(\vec{u}+\vec{v})&=\lambda\cdot(x+x'\,;\,y+y')=\big(\lambda(x+x')\,;\,\lambda(y+y')\big)\\
&=(\lambda x+\lambda x'\,;\,\lambda y+\lambda y')=(\lambda x\,;\,\lambda y)+(\lambda x'\,;\,\lambda y')=\lambda\cdot\vec{u}+\lambda\cdot\vec{v}.
\end{align*}
De même : $(\lambda+\mu)\cdot\vec{u}=\big((\lambda+\mu)x\,;\,(\lambda+\mu)y\big)=\lambda\cdot\vec{u}+\mu\cdot\vec{u}$ ; $\lambda\cdot(\mu\cdot\vec{u})=(\lambda\mu x\,;\,\lambda\mu y)=(\lambda\mu)\cdot\vec{u}$ ; $1\cdot\vec{u}=(x\,;\,y)=\vec{u}$.

\textbf{Conclusion :} $(E,+,\cdot)$ est un \textbf{espace vectoriel sur $\mathbb{R}$} : c'est l'espace vectoriel $\mathbb{R}^2$. Tous les calculs sur les itinéraires sont donc légitimes.

\textbf{2.a) $(\vec{e}_1,\vec{e}_2)$ est une base.} \emph{Famille génératrice :} pour tout $\vec{u}=(x\,;\,y)\in E$ :
\[
\vec{u}=(x\,;\,y)=(x\,;\,0)+(0\,;\,y)=x(1\,;\,0)+y(0\,;\,1)=x\,\vec{e}_1+y\,\vec{e}_2.
\]
Tout déplacement est combinaison linéaire de $\vec{e}_1$ et $\vec{e}_2$. \emph{Famille libre :} si $\lambda\vec{e}_1+\mu\vec{e}_2=\vec{0}$, alors $(\lambda\,;\,\mu)=(0\,;\,0)$, donc $\lambda=\mu=0$. La famille $(\vec{e}_1,\vec{e}_2)$, libre et génératrice, est une \textbf{base} de $E$ (la \emph{base canonique} de $\mathbb{R}^2$).

\textbf{2.b) Décomposition.} $\vec{v}=(4\,;\,-2)=4\,\vec{e}_1-2\,\vec{e}_2$ : le trajet « marché central $\to$ gare » s'obtient en enchaînant 4 fois le trajet élémentaire Est et 2 fois le trajet élémentaire Sud (l'opposé de $\vec{e}_2$).

\begin{popfigure}
\begin{tikzpicture}[>=Stealth, scale=0.7]
\draw[popline, ->] (-1,0) -- (5,0) node[right] {$x$ (Est)};
\draw[popline, ->] (0,-3) -- (0,2) node[above] {$y$ (Nord)};
\draw[popcurve, thick, popblue, ->] (0,0) -- (1,0) node[midway, below] {$\vec{e}_1$};
\draw[popcurve, thick, popgreen, ->] (0,0) -- (0,1) node[midway, left] {$\vec{e}_2$};
\draw[popcurve, thick, poporange, ->] (0,0) -- (4,-2) node[midway, above right] {$\vec{v}=(4;-2)$};
\draw[dashed, popmuted] (4,0) -- (4,-2);
\draw[dashed, popmuted] (0,-2) -- (4,-2);
\node[popmuted] at (2,0.3) {$4\vec{e}_1$};
\node[popmuted] at (4.3,-1) {$-2\vec{e}_2$};
\end{tikzpicture}
\legende{Décomposition de $\vec{v}=(4;-2)$ dans la base canonique $(\vec{e}_1,\vec{e}_2)$.}
\end{popfigure}

\textbf{2.c) Dimension.} $E$ admet une base de 2 vecteurs, donc $\dim E=2$.

\textbf{3.a) Linéarité de $f$.} Soient $\vec{u}=(x\,;\,y)$, $\vec{v}=(x'\,;\,y')$ et $\lambda\in\mathbb{R}$.
\begin{align*}
f(\vec{u}+\vec{v})&=f(x+x'\,;\,y+y')=\big(2(x+x')+(y+y')\,;\,(x+x')-(y+y')\big)\\
&=(2x+y\,;\,x-y)+(2x'+y'\,;\,x'-y')=f(\vec{u})+f(\vec{v}) ;\\
f(\lambda\cdot\vec{u})&=f(\lambda x\,;\,\lambda y)=(2\lambda x+\lambda y\,;\,\lambda x-\lambda y)=\lambda\,(2x+y\,;\,x-y)=\lambda f(\vec{u}).
\end{align*}
$f$ est une application linéaire de $E$ dans $E$ : c'est un \textbf{endomorphisme} de $\mathbb{R}^2$. Concrètement, réorganiser un enchaînement de trajets revient à enchaîner les trajets réorganisés.

\textbf{3.b)} $f(\vec{e}_1)=f(1\,;\,0)=(2\,;\,1)$ et $f(\vec{e}_2)=f(0\,;\,1)=(1\,;\,-1)$.

\begin{popfigure}
\begin{tikzpicture}[>=Stealth, scale=0.7]
\draw[popline, ->] (-1,-1) -- (3,2) node[right] {$x$};
\draw[popline, ->] (-1,-1) -- (2,3) node[above] {$y$};
% Original basis
\draw[popcurve, thick, popblue, ->] (0,0) -- (1,0) node[midway, below] {$\vec{e}_1$};
\draw[popcurve, thick, popgreen, ->] (0,0) -- (0,1) node[midway, left] {$\vec{e}_2$};
% Transformed vectors
\draw[popcurve, thick, poporange, ->, dashed] (0,0) -- (2,1) node[midway, right] {$f(\vec{e}_1)=(2;1)$};
\draw[popcurve, thick, poppurple, ->, dashed] (0,0) -- (1,-1) node[midway, below right] {$f(\vec{e}_2)=(1;-1)$};
\end{tikzpicture}
\legende{Images des vecteurs de la base canonique par l'endomorphisme $f$.}
\end{popfigure}

\textbf{4.a) Noyau.} $\vec{u}=(x\,;\,y)\in\ker f\iff f(x\,;\,y)=(0\,;\,0)$, c'est-à-dire :
\[
\begin{cases} 2x+y=0 \\ x-y=0 \end{cases}
\iff \begin{cases} y=-2x \\ y=x \end{cases}
\iff x=-2x \iff 3x=0 \iff x=0 \text{ et } y=0.
\]
Donc $\ker f=\{(0\,;\,0)\}=\{\vec{0}\}$ : le seul itinéraire transformé en déplacement nul est le déplacement nul. Le noyau est de dimension 0 (il n'admet pas de base, ou sa base est la famille vide).

\textbf{4.b) Bijectivité.} Comme $\ker f=\{\vec{0}\}$, $f$ est \textbf{injective}. Par le théorème du rang : $\dim\mathrm{Im}\,f=\dim E-\dim\ker f=2-0=2=\dim\mathbb{R}^2$. Or $\mathrm{Im}\,f$ est un sous-espace vectoriel de $\mathbb{R}^2$ de même dimension que $\mathbb{R}^2$, donc $\mathrm{Im}\,f=\mathbb{R}^2$ : $f$ est \textbf{surjective}. Étant injective et surjective, $f$ est \textbf{bijective} : c'est un \textbf{automorphisme} de $\mathbb{R}^2$. Interprétation : tout itinéraire transformé provient d'un \emph{unique} itinéraire initial ; la réorganisation du réseau est parfaitement réversible.

\textbf{5.a) Image.} D'après 4.b, $\mathrm{Im}\,f=\mathbb{R}^2$. On peut le retrouver directement : $\mathrm{Im}\,f$ est engendré par $f(\vec{e}_1)=(2\,;\,1)$ et $f(\vec{e}_2)=(1\,;\,-1)$. Or ces deux vecteurs ne sont pas colinéaires ($2\times(-1)-1\times 1=-3\neq 0$) : ils forment une famille libre de 2 vecteurs, donc une base de $\mathbb{R}^2$, et engendrent bien tout $\mathbb{R}^2$.

\textbf{5.b) Problème inverse.} On cherche $(x\,;\,y)$ tel que $f(x\,;\,y)=(5\,;\,1)$ :
\[
\begin{cases} 2x+y=5 \\ x-y=1 \end{cases}
\underset{\text{somme des deux lignes}}{\Longrightarrow} 3x=6 \iff x=2, \quad \text{puis } y=x-1=1.
\]
L'itinéraire initial est $(2\,;\,1)$ : \SI{2}{km} vers l'Est puis \SI{1}{km} vers le Nord. \emph{Vérification :} $f(2\,;\,1)=(2\times 2+1\,;\,2-1)=(5\,;\,1)$. Le résultat est conforme.
\end{exoresolu}

\courssub{Bilan de l'activité}

Cette activité a permis de mettre en évidence les points essentiels de la leçon :
\begin{itemize}
\item la structure d'\cle{espace vectoriel} n'est pas une abstraction gratuite : c'est elle qui garantit que les calculs sur les déplacements (enchaînements, répétitions, trajets retours) sont cohérents ;
\item une \cle{base} joue le rôle d'un « jeu minimal de trajets élémentaires » : avec seulement deux trajets de référence (Est et Nord), on engendre tous les déplacements du plan, d'où la \cle{dimension} 2 ;
\item la \cle{linéarité} de la transformation $f$ se vérifie directement sur l'écriture analytique et signifie concrètement que la réorganisation du réseau respecte les enchaînements et les répétitions de trajets ;
\item le \cle{noyau} répond à une question d'unicité : $\ker f=\{\vec{0}\}$ assure que deux itinéraires distincts ne donnent jamais le même itinéraire transformé ;
\item le théorème du rang relie noyau et \cle{image} : en dimension finie, l'injectivité d'un endomorphisme entraîne sa bijectivité, ce qui rend le problème inverse (retrouver l'itinéraire initial) toujours soluble, par simple résolution d'un système linéaire.
\end{itemize}
Tu as ainsi mobilisé, dans une même situation authentique, la quasi-totalité des savoir-faire du chapitre : structure d'espace vectoriel, base et dimension, application linéaire, noyau, image et bijectivité. C'est exactement la démarche attendue au Probatoire dans les problèmes d'algèbre linéaire.

\begin{savaistu}[Pour aller plus loin]
La transformation $f$ s'écrit matriciellement $\begin{pmatrix} 2 & 1 \\ 1 & -1 \end{pmatrix}\begin{pmatrix} x \\ y \end{pmatrix}$. Son déterminant vaut $2\times(-1)-1\times 1=-3\neq 0$, ce qui confirme d'un coup d'œil sa bijectivité. C'est le même type de calcul qui permet, en géographie et en cartographie numérique (par exemple pour les systèmes d'information géographique utilisés dans la planification urbaine de Yaoundé ou de Douala), de passer d'un système de coordonnées à un autre sans perte d'information.
\end{savaistu}

\newpage

\coursec{Méthodes et exercices résolus}

\coursec{Espaces vectoriels sur $\mathbb{R}$}

\courssub{Lois de composition externe}

\begin{definition}[Loi de composition externe]
Soit $E$ un ensemble non vide. Une \cle{loi de composition externe} sur $E$ est une application de $\mathbb{R}\times E$ dans $E$. On note souvent le couple $(\lambda, u)$ sous la forme $\lambda \cdot u$ ou $\lambda u$ (multiplication d'un vecteur par un scalaire).
\end{definition}

\begin{methode}[Reconnaître et justifier une loi de composition externe]
Pour montrer qu'une loi $\star$ est une loi de composition externe sur un ensemble $E$ :
\begin{enumerate}
\item \textbf{Domaine de départ :} Vérifier que la loi fait intervenir un \emph{réel} $\lambda$ (le scalaire) et un élément $u\in E$. Elle est bien définie sur $\mathbb{R}\times E$.
\item \textbf{Stabilité (arrivée dans $E$) :} Vérifier que pour tout $\lambda\in\mathbb{R}$ et tout $u\in E$, le résultat $\lambda\star u$ appartient encore à $E$.
\item \textbf{Conclusion :} Écrire explicitement : « Pour tout $\lambda\in\mathbb{R}$ et tout $u\in E$, $\lambda\star u = \ldots \in E$, donc $\star$ est une application de $\mathbb{R}\times E$ dans $E$, c'est une loi de composition externe. »
\end{enumerate}
\textbf{Exemples de référence :}
\begin{itemize}
\item Multiplication d'un vecteur du plan par un réel.
\item Sur $\mathbb{R}^2$ : $\lambda(x,y)=(\lambda x,\lambda y)$.
\item Sur $\mathbb{R}^3$ : $\lambda(x,y,z)=(\lambda x,\lambda y,\lambda z)$.
\end{itemize}
\end{methode}

\begin{exoresolu}[Loi externe sur $\mathbb{R}^2$]
On munit $\mathbb{R}^2$ de l'addition usuelle et de la loi $\star$ définie par $\lambda\star(x,y)=(\lambda x,\lambda y)$ pour tout $\lambda\in\mathbb{R}$ et tout $(x,y)\in\mathbb{R}^2$. Montrer que $\star$ est une loi de composition externe sur $\mathbb{R}^2$.
\tcblower
\textbf{Étape 1 : domaine de départ.} La loi $\star$ associe à un réel $\lambda$ (scalaire) et un couple $(x,y)\in\mathbb{R}^2$ le couple $(\lambda x,\lambda y)$. Elle est bien définie sur $\mathbb{R}\times\mathbb{R}^2$.

\textbf{Étape 2 : stabilité.} Soient $\lambda\in\mathbb{R}$ et $(x,y)\in\mathbb{R}^2$ quelconques. Alors
\[ \lambda\star(x,y)=(\lambda x,\lambda y). \]
Comme $x,y,\lambda\in\mathbb{R}$, les produits $\lambda x$ et $\lambda y$ sont des réels, donc $(\lambda x,\lambda y)\in\mathbb{R}^2$.

\textbf{Conclusion.} La loi $\star$ est une application de $\mathbb{R}\times\mathbb{R}^2$ dans $\mathbb{R}^2$ : c'est une loi de composition externe sur $\mathbb{R}^2$.
\end{exoresolu}

\courssub{Espaces vectoriels sur $\mathbb{R}$}

\begin{definition}[Espace vectoriel sur $\mathbb{R}$]
Un \cle{espace vectoriel sur $\mathbb{R}$} est un triplet $(E,+,\cdot)$ où :
\begin{itemize}
\item $E$ est un ensemble non vide (les \emph{vecteurs}),
\item $+$ est une loi de composition interne sur $E$ (l'\emph{addition}),
\item $\cdot$ est une loi de composition externe sur $E$ (la \emph{multiplication par un scalaire}),
\end{itemize}
telles que les huit axiomes suivants soient satisfaits pour tous $u,v,w\in E$ et $\lambda,\mu\in\mathbb{R}$ :
\begin{enumerate}
\item $(E,+)$ est un groupe commutatif : associativité, commutativité, existence d'un élément neutre $0_E$ (vecteur nul), existence d'un opposé $-u$.
\item $\lambda(u+v)=\lambda u+\lambda v$ (distributivité par rapport à l'addition dans $E$).
\item $(\lambda+\mu)u=\lambda u+\mu u$ (distributivité par rapport à l'addition dans $\mathbb{R}$).
\item $\lambda(\mu u)=(\lambda\mu)u$ (associativité de la multiplication externe).
\item $1\cdot u=u$ (élément neutre de $\mathbb{R}$).
\end{enumerate}
\end{definition}

\begin{propriete}[Caractérisation des sous-espaces vectoriels]
Soit $E$ un espace vectoriel sur $\mathbb{R}$ et $F\subset E$, $F\neq\varnothing$. Alors $F$ est un \cle{sous-espace vectoriel} de $E$ (donc un espace vectoriel en lui-même) \textbf{si et seulement si} :
\begin{enumerate}
\item $F$ est stable par l'addition : $\forall u,v\in F,\; u+v\in F$.
\item $F$ est stable par la multiplication par un réel : $\forall \lambda\in\mathbb{R},\; \forall u\in F,\; \lambda u\in F$.
\end{enumerate}
\textbf{Remarque :} La condition $F\neq\varnothing$ est équivalente à $0_E\in F$ (car si $u\in F$, alors $0\cdot u=0_E\in F$ par stabilité). \emph{Au Probatoire, on vérifie systématiquement $0_E\in F$ en premier.}
\end{propriete}

\begin{methode}[Prouver qu'un ensemble est un espace vectoriel]
Deux stratégies selon l'énoncé :
\begin{enumerate}
\item \textbf{Stratégie « Sous-espace » (privilégiée au Probatoire) :} L'ensemble $F$ est inclus dans un espace vectoriel connu $E$ ($\mathbb{R}$, $\mathbb{R}^2$, $\mathbb{R}^3$, vecteurs du plan, vecteurs de l'espace). Vérifier :
      \begin{enumerate}
      \item $F\neq\varnothing$ (souvent en montrant $0_E\in F$).
      \item Stabilité par l'addition : $\forall u,v\in F,\; u+v\in F$.
      \item Stabilité par la multiplication par un réel : $\forall \lambda\in\mathbb{R},\; \forall u\in F,\; \lambda u\in F$.
      \end{enumerate}
\item \textbf{Stratégie « Axiomes directs » :} Si $F$ n'est pas inclus dans un espace connu, vérifier les 8 axiomes de la définition.
\end{enumerate}
\textbf{Réflexe :} Cherche toujours d'abord l'espace « ambiant » connu ($\mathbb{R}^2$, $\mathbb{R}^3$, etc.).
\end{methode}

\begin{exoresolu}[Un plan de $\mathbb{R}^3$ est un espace vectoriel]
Soit $E=\{(x,y,z)\in\mathbb{R}^3 \mid x+y+z=0\}$. Montrer que $E$ est un espace vectoriel sur $\mathbb{R}$.
\tcblower
On montre que $E$ est un sous-espace vectoriel de $\mathbb{R}^3$ (espace vectoriel connu).

\textbf{Étape 1 : $E$ est non vide.} Le triplet $(0,0,0)$ vérifie $0+0+0=0$, donc $(0,0,0)\in E$. (C'est le vecteur nul de $\mathbb{R}^3$).

\textbf{Étape 2 : stabilité par addition.} Soient $u=(x,y,z)\in E$ et $u'=(x',y',z')\in E$. Alors $u+u'=(x+x',\,y+y',\,z+z')$ et
\[ (x+x')+(y+y')+(z+z')=(x+y+z)+(x'+y'+z')=0+0=0, \]
donc $u+u'\in E$.

\textbf{Étape 3 : stabilité par multiplication par un réel.} Soient $\lambda\in\mathbb{R}$ et $u=(x,y,z)\in E$. Alors $\lambda u=(\lambda x,\lambda y,\lambda z)$ et
\[ \lambda x+\lambda y+\lambda z=\lambda(x+y+z)=\lambda\times 0=0, \]
donc $\lambda u\in E$.

\textbf{Conclusion.} $E$ est un sous-espace vectoriel de $\mathbb{R}^3$, donc $E$ est un espace vectoriel sur $\mathbb{R}$.
\end{exoresolu}

\courssub{Familles génératrices, familles libres, familles liées}

\begin{definition}[Combinaison linéaire, famille génératrice, famille libre, famille liée]
Soit $(u_1,\ldots,u_p)$ une famille finie de vecteurs d'un espace vectoriel $E$.
\begin{itemize}
\item Une \cle{combinaison linéaire} de ces vecteurs est un vecteur de la forme $\lambda_1 u_1+\cdots+\lambda_p u_p$ avec $\lambda_i\in\mathbb{R}$.
\item La famille est \cle{génératrice} de $E$ (ou \cle{engendre} $E$) si \emph{tout} vecteur $u\in E$ s'écrit comme combinaison linéaire de $u_1,\ldots,u_p$ : $E=\operatorname{Vect}(u_1,\ldots,u_p)$.
\item La famille est \cle{libre} si la seule combinaison linéaire donnant le vecteur nul est la combinaison triviale : $\lambda_1 u_1+\cdots+\lambda_p u_p=0_E \implies \lambda_1=\cdots=\lambda_p=0$.
\item La famille est \cle{liée} si elle n'est pas libre, c'est-à-dire s'il existe des scalaires $\lambda_1,\ldots,\lambda_p$ \emph{non tous nuls} tels que $\lambda_1 u_1+\cdots+\lambda_p u_p=0_E$. Équivalentement, au moins un vecteur de la famille est combinaison linéaire des autres.
\end{itemize}
\end{definition}

\begin{methode}[Étudier une famille finie de vecteurs]
Soit $\mathcal{F}=(u_1,\ldots,u_p)$ dans un espace vectoriel $E$.
\begin{enumerate}
\item \textbf{Génératrice :} Prendre $u\in E$ quelconque (souvent noté $(x,y)$ ou $(x,y,z)$). Poser $u=\lambda_1 u_1+\cdots+\lambda_p u_p$. Montrer que le système obtenu admet \emph{au moins une solution} pour \emph{tout} $u$ (système compatible quel que soit le second membre).
\item \textbf{Libre :} Poser $\lambda_1 u_1+\cdots+\lambda_p u_p=0_E$. Montrer que le système homogène obtenu n'admet que la \emph{solution nulle} $\lambda_1=\cdots=\lambda_p=0$.
\item \textbf{Liée :} Soit exhiber des scalaires non tous nuls annulant la combinaison, soit montrer qu'un vecteur est combinaison des autres (ex : deux vecteurs non nuls sont liés $\iff$ ils sont colinéaires).
\end{enumerate}
\textbf{Réflexes utiles :}
\begin{itemize}
\item Une famille contenant le vecteur nul $0_E$ est \emph{toujours liée}.
\item Deux vecteurs non nuls sont liés $\iff$ ils sont colinéaires.
\item Si $\dim E=n$ (connue), une famille de $n$ vecteurs est une base $\iff$ elle est libre $\iff$ elle est génératrice.
\end{itemize}
\end{methode}

\begin{exoresolu}[Famille libre et génératrice dans $\mathbb{R}^2$]
Dans $\mathbb{R}^2$, on considère $u=(1,2)$ et $v=(3,-1)$.
\begin{enumerate}
\item Montrer que la famille $(u,v)$ est libre.
\item Montrer que $(u,v)$ est génératrice de $\mathbb{R}^2$.
\end{enumerate}
\tcblower
\textbf{1. Liberté.} Soient $\lambda,\mu\in\mathbb{R}$ tels que $\lambda u+\mu v=0_{\mathbb{R}^2}$. Alors
\[ \lambda(1,2)+\mu(3,-1)=(0,0)\iff\begin{cases}\lambda+3\mu=0\\2\lambda-\mu=0\end{cases} \]
De la deuxième équation, $\mu=2\lambda$. En reportant dans la première : $\lambda+6\lambda=0$, soit $7\lambda=0$, donc $\lambda=0$ puis $\mu=0$. La seule solution est $(\lambda,\mu)=(0,0)$ : la famille $(u,v)$ est \cle{libre}.

\textbf{2. Famille génératrice.} Soit $w=(a,b)$ un vecteur quelconque de $\mathbb{R}^2$. On cherche $\lambda,\mu$ tels que $\lambda u+\mu v=w$ :
\[ \begin{cases}\lambda+3\mu=a\\2\lambda-\mu=b\end{cases} \]
De la deuxième équation, $\mu=2\lambda-b$. En reportant : $\lambda+3(2\lambda-b)=a$, soit $7\lambda=a+3b$, d'où
\[ \lambda=\dfrac{a+3b}{7},\qquad \mu=\dfrac{2a-b}{7}. \]
Ces réels existent pour \emph{tout} couple $(a,b)$, donc tout vecteur de $\mathbb{R}^2$ est combinaison linéaire de $u$ et $v$ : la famille $(u,v)$ est \cle{génératrice} de $\mathbb{R}^2$.

\textbf{Conclusion.} $(u,v)$ est libre et génératrice : c'est une base de $\mathbb{R}^2$.
\end{exoresolu}

\courssub{Bases et dimension d'un espace vectoriel}

\begin{definition}[Base, dimension]
Une \cle{base} d'un espace vectoriel $E$ est une famille de vecteurs à la fois \textbf{libre} et \textbf{génératrice}.
Si $E$ admet une base finie, toutes ses bases ont le même nombre de vecteurs : ce nombre est la \cle{dimension} de $E$, notée $\dim E$.
\end{definition}

\begin{propriete}[Dimensions de référence]
\begin{itemize}
\item $\dim\mathbb{R}=1$ (base : $(1)$).
\item $\dim\mathbb{R}^2=2$ (base canonique : $((1,0),(0,1))$).
\item $\dim\mathbb{R}^3=3$ (base canonique : $((1,0,0),(0,1,0),(0,0,1))$).
\item L'ensemble des vecteurs du plan (muni de l'addition et multiplication par un réel) est de dimension $2$.
\item L'ensemble des vecteurs de l'espace est de dimension $3$.
\end{itemize}
\end{propriete}

\begin{methode}[Montrer qu'une famille est une base et déterminer la dimension]
\begin{enumerate}
\item \textbf{Preuve standard :} Montrer que la famille est \textbf{libre} \textbf{ET} \textbf{génératrice} (deux points indépendants).
\item \textbf{Raccourci (si $\dim E$ connue) :} Si $\dim E=n$, il suffit de montrer \emph{soit} que la famille est libre à $n$ vecteurs, \emph{soit} qu'elle est génératrice à $n$ vecteurs. Une seule vérification suffit (à condition de citer la dimension connue).
\item \textbf{Sous-espace défini par équation(s) :} Exprimer une inconnue en fonction des autres (paramétrisation), décomposer le vecteur générique en combinaison linéaire de vecteurs fixes $\to$ ces vecteurs forment une famille génératrice. Vérifier leur liberté $\to$ base. Le nombre de paramètres = dimension.
\end{enumerate}
\end{methode}

\begin{exoresolu}[Base et dimension d'un sous-espace de $\mathbb{R}^3$]
Soit $F=\{(x,y,z)\in\mathbb{R}^3 \mid x-2y+z=0\}$. On admet que $F$ est un sous-espace vectoriel de $\mathbb{R}^3$. Déterminer une base de $F$ et sa dimension.
\tcblower
\textbf{Étape 1 : vecteur générique et famille génératrice.} Un vecteur $(x,y,z)\in F$ vérifie $x=2y-z$. Donc
\[ (x,y,z)=(2y-z,\;y,\;z)=y(2,1,0)+z(-1,0,1). \]
Posons $u=(2,1,0)$ et $v=(-1,0,1)$. Tout vecteur de $F$ est combinaison linéaire de $u$ et $v$ : la famille $(u,v)$ est \cle{génératrice} de $F$. De plus $u\in F$ ($2-2+0=0$) et $v\in F$ ($-1-0+1=0$).

\textbf{Étape 2 : liberté.} Soient $\lambda,\mu$ tels que $\lambda u+\mu v=0$. Alors
\[ (2\lambda-\mu,\;\lambda,\;\mu)=(0,0,0). \]
La deuxième coordonnée donne $\lambda=0$, la troisième donne $\mu=0$. La famille $(u,v)$ est \cle{libre}.

\textbf{Conclusion.} $(u,v)$ est libre et génératrice de $F$ : c'est une base de $F$, et
\[ \dim F=2. \]
(C'est cohérent : $F$ est un plan vectoriel dans $\mathbb{R}^3$.)
\end{exoresolu}

\courssub{Applications linéaires}

\begin{definition}[Application linéaire, endomorphisme, isomorphisme, automorphisme]
Soient $E,F$ deux espaces vectoriels sur $\mathbb{R}$. Une application $f:E\to F$ est \cle{linéaire} si :
\[ \forall u,v\in E,\; \forall \lambda\in\mathbb{R},\quad f(u+v)=f(u)+f(v)\quad\text{et}\quad f(\lambda u)=\lambda f(u). \]
Équivalentement : $\forall u,v\in E,\; \forall \lambda,\mu\in\mathbb{R},\; f(\lambda u+\mu v)=\lambda f(u)+\mu f(v)$.
\begin{itemize}
\item Si $E=F$, $f$ est un \cle{endomorphisme} de $E$.
\item Si $f$ est linéaire et bijective, $f$ est un \cle{isomorphisme} de $E$ dans $F$.
\item Un endomorphisme bijectif est un \cle{automorphisme}.
\end{itemize}
\end{definition}

\begin{methode}[Prouver la linéarité d'une application]
Soit $f:E\to F$.
\begin{enumerate}
\item Prendre $u,v\in E$ \emph{quelconques}, calculer $f(u+v)$ et $f(u)+f(v)$, vérifier l'égalité.
\item Prendre $\lambda\in\mathbb{R}$ et $u\in E$ \emph{quelconques}, calculer $f(\lambda u)$ et $\lambda f(u)$, vérifier l'égalité.
\item \textbf{Variante en un calcul :} Vérifier directement $f(\lambda u+\mu v)=\lambda f(u)+\mu f(v)$ pour tous $\lambda,\mu\in\mathbb{R}$, $u,v\in E$.
\end{enumerate}
\textbf{Contre-exemple (non-linéarité) :} Il suffit de trouver \emph{un} couple $(u,v)$ ou $(\lambda,u)$ où l'égalité échoue. Souvent, tester $f(0_E)=0_F$ est le plus rapide : si $f(0_E)\neq 0_F$, $f$ n'est pas linéaire.
\end{methode}

\begin{exoresolu}[Linéarité et nature d'une application]
Soit $f:\mathbb{R}^2\to\mathbb{R}^2$ définie par $f(x,y)=(2x-y,\;x+3y)$.
\begin{enumerate}
\item Montrer que $f$ est une application linéaire.
\item Préciser sa nature (endomorphisme, isomorphisme, automorphisme ?).
\end{enumerate}
\tcblower
\textbf{1. Linéarité.} Soient $u=(x,y)$, $v=(x',y')$ dans $\mathbb{R}^2$ et $\lambda\in\mathbb{R}$.

\emph{Addition :} $u+v=(x+x',\,y+y')$, donc
\begin{align*}
f(u+v)&=\big(2(x+x')-(y+y'),\;(x+x')+3(y+y')\big)\\
&=(2x-y,\;x+3y)+(2x'-y',\;x'+3y')=f(u)+f(v).
\end{align*}

\emph{Multiplication par un réel :} $\lambda u=(\lambda x,\lambda y)$, donc
\[ f(\lambda u)=(2\lambda x-\lambda y,\;\lambda x+3\lambda y)=\lambda(2x-y,\;x+3y)=\lambda f(u). \]
$f$ est donc une application linéaire.

\textbf{2. Nature.} $f$ va de $\mathbb{R}^2$ dans $\mathbb{R}^2$ : c'est un \cle{endomorphisme} de $\mathbb{R}^2$. Pour la bijectivité, résolvons $f(x,y)=(a,b)$ :
\[ \begin{cases}2x-y=a\\x+3y=b\end{cases}\iff\begin{cases}y=2x-a\\x+3(2x-a)=b\end{cases}\iff 7x=3a+b,\; 7y=2b-a. \]
Ainsi $x=\dfrac{3a+b}{7},\; y=\dfrac{2b-a}{7}$. Tout $(a,b)$ a un unique antécédent : $f$ est bijective. C'est donc un isomorphisme de $\mathbb{R}^2$ sur $\mathbb{R}^2$, c'est-à-dire un \cle{automorphisme} de $\mathbb{R}^2$.
\end{exoresolu}

\courssub{Noyau et image d'une application linéaire}

\begin{definition}[Noyau, Image]
Soit $f:E\to F$ une application linéaire.
\begin{itemize}
\item Le \cle{noyau} de $f$ est $\ker f=\{u\in E \mid f(u)=0_F\}$. C'est un sous-espace vectoriel de $E$.
\item L'\cle{image} de $f$ est $\operatorname{Im} f=\{f(u)\mid u\in E\}$. C'est un sous-espace vectoriel de $F$.
\end{itemize}
\end{definition}

\begin{methode}[Noyau et image : équations caractéristiques et bases]
Soit $f:E\to F$ linéaire (avec $E,F\in\{\mathbb{R},\mathbb{R}^2,\mathbb{R}^3,\text{vecteurs plan},\text{vecteurs espace}\}$).
\begin{enumerate}
\item \textbf{Noyau ($\ker f$) :}
      \begin{itemize}
      \item Écrire $f(x,y,\ldots)=(0,0,\ldots)$ : c'est le système homogène, appelé \cle{équation(s) caractéristique(s) du noyau}.
      \item Résoudre le système, exprimer la solution générale en fonction de paramètres.
      \item Les vecteurs directeurs des paramètres forment une famille génératrice de $\ker f$ ; en extraire une base (vérifier la liberté).
      \item $\dim\ker f$ = nombre de paramètres (degrés de liberté).
      \end{itemize}
\item \textbf{Image ($\operatorname{Im} f$) :}
      \begin{itemize}
      \item Écrire le vecteur générique $f(x,y,\ldots)$ comme combinaison linéaire de vecteurs fixes (souvent les images des vecteurs d'une base de $E$).
      \item Ces vecteurs engendrent $\operatorname{Im} f$ ; en extraire une famille libre (base).
      \item L'\cle{équation caractéristique de l'image} s'obtient en cherchant la condition sur $(a,b,\ldots)\in F$ pour que $f(x,y,\ldots)=(a,b,\ldots)$ ait une solution (compatibilité du système).
      \end{itemize}
\item \textbf{Bijectivité :} $f$ injective $\iff \ker f=\{0_E\}$ ; $f$ surjective $\iff \operatorname{Im} f=F$.
\item \textbf{Théorème du rang (vérification) :} $\dim E=\dim\ker f+\dim\operatorname{Im} f$.
\end{enumerate}
\textbf{Important :} Pour un \emph{endomorphisme} en dimension finie, injectif $\iff$ surjectif $\iff$ bijectif.
\end{methode}

\begin{exoresolu}[Noyau, image et bijectivité]
Soit $f:\mathbb{R}^3\to\mathbb{R}^2$ définie par $f(x,y,z)=(x+y,\;y+z)$.
\begin{enumerate}
\item Déterminer une équation caractéristique et une base de $\ker f$.
\item Déterminer une base de $\operatorname{Im} f$. L'application $f$ est-elle surjective ? injective ?
\end{enumerate}
\tcblower
\textbf{1. Noyau.} $(x,y,z)\in\ker f\iff f(x,y,z)=(0,0)$, c'est-à-dire
\[ \begin{cases}x+y=0\\y+z=0\end{cases}\quad\text{(équations caractéristiques de }\ker f\text{)}. \]
On obtient $x=-y$ et $z=-y$. Le vecteur générique du noyau est donc
\[ (x,y,z)=(-y,\;y,\;-y)=y(-1,1,-1). \]
Ainsi $\ker f=\operatorname{Vect}\big((-1,1,-1)\big)$. Le vecteur $(-1,1,-1)$ est non nul, donc la famille $\big((-1,1,-1)\big)$ est libre : c'est une \cle{base} de $\ker f$, et $\dim\ker f=1$.

Comme $\ker f\neq\{(0,0,0)\}$, $f$ n'est \textbf{pas injective}.

\textbf{2. Image.} Pour tout $(x,y,z)\in\mathbb{R}^3$ :
\[ f(x,y,z)=(x+y,\;y+z)=x(1,0)+y(1,1)+z(0,1). \]
Donc $\operatorname{Im} f=\operatorname{Vect}\big((1,0),(1,1),(0,1)\big)$. Or $(1,1)=(1,0)+(0,1)$, donc $\operatorname{Im} f=\operatorname{Vect}\big((1,0),(0,1)\big)$. Les vecteurs $(1,0)$ et $(0,1)$ ne sont pas colinéaires : ils forment une famille libre, donc une \cle{base} de $\operatorname{Im} f$, et $\dim\operatorname{Im} f=2$.

Comme $\dim\operatorname{Im} f=2=\dim\mathbb{R}^2$, on a $\operatorname{Im} f=\mathbb{R}^2$ : $f$ est \textbf{surjective}.

\textbf{Vérification (théorème du rang) :} $\dim\ker f+\dim\operatorname{Im} f=1+2=3=\dim\mathbb{R}^3$. Cohérent.

\textbf{Conclusion.} $\ker f$ admet pour base $\big((-1,1,-1)\big)$, $\operatorname{Im} f$ admet pour base $\big((1,0),(0,1)\big)$ ; $f$ est surjective mais non injective, donc non bijective.
\end{exoresolu}

\courssub{Expression analytique d'une application linéaire dans une base}

\begin{methode}[Retrouver l'écriture analytique d'une application linéaire]
Une application linéaire est entièrement déterminée par les images des vecteurs d'une base. Pour obtenir son \cle{expression analytique} $f(x,y,\ldots)$ :
\begin{enumerate}
\item Écrire le vecteur générique $u=(x,y)$ comme combinaison linéaire des vecteurs de la base donnée : $u=\lambda e_1+\mu e_2$ (résoudre le système pour trouver $\lambda,\mu$ en fonction de $x,y$).
\item Utiliser la linéarité : $f(u)=\lambda f(e_1)+\mu f(e_2)$.
\item Développer et regrouper les coordonnées pour obtenir $f(x,y)=(\ldots,\ldots)$.
\end{enumerate}
\textbf{Réflexe (base canonique) :} Si la base est la base canonique $e_1=(1,0)$, $e_2=(0,1)$, alors directement $f(x,y)=x\,f(e_1)+y\,f(e_2)$ (pas de système à résoudre).
\end{methode}

\begin{exoresolu}[Déterminer l'expression analytique]
Soit $f$ l'endomorphisme de $\mathbb{R}^2$ tel que $f(1,0)=(2,1)$ et $f(1,1)=(0,3)$. Déterminer l'expression analytique de $f(x,y)$ pour tout $(x,y)\in\mathbb{R}^2$.
\tcblower
\textbf{Étape 1 : décomposition sur la famille $((1,0),(1,1))$.} Cette famille est une base de $\mathbb{R}^2$ (deux vecteurs non colinéaires). Soit $(x,y)\in\mathbb{R}^2$ ; on cherche $\lambda,\mu$ tels que
\[ (x,y)=\lambda(1,0)+\mu(1,1)\iff\begin{cases}\lambda+\mu=x\\\mu=y\end{cases}\iff \mu=y,\quad \lambda=x-y. \]
Donc $(x,y)=(x-y)(1,0)+y(1,1)$.

\textbf{Étape 2 : linéarité.}
\begin{align*}
f(x,y)&=(x-y)f(1,0)+y\,f(1,1)\\
&=(x-y)(2,1)+y(0,3)\\
&=\big(2(x-y),\;(x-y)+3y\big)\\
&=(2x-2y,\;x+2y).
\end{align*}

\textbf{Conclusion.} Pour tout $(x,y)\in\mathbb{R}^2$, $f(x,y)=(2x-2y,\;x+2y)$.

\textbf{Vérification :} $f(1,0)=(2,1)$ et $f(1,1)=(0,3)$ : conforme aux données.
\end{exoresolu}

\begin{savaistu}[Algèbre linéaire et télécommunications au Cameroun]
Les espaces vectoriels et les applications linéaires sont au cœur du traitement du signal numérique. Lorsque vous passez un appel sur le réseau MTN ou Orange, votre voix est échantillonnée, transformée en vecteurs (trames de bits), puis codée par des matrices (applications linéaires) pour la compression (codecs AMR, EVS) et la correction d'erreurs (codes de Reed-Solomon, LDPC). Le noyau de la matrice de codage correspond aux syndromes d'erreur, et son image à l'espace des mots de code valides. La dimension de ces espaces détermine le débit et la robustesse de la communication. L'algèbre linéaire, c'est aussi ce qui permet au GPS de votre téléphone de trianguler votre position à Yaoundé ou Douala en résolvant des systèmes linéaires issus de mesures de distances aux satellites !
\end{savaistu}

\begin{attention}[Les 5 erreurs qui coûtent des points au Probatoire]
\begin{enumerate}
\item \textbf{Oublier de vérifier que l'ensemble est non vide} dans la caractérisation des sous-espaces vectoriels. La stabilité seule ne suffit pas : l'ensemble vide est « stable » mais n'est pas un sous-espace vectoriel. Écris toujours « $0_E\in F$ car \ldots ».
\item \textbf{Confondre libre et génératrice.} Une famille libre ne « couvre » pas forcément l'espace ; une famille génératrice peut contenir des vecteurs « en trop ». Pour une base, il faut prouver \emph{les deux} propriétés, sauf si tu connais déjà la dimension (alors une seule suffit, à condition de le justifier explicitement : « car $\dim E=n$ »).
\item \textbf{Mal rédiger la liberté.} Il faut partir de l'hypothèse $\lambda_1 u_1+\cdots+\lambda_p u_p=0_E$ et \emph{conclure} que tous les $\lambda_i$ sont nuls. Partir de « $\lambda_1=\cdots=\lambda_p=0$ donne $0_E$ » prouve le contraire de ce qui est demandé (raisonnement à l'envers), sanctionné systématiquement.
\item \textbf{Croire que $\ker f=\{0\}$ suffit pour la bijectivité.} Cela ne donne que l'\emph{injectivité}. Il faut aussi la surjectivité ($\operatorname{Im} f=F$) — sauf pour un \emph{endomorphisme en dimension finie}, où l'une entraîne l'autre (théorème du rang), ce qui doit être \emph{mentionné} dans la copie.
\item \textbf{Erreurs de calcul dans les systèmes.} Une erreur de signe dans la résolution du système du noyau fausse toute la base. Vérifie systématiquement que chaque vecteur de la base trouvée annule bien $f$ (pour le noyau) ou appartient bien à l'image, et contrôle tes dimensions avec le théorème du rang : $\dim E=\dim\ker f+\dim\operatorname{Im} f$.
\end{enumerate}
\end{attention}

\newpage

\coursec{Exercices}

\courssub{Je vérifie mes connaissances}

\begin{exercice}[QCM : espaces vectoriels]
Pour chaque affirmation, choisir la (ou les) bonne(s) réponse(s) en justifiant brièvement.
\begin{enumerate}
\item L'ensemble $F=\{(x,y)\in\mathbb{R}^2 \mid x+y=1\}$ est :
\begin{enumerate}
\item un sous-espace vectoriel de $\mathbb{R}^2$ ;
\item une partie de $\mathbb{R}^2$ stable par addition ;
\item une partie de $\mathbb{R}^2$ contenant le vecteur nul ;
\item aucune des réponses précédentes.
\end{enumerate}
\item La famille $\big((1,2),(2,4)\big)$ de $\mathbb{R}^2$ est :
\begin{enumerate}
\item libre ; \item liée ; \item génératrice de $\mathbb{R}^2$ ; \item une base de $\mathbb{R}^2$.
\end{enumerate}
\item Si $f:\mathbb{R}^3\to\mathbb{R}^2$ est une application linéaire surjective, alors :
\begin{enumerate}
\item $\dim(\ker f)=1$ ; \item $\dim(\ker f)=2$ ; \item $f$ est injective ; \item $\dim(\operatorname{Im} f)=2$.
\end{enumerate}
\item Un endomorphisme de $\mathbb{R}^2$ dont le noyau est réduit au vecteur nul est :
\begin{enumerate}
\item nécessairement nul ; \item injectif ; \item un automorphisme ; \item surjectif.
\end{enumerate}
\end{enumerate}
\end{exercice}

\begin{exercice}[Vrai ou faux]
Répondre par \textbf{vrai} ou \textbf{faux} en justifiant chaque réponse par un argument ou un contre-exemple.
\begin{enumerate}
\item Toute famille de trois vecteurs de $\mathbb{R}^2$ est liée.
\item Toute famille libre de deux vecteurs de $\mathbb{R}^2$ est une base de $\mathbb{R}^2$.
\item L'ensemble $\{(x,y)\in\mathbb{R}^2 \mid xy=0\}$ est un sous-espace vectoriel de $\mathbb{R}^2$.
\item Si $f$ et $g$ sont deux applications linéaires de $\mathbb{R}^2$ dans $\mathbb{R}^2$, alors $f+g$ est linéaire.
\item L'image d'une application linéaire $f:E\to F$ est un sous-espace vectoriel de $F$.
\end{enumerate}
\end{exercice}

\begin{exercice}[Définitions et vocabulaire]
\begin{enumerate}
\item Donner la définition d'une loi de composition externe sur un ensemble $E$, à domaine d'opérateurs $\mathbb{R}$.
\item Donner la définition d'une application linéaire $f:E\to F$ entre deux espaces vectoriels réels.
\item Compléter : un endomorphisme bijectif s'appelle un \ldots\ ; une application linéaire bijective de $E$ vers $F$ s'appelle un \ldots
\item Énoncer le théorème du rang pour une application linéaire $f:E\to F$ où $E$ est de dimension finie.
\end{enumerate}
\end{exercice}

\begin{exercice}[Calculs directs]
Soit $f:\mathbb{R}^2\to\mathbb{R}^2$ définie par $f(x,y)=(3x-y,\;x+2y)$.
\begin{enumerate}
\item Calculer $f(1,0)$, $f(0,1)$ et $f(2,-1)$.
\item Calculer $f\big(2(1,1)+3(0,2)\big)$ de deux manières différentes.
\item Déterminer les couples $(x,y)$ tels que $f(x,y)=(0,0)$. Que peut-on en conclure ?
\end{enumerate}
\end{exercice}

\courssub{Je m'entraîne}

\begin{exercice}[Un sous-espace vectoriel de $\mathbb{R}^3$]
On considère l'ensemble $F=\{(x,y,z)\in\mathbb{R}^3 \mid x+2y-z=0\}$.
\begin{enumerate}
\item Montrer que $F$ est un sous-espace vectoriel de $\mathbb{R}^3$.
\item Montrer que $F=\{(x,y,x+2y)\mid (x,y)\in\mathbb{R}^2\}$, puis en déduire une famille génératrice de $F$ formée de deux vecteurs.
\item Montrer que cette famille est libre. Conclure sur une base et la dimension de $F$.
\end{enumerate}
\end{exercice}

\begin{exercice}[Libre, liée, base ?]
On considère dans $\mathbb{R}^3$ la famille $\mathcal{F}=\big((1,0,1),(1,1,0),(0,1,1)\big)$.
\begin{enumerate}
\item La famille $\mathcal{F}$ est-elle libre ? (On posera $\alpha(1,0,1)+\beta(1,1,0)+\gamma(0,1,1)=(0,0,0)$ et on résoudra le système obtenu.)
\item La famille $\mathcal{F}$ est-elle une base de $\mathbb{R}^3$ ? Justifier.
\item Exprimer le vecteur $(2,1,3)$ dans cette base.
\end{enumerate}
\end{exercice}

\begin{exercice}[Un endomorphisme de $\mathbb{R}^2$]
Soit $f:\mathbb{R}^2\to\mathbb{R}^2$ définie par $f(x,y)=(3x-y,\;x+2y)$.
\begin{enumerate}
\item Montrer que $f$ est une application linéaire. Que dire alors de $f$ ?
\item Déterminer $\ker f$ (on donnera une équation caractéristique).
\item En déduire que $f$ est un automorphisme de $\mathbb{R}^2$.
\item Déterminer l'antécédent du vecteur $(5,3)$ par $f$.
\end{enumerate}
\end{exercice}

\begin{exercice}[Noyau et image]
Soit $g:\mathbb{R}^3\to\mathbb{R}^2$ définie par $g(x,y,z)=(x+y,\;y-z)$.
\begin{enumerate}
\item Vérifier que $g$ est une application linéaire.
\item Déterminer une équation caractéristique de $\ker g$, puis une base de $\ker g$.
\item Déterminer une base de $\operatorname{Im} g$ (on pourra calculer les images des vecteurs de la base canonique de $\mathbb{R}^3$).
\item $g$ est-elle injective ? surjective ? Vérifier le théorème du rang.
\end{enumerate}
\end{exercice}

\courssub{Je me prépare au Probatoire}

\begin{exercice}[La droite des coûts (type Probatoire)]
Une menuiserie de Bafoussam fabrique deux articles. Le coût de production (en milliers de francs CFA) de $x$ unités du premier article et $y$ unités du second est modélisé par le couple $(x,y)\in\mathbb{R}^2$. On dit qu'une production est \og équilibrée \fg{} lorsque $y=3x$. On note $E=\{(x,y)\in\mathbb{R}^2 \mid y=3x\}$.

\textbf{Partie A : structure de $E$}
\begin{enumerate}
\item Montrer que $E$, muni de l'addition usuelle de $\mathbb{R}^2$ et de la multiplication par un réel, est un espace vectoriel sur $\mathbb{R}$.
\item Donner une base de $E$ et préciser sa dimension.
\end{enumerate}

\textbf{Partie B : une application linéaire}
On considère l'application $h:\mathbb{R}^2\to\mathbb{R}^2$ définie par $h(x,y)=(x-2y,\;-2x+4y)$.
\begin{enumerate}
\item Montrer que $h$ est un endomorphisme de $\mathbb{R}^2$.
\item Déterminer une équation caractéristique de $\ker h$, puis une base de $\ker h$.
\item Déterminer une équation caractéristique de $\operatorname{Im} h$, puis une base de $\operatorname{Im} h$.
\item $h$ est-elle bijective ? Justifier.
\item Résoudre dans $\mathbb{R}^2$ l'équation $h(x,y)=(3,-6)$. Interpréter géométriquement l'ensemble des solutions.
\end{enumerate}
\end{exercice}

\begin{exercice}[Symétrie vectorielle et automorphisme (type Probatoire)]
Le plan vectoriel $\mathcal{P}$ est muni d'une base $(\vec{i},\vec{j})$. On considère la symétrie vectorielle $s$ par rapport à la droite de base $\vec{i}$, parallèlement à la droite de base $\vec{j}$ : pour tout vecteur $\vec{u}=x\vec{i}+y\vec{j}$, on a $s(\vec{u})=x\vec{i}-y\vec{j}$.

\textbf{Partie A : étude de $s$}
\begin{enumerate}
\item Donner l'expression analytique de $s$ dans la base $(\vec{i},\vec{j})$.
\item Montrer que $s$ est un endomorphisme de $\mathcal{P}$.
\item Déterminer $\ker s$ et $\operatorname{Im} s$.
\item En déduire que $s$ est un automorphisme de $\mathcal{P}$.
\item Calculer $s\circ s$. Quelle propriété remarquable de $s$ retrouve-t-on ?
\end{enumerate}

\textbf{Partie B : une famille de vecteurs}
On pose $\vec{u}=\vec{i}+\vec{j}$ et $\vec{v}=s(\vec{u})$.
\begin{enumerate}
\item Montrer que $(\vec{u},\vec{v})$ est une base de $\mathcal{P}$.
\item Déterminer les coordonnées de $s(\vec{u})$ et $s(\vec{v})$ dans la base $(\vec{u},\vec{v})$.
\end{enumerate}
\end{exercice}

\begin{exercice}[Réseau de distribution d'eau (type Probatoire)]
Une société de distribution d'eau étudie les flux entre trois quartiers de Douala. Les volumes d'eau (en m$^3$) distribués sont modélisés par des triplets $(x,y,z)\in\mathbb{R}^3$. On considère l'application $f:\mathbb{R}^3\to\mathbb{R}^3$ définie par
\[ f(x,y,z)=(x+y+z,\;x-y,\;2x+z). \]

\textbf{Partie A : linéarité}
\begin{enumerate}
\item Montrer que $f$ est un endomorphisme de $\mathbb{R}^3$.
\item Calculer les images par $f$ des vecteurs de la base canonique $(\vec{e_1},\vec{e_2},\vec{e_3})$ de $\mathbb{R}^3$.
\end{enumerate}

\textbf{Partie B : noyau et image}
\begin{enumerate}
\item Déterminer une équation caractéristique de $\ker f$ et montrer que $\ker f=\{(0,0,0)\}$.
\item En déduire que $f$ est un automorphisme de $\mathbb{R}^3$.
\item Quelle est la dimension de $\operatorname{Im} f$ ? Donner une base de $\operatorname{Im} f$.
\end{enumerate}

\textbf{Partie C : exploitation}
\begin{enumerate}
\item Déterminer l'antécédent par $f$ du triplet $(6,2,7)$, qui représente les volumes souhaités dans les trois quartiers.
\item On considère la famille $\mathcal{B}=\big((1,1,2),(1,-1,0),(1,0,1)\big)$. Montrer que $\mathcal{B}$ est une base de $\mathbb{R}^3$. Que représente cette famille pour $f$ ?
\end{enumerate}
\end{exercice}

\newpage

\coursec{Corrigés des exercices}

\coursec{Exercices corrigés — Espaces vectoriels sur $\mathbb{R}$}

\begin{corrige}[Exercice 1 — QCM : espaces vectoriels]
\textbf{1.} L'ensemble $F=\{(x,y)\in\mathbb{R}^2 \mid x+y=1\}$.
\begin{itemize}
    \item \textbf{Contient le vecteur nul ?} $\vec{0}=(0,0)$ ne satisfait pas $0+0=1$. Donc $\vec{0}\notin F$. \textbf{Faux}.
    \item \textbf{Stable par addition ?} Prenons $(1,0)\in F$ et $(0,1)\in F$. Leur somme $(1,1)$ ne satisfait pas $1+1=1$. \textbf{Faux}.
    \item \textbf{Sous-espace vectoriel ?} Non, car ne contient pas $\vec{0}$ (condition nécessaire). \textbf{Faux}.
\end{itemize}
\cle{Bonne réponse :} \textbf{aucune des réponses précédentes}.

\textbf{2.} Famille $\mathcal{F}=\big((1,2),(2,4)\big)$ dans $\mathbb{R}^2$.
On observe que $(2,4)=2\times(1,2)$. Les vecteurs sont \cle{colinéaires} (proportionnels).
\begin{itemize}
    \item \textbf{Libre ?} Non, car $\exists (\alpha,\beta)\neq(0,0)$ tel que $\alpha(1,2)+\beta(2,4)=\vec{0}$ (ex: $\alpha=2,\beta=-1$). \textbf{Faux}.
    \item \textbf{Liée ?} \textbf{Vrai}.
    \item \textbf{Génératrice de $\mathbb{R}^2$ ?} Non, elle ne génère qu'une droite (dimension 1). \textbf{Faux}.
    \item \textbf{Base de $\mathbb{R}^2$ ?} Non (ni libre, ni génératrice). \textbf{Faux}.
\end{itemize}
\cle{Bonne réponse :} \textbf{liée}.

\textbf{3.} $f:\mathbb{R}^3\to\mathbb{R}^2$ linéaire et surjective.
$\dim(\mathbb{R}^3)=3$, $\dim(\mathbb{R}^2)=2$. Surjective $\iff \operatorname{Im} f = \mathbb{R}^2 \iff \dim(\operatorname{Im} f)=2$.
\cle{Théorème du rang} : $\dim(\ker f) + \dim(\operatorname{Im} f) = \dim(\mathbb{R}^3) = 3$.
Donc $\dim(\ker f) = 3 - 2 = 1$.
\begin{itemize}
    \item $\dim(\ker f)=1$ : \textbf{Vrai}.
    \item $\dim(\ker f)=2$ : \textbf{Faux}.
    \item $f$ injective ? Non, car $\ker f \neq \{\vec{0}\}$ (dimension 1). \textbf{Faux}.
    \item $\dim(\operatorname{Im} f)=2$ : \textbf{Vrai} (par surjectivité).
\end{itemize}
\cle{Bonnes réponses :} $\dim(\ker f)=1$ et $\dim(\operatorname{Im} f)=2$.

\textbf{4.} Endomorphisme $f$ de $\mathbb{R}^2$ tel que $\ker f = \{\vec{0}\}$.
$\ker f = \{\vec{0}\} \iff f$ \cle{injective}. Comme $f$ est un endomorphisme d'un espace de dimension finie, injective $\iff$ bijective $\iff$ \cle{automorphisme} $\iff$ surjective.
\begin{itemize}
    \item Nécessairement nul ? \textbf{Faux} (ex: identité).
    \item Injectif ? \textbf{Vrai}.
    \item Automorphisme ? \textbf{Vrai}.
    \item Surjectif ? \textbf{Vrai}.
\end{itemize}
\cle{Bonnes réponses :} injectif, automorphisme, surjectif.
\end{corrige}

\begin{corrige}[Exercice 2 — Vrai ou faux]
\textbf{1. VRAI.} $\dim(\mathbb{R}^2)=2$. Toute famille de $3$ vecteurs dans un espace de dimension $2$ est \cle{liée} (théorème : toute famille de plus de $n$ vecteurs dans un espace de dimension $n$ est liée).

\textbf{2. VRAI.} Une famille \cle{libre} de $2$ vecteurs dans $\mathbb{R}^2$ (dimension $2$) est une \cle{base} (théorème de la base incomplète : toute famille libre de cardinal égal à la dimension est une base).

\textbf{3. FAUX.} $E=\{(x,y)\in\mathbb{R}^2 \mid xy=0\}$ est la réunion des deux axes. Prenons $\vec{u}=(1,0)\in E$ et $\vec{v}=(0,1)\in E$. $\vec{u}+\vec{v}=(1,1)\notin E$ car $1\times 1 \neq 0$. $E$ n'est pas stable par addition, ce n'est pas un sous-espace vectoriel.

\textbf{4. VRAI.} Soit $f,g$ linéaires. Pour tous $\vec{u},\vec{v}\in\mathbb{R}^2$, $\lambda\in\mathbb{R}$ :
$(f+g)(\vec{u}+\vec{v}) = f(\vec{u}+\vec{v})+g(\vec{u}+\vec{v}) = f(\vec{u})+f(\vec{v})+g(\vec{u})+g(\vec{v}) = (f+g)(\vec{u})+(f+g)(\vec{v})$.
$(f+g)(\lambda\vec{u}) = f(\lambda\vec{u})+g(\lambda\vec{u}) = \lambda f(\vec{u})+\lambda g(\vec{u}) = \lambda(f+g)(\vec{u})$.
Donc $f+g$ est linéaire.

\textbf{5. VRAI.} Propriété fondamentale : l'image d'une application linéaire $f:E\to F$ est un \cle{sous-espace vectoriel} de $F$ (contient $\vec{0}_F$, stable par addition et multiplication par un scalaire).
\end{corrige}

\begin{corrige}[Exercice 3 — Définitions et vocabulaire]
\textbf{1.} Une \cle{loi de composition externe} sur un ensemble $E$ à domaine d'opérateurs $\mathbb{R}$ est une application
\[ \cdot : \mathbb{R} \times E \longrightarrow E, \quad (\lambda, x) \longmapsto \lambda \cdot x \]
telle que pour tous $\lambda,\mu\in\mathbb{R}$ et $x,y\in E$ :
\begin{itemize}
    \item $\lambda \cdot (x+y) = \lambda \cdot x + \lambda \cdot y$ (distributivité par rapport à l'addition dans $E$),
    \item $(\lambda+\mu) \cdot x = \lambda \cdot x + \mu \cdot x$ (distributivité par rapport à l'addition dans $\mathbb{R}$),
    \item $(\lambda\mu) \cdot x = \lambda \cdot (\mu \cdot x)$ (associativité / compatibilité avec la multiplication dans $\mathbb{R}$),
    \item $1 \cdot x = x$ (neutre de la multiplication dans $\mathbb{R}$).
\end{itemize}

\textbf{2.} Soient $E,F$ deux espaces vectoriels sur $\mathbb{R}$. Une application $f\colon E\to F$ est \cle{linéaire} si elle vérifie :
\begin{enumerate}
    \item $\forall \vec{u},\vec{v}\in E,\quad f(\vec{u}+\vec{v}) = f(\vec{u})+f(\vec{v})$ (stabilité par addition / additivité),
    \item $\forall \lambda\in\mathbb{R}, \forall \vec{u}\in E,\quad f(\lambda\vec{u}) = \lambda f(\vec{u})$ (stabilité par multiplication par un scalaire / homogénéité).
\end{enumerate}
Équivalentement : $\forall \vec{u},\vec{v}\in E, \forall \lambda,\mu\in\mathbb{R},\quad f(\lambda\vec{u}+\mu\vec{v}) = \lambda f(\vec{u})+\mu f(\vec{v})$.

\textbf{3.}
\begin{itemize}
    \item Un endomorphisme bijectif s'appelle un \textbf{automorphisme}.
    \item Une application linéaire bijective de $E$ vers $F$ s'appelle un \textbf{isomorphisme}.
\end{itemize}

\textbf{4.} \cle{Théorème du rang} : Soit $f\colon E\to F$ une application linéaire avec $E$ de dimension finie $n$. Alors
\[ \dim(E) = \dim(\ker f) + \dim(\operatorname{Im} f). \]
On note $\operatorname{rg}(f) = \dim(\operatorname{Im} f)$ le \cle{rang} de $f$.
\end{corrige}

\begin{corrige}[Exercice 4 — Calculs directs]
Soit $f(x,y)=(3x-y,\;x+2y)$.

\textbf{1.}
\begin{align*}
f(1,0) &= (3\times1-0,\;1+2\times0) = (3,1), \\
f(0,1) &= (3\times0-1,\;0+2\times1) = (-1,2), \\
f(2,-1) &= (3\times2-(-1),\;2+2\times(-1)) = (6+1,\;2-2) = (7,0).
\end{align*}

\textbf{2.} Calculons $f\big(2(1,1)+3(0,2)\big)$.
\begin{itemize}
    \item \textbf{Méthode directe :} $2(1,1)+3(0,2) = (2,2)+(0,6) = (2,8)$.
    $f(2,8) = (3\times2-8,\;2+2\times8) = (6-8,\;2+16) = (-2,18)$.
    \item \textbf{Par linéarité :} $f\big(2(1,1)+3(0,2)\big) = 2f(1,1) + 3f(0,2)$.
    $f(1,1) = (3-1,\;1+2) = (2,3)$.
    $f(0,2) = (0-2,\;0+4) = (-2,4)$.
    $2(2,3) + 3(-2,4) = (4,6) + (-6,12) = (-2,18)$.
\end{itemize}
On retrouve le même résultat, confirmant la linéarité de $f$.

\textbf{3.} Résolution de $f(x,y)=(0,0)$ :
\[ \begin{cases} 3x - y = 0 \\ x + 2y = 0 \end{cases} \iff \begin{cases} y = 3x \\ x + 6x = 0 \end{cases} \iff \begin{cases} 7x = 0 \\ y = 0 \end{cases} \iff (x,y)=(0,0). \]
Donc $\ker f = \{\vec{0}\}$. On en conclut que $f$ est \cle{injective} (noyau réduit au vecteur nul). Comme $f$ est un endomorphisme de $\mathbb{R}^2$ (espace de dimension finie), $f$ est bijective, donc c'est un \cle{automorphisme}.
\end{corrige}

\begin{corrige}[Exercice 5 — Un sous-espace vectoriel de $\mathbb{R}^3$]
$F=\{(x,y,z)\in\mathbb{R}^3 \mid x+2y-z=0\}$.

\textbf{1.} Montrons que $F$ est un sous-espace vectoriel de $\mathbb{R}^3$ (critère de stabilité).
\begin{itemize}
    \item \textbf{Non vide / Vecteur nul :} $(0,0,0)$ vérifie $0+2\times0-0=0$, donc $\vec{0}\in F$.
    \item \textbf{Stable par addition :} Soient $\vec{u}=(x_1,y_1,z_1)$ et $\vec{v}=(x_2,y_2,z_2)$ dans $F$.
    $x_1+2y_1-z_1=0$ et $x_2+2y_2-z_2=0$.
    $\vec{u}+\vec{v} = (x_1+x_2,\; y_1+y_2,\; z_1+z_2)$.
    $(x_1+x_2)+2(y_1+y_2)-(z_1+z_2) = (x_1+2y_1-z_1)+(x_2+2y_2-z_2)=0+0=0$.
    Donc $\vec{u}+\vec{v}\in F$.
    \item \textbf{Stable par multiplication par un scalaire :} Soit $\lambda\in\mathbb{R}$, $\vec{u}=(x,y,z)\in F$.
    $\lambda\vec{u}=(\lambda x,\lambda y,\lambda z)$.
    $\lambda x + 2\lambda y - \lambda z = \lambda(x+2y-z)=\lambda\times0=0$.
    Donc $\lambda\vec{u}\in F$.
\end{itemize}
$F$ est un sous-espace vectoriel de $\mathbb{R}^3$.

\textbf{2.} L'équation $x+2y-z=0 \iff z = x+2y$.
Tout vecteur de $F$ s'écrit $(x,y,x+2y)$ avec $(x,y)\in\mathbb{R}^2$.
Réciproquement, tout couple $(x,y)$ définit un vecteur de $F$.
Donc $F = \{(x,y,x+2y) \mid (x,y)\in\mathbb{R}^2\}$.
Décomposition : $(x,y,x+2y) = x(1,0,1) + y(0,1,2)$.
La famille $\mathcal{G} = \big((1,0,1),\;(0,1,2)\big)$ est une \cle{famille génératrice} de $F$.

\textbf{3.} Montrons que $\mathcal{G}$ est \cle{libre}.
Soient $\alpha,\beta\in\mathbb{R}$ tels que $\alpha(1,0,1)+\beta(0,1,2)=\vec{0}$.
$(\alpha,\;\beta,\;\alpha+2\beta) = (0,0,0) \implies \alpha=0,\; \beta=0,\; \alpha+2\beta=0$.
Le système donne $\alpha=0$ et $\beta=0$. La famille est libre.
$\mathcal{G}$ est une famille libre et génératrice de $F$, c'est une \cle{base} de $F$.
$\dim(F) = \text{card}(\mathcal{G}) = 2$.
\end{corrige}

\begin{corrige}[Exercice 6 — Libre, liée, base ?]
$\mathcal{F}=\big(\vec{u_1}=(1,0,1),\;\vec{u_2}=(1,1,0),\;\vec{u_3}=(0,1,1)\big)$ dans $\mathbb{R}^3$.

\textbf{1.} Étude de la liberté.
$\alpha\vec{u_1}+\beta\vec{u_2}+\gamma\vec{u_3}=\vec{0} \iff \alpha(1,0,1)+\beta(1,1,0)+\gamma(0,1,1)=(0,0,0)$.
\[ \begin{cases} \alpha + \beta = 0 \\ \beta + \gamma = 0 \\ \alpha + \gamma = 0 \end{cases} \]
Somme des trois équations : $2(\alpha+\beta+\gamma)=0 \implies \alpha+\beta+\gamma=0$.
En soustrayant la 1ère : $\gamma=0$. Alors $\beta=0$ (2ème), $\alpha=0$ (1ère).
La seule solution est triviale : la famille $\mathcal{F}$ est \cle{libre}.

\textbf{2.} $\mathcal{F}$ est une famille libre de $3$ vecteurs dans $\mathbb{R}^3$ qui est de dimension $3$.
D'après le \cle{théorème de la base incomplète}, toute famille libre de cardinal égal à la dimension de l'espace est une \cle{base}.
Donc $\mathcal{F}$ est une base de $\mathbb{R}^3$.

\textbf{3.} Coordonnées de $\vec{v}=(2,1,3)$ dans la base $\mathcal{F}$.
On cherche $x,y,z$ tels que $x\vec{u_1}+y\vec{u_2}+z\vec{u_3}=\vec{v}$.
\[ \begin{cases} x + y = 2 \\ y + z = 1 \\ x + z = 3 \end{cases} \]
Somme : $2(x+y+z)=6 \implies x+y+z=3$.
$z = 3 - (x+y) = 3-2 = 1$.
$y = 1 - z = 0$.
$x = 2 - y = 2$.
Vérification 3ème équation : $x+z=2+1=3$ (OK).
Coordonnées : $\boxed{(2,\;0,\;1)}$.
\end{corrige}

\begin{corrige}[Exercice 7 — Un endomorphisme de $\mathbb{R}^2$]
$f(x,y)=(3x-y,\;x+2y)$.

\textbf{1.} Linéarité :
$\forall (x_1,y_1),(x_2,y_2)\in\mathbb{R}^2,\; \forall \lambda,\mu\in\mathbb{R}$,
$f(\lambda(x_1,y_1)+\mu(x_2,y_2)) = f(\lambda x_1+\mu x_2,\; \lambda y_1+\mu y_2)$
$= \big(3(\lambda x_1+\mu x_2)-(\lambda y_1+\mu y_2),\; (\lambda x_1+\mu x_2)+2(\lambda y_1+\mu y_2)\big)$
$= \lambda(3x_1-y_1,\;x_1+2y_1) + \mu(3x_2-y_2,\;x_2+2y_2)$
$= \lambda f(x_1,y_1) + \mu f(x_2,y_2)$.
$f$ est linéaire. Comme $f:\mathbb{R}^2\to\mathbb{R}^2$, $f$ est un \cle{endomorphisme} de $\mathbb{R}^2$.

\textbf{2.} $\ker f = \{(x,y)\in\mathbb{R}^2 \mid f(x,y)=(0,0)\}$.
Système : $\begin{cases} 3x-y=0 \\ x+2y=0 \end{cases} \implies y=3x,\; x+6x=0 \implies 7x=0 \implies x=0,y=0$.
\cle{Équation caractéristique} : $\begin{cases} 3x-y=0 \\ x+2y=0 \end{cases}$.
$\ker f = \{\vec{0}\}$.

\textbf{3.} $\ker f = \{\vec{0}\} \implies f$ injective. $f$ endomorphisme de $\mathbb{R}^2$ (dimension finie), injective $\implies$ bijective $\implies$ \cle{automorphisme} de $\mathbb{R}^2$.

\textbf{4.} Antécédent de $(5,3)$ : résoudre $f(x,y)=(5,3)$.
$\begin{cases} 3x-y=5 \\ x+2y=3 \end{cases} \implies y=3x-5$.
$x+2(3x-5)=3 \implies 7x-10=3 \implies 7x=13 \implies x=\frac{13}{7}$.
$y = 3\times\frac{13}{7}-5 = \frac{39-35}{7} = \frac{4}{7}$.
L'antécédent est $\boxed{\left(\frac{13}{7},\;\frac{4}{7}\right)}$ (unique car $f$ bijective).
\end{corrige}

\begin{corrige}[Exercice 8 — Noyau et image]
$g:\mathbb{R}^3\to\mathbb{R}^2,\; g(x,y,z)=(x+y,\;y-z)$.

\textbf{1.} Linéarité : immédiate car chaque composante est une forme linéaire (combinaison linéaire des coordonnées). $g$ est une application linéaire.

\textbf{2.} $\ker g = \{(x,y,z)\mid x+y=0,\; y-z=0\}$.
\cle{Équation caractéristique} : $\begin{cases} x+y=0 \\ y-z=0 \end{cases} \iff \begin{cases} x=-y \\ z=y \end{cases}$.
Vecteurs : $(-y,\;y,\;y) = y(-1,1,1),\quad y\in\mathbb{R}$.
$\ker g = \operatorname{Vect}\big((-1,1,1)\big)$.
Une base de $\ker g$ : $\boxed{\mathcal{B}_{\ker}=((-1,1,1))}$. $\dim(\ker g)=1$.

\textbf{3.} $\operatorname{Im} g = \operatorname{Vect}\big(g(\vec{e_1}), g(\vec{e_2}), g(\vec{e_3})\big)$ (base canonique de $\mathbb{R}^3$).
$g(1,0,0)=(1,0)$, $g(0,1,0)=(1,1)$, $g(0,0,1)=(0,-1)$.
$\operatorname{Im} g = \operatorname{Vect}\big((1,0),\;(1,1),\;(0,-1)\big)$.
$(1,0)$ et $(0,-1)$ sont libres (non colinéaires) et génèrent $\mathbb{R}^2$ (dimension 2).
Donc $\operatorname{Im} g = \mathbb{R}^2$.
Une base de $\operatorname{Im} g$ : $\boxed{\mathcal{B}_{\operatorname{Im}}=((1,0),\;(0,1))}$ (base canonique) ou $((1,0),(1,1))$. $\dim(\operatorname{Im} g)=2$.

\textbf{4.}
\begin{itemize}
    \item \textbf{Injective ?} Non, car $\ker g \neq \{\vec{0}\}$ ($\dim(\ker g)=1$).
    \item \textbf{Surjective ?} Oui, car $\operatorname{Im} g = \mathbb{R}^2$ (codomaine).
    \item \textbf{Théorème du rang :} $\dim(\mathbb{R}^3) = 3$. $\dim(\ker g) + \dim(\operatorname{Im} g) = 1 + 2 = 3$. \textbf{Vérifié}.
\end{itemize}
\end{corrige}

\begin{corrige}[Exercice 9 — La droite des coûts (type Probatoire)]
\textbf{Partie A : Structure de $E$}
$E=\{(x,y)\in\mathbb{R}^2 \mid y=3x\}$.

\textbf{1.} $E$ est une droite passant par l'origine (équation cartésienne $y-3x=0$).
\begin{itemize}
    \item $\vec{0}=(0,0)$ vérifie $0=3\times0$, donc $\vec{0}\in E$.
    \item Stable par addition : $(x_1,3x_1)+(x_2,3x_2)=(x_1+x_2,\;3(x_1+x_2))\in E$.
    \item Stable par multiplication par un scalaire : $\lambda(x,3x)=(\lambda x,\;3\lambda x)\in E$.
\end{itemize}
$E$ est un sous-espace vectoriel de $\mathbb{R}^2$, donc un \cle{espace vectoriel} sur $\mathbb{R}$ (muni des lois induites).

\textbf{2.} Tout vecteur de $E$ s'écrit $(x,3x)=x(1,3)$.
La famille $\mathcal{B}_E = \big((1,3)\big)$ est génératrice et libre (vecteur non nul).
C'est une \cle{base} de $E$. $\dim(E)=1$.

\vspace{0.5em}
\textbf{Partie B : Application linéaire $h$}
$h(x,y)=(x-2y,\;-2x+4y)$.

\textbf{1.} $h$ est linéaire (composantes formes linéaires). $h:\mathbb{R}^2\to\mathbb{R}^2$, donc $h$ est un \cle{endomorphisme} de $\mathbb{R}^2$.

\textbf{2.} $\ker h = \{(x,y)\mid x-2y=0,\; -2x+4y=0\}$.
La 2ème équation est $-2\times$(1ère). Système équivalent à $x-2y=0 \iff x=2y$.
\cle{Équation caractéristique} : $x-2y=0$.
Vecteurs : $(2y,y)=y(2,1)$.
Base de $\ker h$ : $\boxed{\mathcal{B}_{\ker}=((2,1))}$. $\dim(\ker h)=1$.

\textbf{3.} $\operatorname{Im} h = \operatorname{Vect}\big(h(1,0), h(0,1)\big) = \operatorname{Vect}\big((1,-2),\;(-2,4)\big)$.
$(-2,4) = -2(1,-2)$. Les deux vecteurs sont colinéaires.
$\operatorname{Im} h = \operatorname{Vect}\big((1,-2)\big)$.
\cle{Équation caractéristique} : $2x+y=0$ (car $(x,y)=\lambda(1,-2) \implies y=-2x$).
Base de $\operatorname{Im} h$ : $\boxed{\mathcal{B}_{\operatorname{Im}}=((1,-2))}$. $\dim(\operatorname{Im} h)=1$.

\textbf{4.} $h$ n'est pas injective ($\ker h \neq \{\vec{0}\}$), n'est pas surjective ($\operatorname{Im} h \neq \mathbb{R}^2$, dimension 1).
Donc $h$ n'est \textbf{pas bijective}.

\textbf{5.} Résolution de $h(x,y)=(3,-6)$.
$\begin{cases} x-2y=3 \\ -2x+4y=-6 \end{cases} \iff x-2y=3$ (2ème = $-2\times$1ère).
Ensemble des solutions : $x = 3+2y,\; y\in\mathbb{R}$.
$\boxed{S = \{(3+2t,\;t) \mid t\in\mathbb{R}\} = (3,0) + \mathbb{R}(2,1)}$.
\cle{Interprétation géométrique :} C'est la droite passant par $(3,0)$ de vecteur directeur $(2,1)$ (qui est un vecteur directeur de $\ker h$). C'est une \cle{droite affine} parallèle à $\ker h$ (fibre de l'antécédent).

\vspace{0.5em}
\noindent\textbf{Barème indicatif (Probatoire) :}
\begin{itemize}
    \item Partie A : 4 pts (Stabilité 2 pts, Base/Dim 2 pts).
    \item Partie B : 16 pts (Linéarité 2 pts, Noyau 4 pts, Image 4 pts, Bijectivité 2 pts, Résolution/Interprétation 4 pts).
\end{itemize}
\noindent\textbf{Points de vigilance :}
\begin{itemize}
    \item Citer explicitement les 3 conditions de sous-espace vectoriel (ou critère de stabilité).
    \item Pour $\ker h$ et $\operatorname{Im} h$, donner une \textbf{équation caractéristique} (système d'équations cartésiennes) ET une \textbf{base}.
    \item Justifier la non-bijectivité par $\dim(\ker h)\neq 0$ OU $\dim(\operatorname{Im} h)\neq 2$.
    \item Pour l'interprétation : « droite affine parallèle au noyau » ou « fibre de l'application ».
\end{itemize}
\end{corrige}

\begin{corrige}[Exercice 10 — Symétrie vectorielle et automorphisme (type Probatoire)]
$(\vec{i},\vec{j})$ base du plan $\mathcal{P}$. $s(x\vec{i}+y\vec{j}) = x\vec{i}-y\vec{j}$.

\textbf{Partie A : Étude de $s$}

\textbf{1.} Dans la base $(\vec{i},\vec{j})$, les coordonnées de $\vec{u}$ sont $(x,y)$.
$s(\vec{u})$ a pour coordonnées $(x,\;-y)$.
Expression analytique : $\boxed{s(x,y) = (x,\;-y)}$.
Matrice : $\begin{pmatrix} 1 & 0 \\ 0 & -1 \end{pmatrix}$.

\textbf{2.} Linéarité : $s(\lambda(x_1,y_1)+\mu(x_2,y_2)) = s(\lambda x_1+\mu x_2,\; \lambda y_1+\mu y_2) = (\lambda x_1+\mu x_2,\; -\lambda y_1-\mu y_2) = \lambda(x_1,-y_1)+\mu(x_2,-y_2) = \lambda s(x_1,y_1)+\mu s(x_2,y_2)$.
$s$ est un \cle{endomorphisme} de $\mathcal{P}$.

\textbf{3.}
\begin{itemize}
    \item $\ker s = \{(x,y)\mid (x,-y)=(0,0)\} = \{\vec{0}\}$.
    \item $\operatorname{Im} s = \{s(x,y)\mid (x,y)\in\mathcal{P}\} = \{(x,-y)\mid (x,y)\in\mathcal{P}\} = \mathcal{P}$ (car l'application $(x,y)\mapsto(x,-y)$ est surjective sur $\mathbb{R}^2$).
\end{itemize}

\textbf{4.} $\ker s = \{\vec{0}\} \implies s$ injective. Endomorphisme injectif d'un espace de dimension finie $\implies$ bijective $\implies$ \cle{automorphisme} de $\mathcal{P}$.

\textbf{5.} $s\circ s (x,y) = s(s(x,y)) = s(x,-y) = (x,y) = \operatorname{Id}_{\mathcal{P}}(x,y)$.
$s\circ s = \operatorname{Id}$. $s$ est une \cle{involution} (symétrie : $s^{-1}=s$).

\vspace{0.5em}
\textbf{Partie B : Une famille de vecteurs}
$\vec{u}=\vec{i}+\vec{j}$, $\vec{v}=s(\vec{u})=\vec{i}-\vec{j}$.

\textbf{1.} Montrer que $(\vec{u},\vec{v})$ est une base.
$\alpha\vec{u}+\beta\vec{v}=\vec{0} \iff \alpha(\vec{i}+\vec{j})+\beta(\vec{i}-\vec{j})=\vec{0} \iff (\alpha+\beta)\vec{i}+(\alpha-\beta)\vec{j}=\vec{0}$.
$(\vec{i},\vec{j})$ base $\implies \begin{cases} \alpha+\beta=0 \\ \alpha-\beta=0 \end{cases} \implies \alpha=\beta=0$.
Famille libre de 2 vecteurs dans un espace de dimension 2 $\implies$ \cle{base} de $\mathcal{P}$.

\textbf{2.} Coordonnées dans la base $(\vec{u},\vec{v})$.
\begin{itemize}
    \item $s(\vec{u}) = \vec{v} = 0\vec{u} + 1\vec{v}$. Coordonnées : $\boxed{\begin{pmatrix} 0 \\ 1 \end{pmatrix}_{(\vec{u},\vec{v})}}$.
    \item $s(\vec{v}) = s(\vec{i}-\vec{j}) = \vec{i}+\vec{j} = \vec{u} = 1\vec{u} + 0\vec{v}$. Coordonnées : $\boxed{\begin{pmatrix} 1 \\ 0 \end{pmatrix}_{(\vec{u},\vec{v})}}$.
\end{itemize}
Matrice de $s$ dans la base $(\vec{u},\vec{v})$ : $\begin{pmatrix} 0 & 1 \\ 1 & 0 \end{pmatrix}$ (matrice de permutation).

\vspace{0.5em}
\noindent\textbf{Barème indicatif (Probatoire) :}
\begin{itemize}
    \item Partie A : 12 pts (Expression 2, Linéarité 2, Noyau/Image 3, Automorphe 2, $s\circ s$ 3).
    \item Partie B : 8 pts (Base 4, Coordonnées 4).
\end{itemize}
\noindent\textbf{Points de vigilance :}
\begin{itemize}
    \item Distinguer « expression analytique » (coordonnées dans la base canonique) et « matrice ».
    \item Pour $\ker s$ et $\operatorname{Im} s$, écrire explicitement les ensembles.
    \item Justifier l'automorphisme par $\ker=\{\vec{0}\}$ + dimension finie (ou $\operatorname{Im}=\mathcal{P}$).
    \item Pour les coordonnées dans la nouvelle base, exprimer le vecteur image en combinaison linéaire de $\vec{u},\vec{v}$.
\end{itemize}
\end{corrige}

\begin{corrige}[Exercice 11 — Réseau de distribution d'eau (type Probatoire)]
$f(x,y,z)=(x+y+z,\;x-y,\;2x+z)$.

\textbf{Partie A : Linéarité}

\textbf{1.} $f$ est linéaire car chaque composante est une forme linéaire en $(x,y,z)$.
$f:\mathbb{R}^3\to\mathbb{R}^3$, donc $f$ est un \cle{endomorphisme} de $\mathbb{R}^3$.

\textbf{2.} Base canonique $(\vec{e_1},\vec{e_2},\vec{e_3}) = ((1,0,0),(0,1,0),(0,0,1))$.
\begin{align*}
f(\vec{e_1}) &= f(1,0,0) = (1+0+0,\;1-0,\;2+0) = \boxed{(1,1,2)}, \\
f(\vec{e_2}) &= f(0,1,0) = (0+1+0,\;0-1,\;0+0) = \boxed{(1,-1,0)}, \\
f(\vec{e_3}) &= f(0,0,1) = (0+0+1,\;0-0,\;0+1) = \boxed{(1,0,1)}.
\end{align*}

\vspace{0.5em}
\textbf{Partie B : Noyau et Image}

\textbf{1.} $\ker f = \{(x,y,z)\mid f(x,y,z)=(0,0,0)\}$.
Système (équation caractéristique) :
\[ \begin{cases} x + y + z = 0 & (1) \\ x - y = 0 & (2) \\ 2x + z = 0 & (3) \end{cases} \]
De (2) : $x=y$.
De (3) : $z = -2x$.
Substituons dans (1) : $x + x - 2x = 0 \iff 0=0$ (toujours vrai).
Les solutions sont : $(x,\;x,\;-2x) = x(1,1,-2),\quad x\in\mathbb{R}$.
$\ker f = \operatorname{Vect}\big((1,1,-2)\big) \neq \{\vec{0}\}$.
\textbf{Attention :} L'énoncé original (sujet d'examen) pouvait demander de « montrer que $\ker f=\{\vec{0}\}$ ». C'est \textbf{erroné} : le noyau est une droite (dimension 1). Le déterminant de la matrice de $f$ est nul ($\det=0$), $f$ n'est pas injective.

\textbf{2.} Comme $\ker f \neq \{\vec{0}\}$, $f$ n'est pas injective, donc n'est \textbf{pas un automorphisme} de $\mathbb{R}^3$.
(Correction de l'énoncé : on ne peut pas en déduire que $f$ est un automorphisme).

\textbf{3.} \cle{Théorème du rang} : $\dim(\mathbb{R}^3) = \dim(\ker f) + \dim(\operatorname{Im} f)$.
$3 = 1 + \dim(\operatorname{Im} f) \implies \dim(\operatorname{Im} f) = 2$.
$\operatorname{Im} f = \operatorname{Vect}\big(f(\vec{e_1}), f(\vec{e_2}), f(\vec{e_3})\big) = \operatorname{Vect}\big((1,1,2),\;(1,-1,0),\;(1,0,1)\big)$.
Les deux premiers vecteurs sont linéairement indépendants (non proportionnels).
Une base de $\operatorname{Im} f$ : $\boxed{\mathcal{B}_{\operatorname{Im}} = \big((1,1,2),\;(1,-1,0)\big)}$.
(On peut aussi prendre $(1,1,2)$ et $(1,0,1)$, ou $(1,-1,0)$ et $(1,0,1)$).

\vspace{0.5em}
\textbf{Partie C : Exploitation}

\textbf{1.} Antécédent de $(6,2,7)$ : résoudre $f(x,y,z)=(6,2,7)$.
\[ \begin{cases} x + y + z = 6 & (1) \\ x - y = 2 & (2) \\ 2x + z = 7 & (3) \end{cases} \]
De (2) : $y = x-2$.
De (3) : $z = 7-2x$.
Substituons dans (1) : $x + (x-2) + (7-2x) = 6 \iff 5 = 6$.
\cle{Contradiction.} Le système n'a \textbf{aucune solution}.
Le triplet $(6,2,7)$ n'appartient pas à $\operatorname{Im} f$ (vérification : $6+2-7=1\neq0$, alors que $\operatorname{Im} f = \{(u,v,w)\mid u+v-w=0\}$).
Il n'existe \textbf{pas d'antécédent}. La distribution demandée est impossible avec ce modèle linéaire.

\textbf{2.} Famille $\mathcal{B}=\big((1,1,2),\;(1,-1,0),\;(1,0,1)\big)$.
Ce sont exactement les vecteurs $f(\vec{e_1}), f(\vec{e_2}), f(\vec{e_3})$ calculés en A.2.
Calculons le déterminant de la matrice $M = \begin{pmatrix} 1 & 1 & 1 \\ 1 & -1 & 0 \\ 2 & 0 & 1 \end{pmatrix}$.
$\det M = 1(-1\cdot1 - 0) - 1(1\cdot1 - 0) + 1(0 - (-1)\cdot2) = -1 - 1 + 2 = 0$.
Le déterminant est nul, les trois vecteurs sont \cle{liés}.
$\mathcal{B}$ n'est \textbf{pas une base} de $\mathbb{R}^3$ (elle ne contient que 2 vecteurs indépendants).
\cle{Interprétation :} $\mathcal{B}$ est la famille des images des vecteurs de la base canonique par $f$. Elle \cle{génère l'image} de $f$ ($\operatorname{Im} f = \operatorname{Vect}(\mathcal{B})$) mais n'en est pas une base (famille liée). Une base de $\operatorname{Im} f$ s'en dégage en enlevant un vecteur (ex: les deux premiers).

\vspace{0.5em}
\noindent\textbf{Barème indicatif (Probatoire) :}
\begin{itemize}
    \item Partie A : 4 pts (Linéarité 1, Images base canonique 3).
    \item Partie B : 10 pts (Noyau 4, Automorphe 2 — \textbf{attention piège}, Rang/Image 4).
    \item Partie C : 6 pts (Antécédent 3, Famille B 3).
\end{itemize}
\noindent\textbf{Points de vigilance :}
\begin{itemize}
    \item \textbf{Calculer soigneusement le noyau.} Ne pas affirmer $\ker f=\{\vec{0}\}$ si le déterminant est nul. C'est un piège classique : vérifier le rang de la matrice.
    \item Si $\ker f \neq \{\vec{0}\}$, $f$ n'est \textbf{pas} un automorphisme. Ne pas écrire « on en déduit que $f$ est un automorphisme » si c'est faux.
    \item Pour l'antécédent, vérifier la compatibilité (le vecteur appartient-il à l'image ?). Si non, le dire clairement.
    \item La famille $\mathcal{B}$ est la famille des vecteurs colonnes de la matrice de $f$. Son rang = rang de $f$ = dimension de l'image.
\end{itemize}
\end{corrige}

\newpage

\coursec{Fiche bilan}

\begin{popfigure}
\begin{center}
\begin{center}\tcbox[colback=popink,colframe=popink,coltext=white,arc=9pt,boxsep=4pt,left=6pt,right=6pt]{\bfseries\small Espaces vectoriels sur $\mathbb{R}$}\end{center}
\vspace{-2pt}
\begin{tcbraster}[raster columns=3,raster equal height=rows,raster column skip=6pt,raster row skip=6pt,enhanced,blankest]
\begin{tcolorbox}[enhanced,colback=white,colframe=popblue,boxrule=0.9pt,arc=3pt,title={Lois de composition},fonttitle=\bfseries\scriptsize,coltitle=white,colbacktitle=popblue,left=4pt,right=4pt,top=2pt,bottom=2pt,boxsep=1pt,toptitle=1.5pt,bottomtitle=1.5pt,halign=left]
\begin{itemize}[nosep,leftmargin=*,itemsep=1pt,font=\scriptsize]
\item Loi externe : $\mathbb{R}\times E\to E$
\item 8 axiomes d'un e.v.
\end{itemize}
\end{tcolorbox}
\begin{tcolorbox}[enhanced,colback=white,colframe=poporange,boxrule=0.9pt,arc=3pt,title={Sous-espaces vectoriels},fonttitle=\bfseries\scriptsize,coltitle=white,colbacktitle=poporange,left=4pt,right=4pt,top=2pt,bottom=2pt,boxsep=1pt,toptitle=1.5pt,bottomtitle=1.5pt,halign=left]
\begin{itemize}[nosep,leftmargin=*,itemsep=1pt,font=\scriptsize]
\item $F\neq\emptyset$ ; stable par $+$ et $\cdot$
\item Test : $\lambda u+v\in F$
\end{itemize}
\end{tcolorbox}
\begin{tcolorbox}[enhanced,colback=white,colframe=popgreen,boxrule=0.9pt,arc=3pt,title={Familles de vecteurs},fonttitle=\bfseries\scriptsize,coltitle=white,colbacktitle=popgreen,left=4pt,right=4pt,top=2pt,bottom=2pt,boxsep=1pt,toptitle=1.5pt,bottomtitle=1.5pt,halign=left]
\begin{itemize}[nosep,leftmargin=*,itemsep=1pt,font=\scriptsize]
\item G\'en\'eratrice : tout $u$ est combinaison
\item Libre : $\sum\lambda_i u_i=0 \Rightarrow \lambda_i=0$
\item Li\'ee = non libre
\end{itemize}
\end{tcolorbox}
\begin{tcolorbox}[enhanced,colback=white,colframe=popgold,boxrule=0.9pt,arc=3pt,title={Bases et dimension},fonttitle=\bfseries\scriptsize,coltitle=white,colbacktitle=popgold,left=4pt,right=4pt,top=2pt,bottom=2pt,boxsep=1pt,toptitle=1.5pt,bottomtitle=1.5pt,halign=left]
\begin{itemize}[nosep,leftmargin=*,itemsep=1pt,font=\scriptsize]
\item Base = libre + g\'en\'eratrice
\item Toutes les bases : m\^eme cardinal
\item $\dim\mathbb{R}^n=n$
\end{itemize}
\end{tcolorbox}
\begin{tcolorbox}[enhanced,colback=white,colframe=poppurple,boxrule=0.9pt,arc=3pt,title={Applications lin\'eaires},fonttitle=\bfseries\scriptsize,coltitle=white,colbacktitle=poppurple,left=4pt,right=4pt,top=2pt,bottom=2pt,boxsep=1pt,toptitle=1.5pt,bottomtitle=1.5pt,halign=left]
\begin{itemize}[nosep,leftmargin=*,itemsep=1pt,font=\scriptsize]
\item $f(\lambda u+v)=\lambda f(u)+f(v)$
\item Endo, iso, automorphisme
\item Image d'une base
\end{itemize}
\end{tcolorbox}
\begin{tcolorbox}[enhanced,colback=white,colframe=poppink,boxrule=0.9pt,arc=3pt,title={Noyau et image},fonttitle=\bfseries\scriptsize,coltitle=white,colbacktitle=poppink,left=4pt,right=4pt,top=2pt,bottom=2pt,boxsep=1pt,toptitle=1.5pt,bottomtitle=1.5pt,halign=left]
\begin{itemize}[nosep,leftmargin=*,itemsep=1pt,font=\scriptsize]
\item $\ker f=\{u\mid f(u)=0\}$
\item $\operatorname{Im}f=\{f(u)\}$
\item $f$ injective $\iff\ker f=\{0\}$
\end{itemize}
\end{tcolorbox}
\end{tcbraster}

\legende{Carte mentale de la le\c{c}on : les six blocs de comp\'etences des espaces vectoriels.}
\end{center}
\end{popfigure}

\begin{aretenir}[L'essentiel de la le\c{c}on]
\textbf{D\'efinitions cl\'es.}
\begin{itemize}
  \item \cle{Loi de composition externe} sur $E$ : application $\mathbb{R}\times E\to E$, $(\lambda,u)\mapsto\lambda\cdot u$.
  \item \cle{Espace vectoriel sur $\mathbb{R}$} : ensemble $E$ muni d'une addition interne et d'une multiplication par les r\'eels v\'erifiant les 8 axiomes (associativit\'e, commutativit\'e, neutre $0_E$, oppos\'e, distributivit\'es, $1\cdot u=u$, $(\lambda\mu)u=\lambda(\mu u)$).
  \item \cle{Sous-espace vectoriel} : partie $F\subset E$ non vide telle que pour tous $u,v\in F$ et tout $\lambda\in\mathbb{R}$ : $\lambda u+v\in F$.
  \item Famille \cle{g\'en\'eratrice} : tout vecteur de $E$ s'\'ecrit comme combinaison lin\'eaire des vecteurs de la famille.
  \item Famille \cle{libre} : $\lambda_1u_1+\dots+\lambda_pu_p=0_E \Rightarrow \lambda_1=\dots=\lambda_p=0$. Sinon elle est \cle{li\'ee}.
  \item \cle{Base} : famille \`a la fois libre et g\'en\'eratrice. Tout vecteur s'\'ecrit alors de fa\c{c}on \emph{unique} comme combinaison lin\'eaire des vecteurs de la base.
  \item \cle{Application lin\'eaire} $f:E\to F$ : pour tous $u,v\in E$, $\lambda\in\mathbb{R}$ : $f(u+v)=f(u)+f(v)$ et $f(\lambda u)=\lambda f(u)$.
  \item \cle{Noyau} : $\ker f=\{u\in E\mid f(u)=0_F\}$ ; \cle{Image} : $\operatorname{Im}f=\{f(u)\mid u\in E\}$.
\end{itemize}

\textbf{Th\'eor\`emes et formules \`a conna\^itre par c\oe ur.}
\begin{center}
\begin{tabularx}{\linewidth}{|l|X|}
\hline
\rowcolor{popblueL} \textbf{R\'esultat} & \textbf{\'Enonc\'e} \\
\hline
Dimension & Si $E$ admet une famille g\'en\'eratrice finie, toutes ses bases ont le m\^eme nombre de vecteurs : la \cle{dimension} de $E$. \\
\hline
Dimensions usuelles & $\dim\mathbb{R}=1$ ; $\dim\mathbb{R}^2=2$ (base canonique $(e_1,e_2)$) ; $\dim\mathbb{R}^3=3$ ; plan vectoriel : $2$ ; espace vectoriel : $3$. \\
\hline
Crit\`ere de base & En dimension $n$, toute famille libre de $n$ vecteurs (ou g\'en\'eratrice de $n$ vecteurs) est une base. \\
\hline
Noyau, image & $\ker f$ est un sous-espace de $E$ ; $\operatorname{Im}f$ est un sous-espace de $F$, engendr\'e par les images des vecteurs d'une base de $E$. \\
\hline
Injectivit\'e & $f$ lin\'eaire injective $\iff \ker f=\{0_E\}$. \\
\hline
Bijectivit\'e & $f:E\to F$ lin\'eaire avec $\dim E=\dim F$ finie : injective $\iff$ surjective $\iff$ bijective (isomorphisme). \\
\hline
\end{tabularx}
\end{center}

\textbf{M\'ethodes express.}
\begin{enumerate}
  \item \emph{Prouver qu'une partie est un sous-espace} : v\'erifier $0_E\in F$, puis $\lambda u+v\in F$ pour $u,v\in F$ g\'en\'eriques.
  \item \emph{Libert\'e d'une famille} : poser $\sum\lambda_i u_i=0$, traduire en syst\`eme, montrer que la seule solution est nulle.
  \item \emph{Base du noyau} : r\'esoudre $f(u)=0$ (\'equation(s) caract\'eristique(s)), param\'etrer les solutions, lire une famille libre g\'en\'eratrice.
  \item \emph{Base de l'image} : calculer les images des vecteurs de la base canonique, extraire une sous-famille libre maximale.
  \item \emph{Lin\'earit\'e} : tester $f(u+v)=f(u)+f(v)$ et $f(\lambda u)=\lambda f(u)$, ou directement $f(\lambda u+v)=\lambda f(u)+f(v)$.
\end{enumerate}
\end{aretenir}

\begin{attention}[Pr\'ecisions et pi\`eges \`a conna\^itre]
\begin{itemize}
  \item L'espace nul $\{0_E\}$ est un espace vectoriel ; par convention $\dim\{0_E\}=0$ (il n'admet pas de base, ou une base vide).
  \item Pour le test de sous-espace, la condition $F\neq\emptyset$ se v\'erifie en pratique en montrant que $0_E\in F$ : ne l'oublie jamais, c'est un point fr\'equemment sanctionn\'e au Probatoire.
  \item Une famille contenant le vecteur nul est toujours li\'ee ; deux vecteurs sont li\'es si et seulement si ils sont colin\'eaires.
  \item Le crit\`ere << libre de $n$ vecteurs en dimension $n$ $\Rightarrow$ base >> exige de conna\^itre au pr\'ealable la dimension de l'espace : cite-le explicitement dans ta r\'edaction.
  \item Le crit\`ere << injective $\iff$ surjective $\iff$ bijective >> n'est valable que si $\dim E=\dim F$ \emph{finie} ; il ne s'applique pas \`a une application entre espaces de dimensions diff\'erentes.
  \item Pour prouver qu'une application n'est \emph{pas} lin\'eaire, un seul contre-exemple suffit (souvent $f(0_E)\neq 0_F$).
\end{itemize}
\end{attention}

\begin{exemplebox}[Deux illustrations rapides]
\textbf{1. Base du noyau.} Soit $f:\mathbb{R}^3\to\mathbb{R}$, $f(x,y,z)=x-2y+z$. L'\'equation caract\'eristique du noyau est $x-2y+z=0$, soit $x=2y-z$. Tout vecteur de $\ker f$ s'\'ecrit $(2y-z,\,y,\,z)=y(2,1,0)+z(-1,0,1)$. La famille $\big((2,1,0),(-1,0,1)\big)$ est libre et engendre $\ker f$ : c'est une base de $\ker f$, et $\dim\ker f=2$.

\textbf{2. Base de l'image.} Soit $g:\mathbb{R}^2\to\mathbb{R}^2$, $g(x,y)=(x+y,\,2x+2y)$. On a $g(e_1)=(1,2)$ et $g(e_2)=(1,2)$ : $\operatorname{Im}g$ est engendr\'e par $(1,2)$ seul, qui est non nul donc libre. Une base de $\operatorname{Im}g$ est $\big((1,2)\big)$ et $\dim\operatorname{Im}g=1$. Comme $\ker g=\{(x,y)\mid x+y=0\}\neq\{0\}$, $g$ n'est pas injective, donc pas bijective.
\end{exemplebox}

\begin{aretenir}[Checklist avant le Probatoire]
\begin{itemize}
  \item[$\square$] Je sais citer les 8 axiomes d'un espace vectoriel sur $\mathbb{R}$ et les v\'erifier sur un exemple.
  \item[$\square$] Je sais montrer qu'une loi $\mathbb{R}\times E\to E$ est une loi de composition externe.
  \item[$\square$] Je sais prouver qu'une partie est un sous-espace vectoriel avec le test $\lambda u+v\in F$ (sans oublier $0_E\in F$).
  \item[$\square$] Je sais d\'ecider si une famille est libre ou li\'ee en r\'esolvant un syst\`eme homog\`ene.
  \item[$\square$] Je sais montrer qu'une famille est g\'en\'eratrice de $\mathbb{R}^2$ ou $\mathbb{R}^3$.
  \item[$\square$] Je sais qu'une base est une famille libre \emph{et} g\'en\'eratrice, et je connais les bases canoniques de $\mathbb{R}$, $\mathbb{R}^2$, $\mathbb{R}^3$.
  \item[$\square$] Je sais que toutes les bases d'un m\^eme espace ont le m\^eme cardinal : la dimension.
  \item[$\square$] Je sais v\'erifier qu'une application est lin\'eaire et reconna\^itre endomorphisme, isomorphisme, automorphisme.
  \item[$\square$] Je sais d\'eterminer une \'equation caract\'eristique et une base du noyau d'une application lin\'eaire.
  \item[$\square$] Je sais d\'eterminer une base de l'image \`a partir des images des vecteurs d'une base.
  \item[$\square$] Je sais que $f$ injective $\iff\ker f=\{0\}$, et j'utilise ce crit\`ere pour prouver la bijectivit\'e.
  \item[$\square$] Je r\'edige proprement : quantificateurs, vecteurs g\'en\'eriques, conclusion explicite.
\end{itemize}
\end{aretenir}

\findecours

\end{document}
